% ============================================================================
%  ASSUMPTIONS -- ONLINE SUPPLEMENT.  project `panel-npiv-discrete`
%
%  WHAT THIS FILE HOLDS.  Everything the assumption block prints that is not a
%  statement: the per-condition interpretation, the retirement records for A3,
%  A6, B1(vi) and A2'(ii), the status list for the five clauses no theorem rests
%  on, and the smoothness-budget table.  Copied verbatim from assumptions.tex;
%  nothing here was rewritten.
%
%  IT IS READ AFTER assumptions_main.tex AND NOT INSTEAD OF IT.  The paragraphs
%  below cite the statements they discuss, so the supplement inputs the
%  statements first.  Together the two files are the undivided assumptions.tex.
%
%%  DO NOT \input THIS FILE ALONGSIDE results.tex OR assumptions.tex.
%%  results.tex and assumptions.tex SURVIVE UNDIVIDED and remain the source of
%%  record: scripts/build-theory/build.sh compiles them, scripts/render-sync/
%%  check_rendering_sync.py reads them, and ~100 files under scripts/symbolic/
%%  specs/ cite them BY LINE NUMBER.  Neither was edited.  Every \label in this
%%  file is therefore declared TWICE on disk -- here and in its parent -- and the
%%  two must never reach the same document.
%
%%  THE PARAGRAPH ABOVE IS STALE AS OF 2026-08-05 AND IS NOT AMENDED IN
%%  DISPATCH 25, WHICH WOULD BE OUT OF SCOPE.  quality.md's binary gate 1 was
%%  repointed that day: results.tex and assumptions.tex are compiled by NO
%%  document (build.sh:46 records them RETIRED; driver.tex inputs the eight split
%%  files), and the two copies have DIVERGED at cor:a1cprim.  DISPATCH 25 EDITED
%%  NEITHER, so the divergence widens by everything this dispatch did.  That is
%%  the intended state -- the split files are the corpus -- and it is recorded
%%  here so the next reader of the "declared TWICE" sentence does not try to
%%  restore parity by hand.
% ============================================================================

% ============================================================================
%  DISPATCH-25 CHANGE LOG (2026-08-08).  PHASE 1 OF THE MULTINOMIAL-INSTRUMENT
%  PLAN.  The full log, including everything done to assumptions_main.tex, is at
%  the foot of that file; this is what changed HERE.
%
%  NEW PARAGRAPH: "A1 (conditional exogeneity)", placed first.  It proves the
%  three facts the sans-L-category representation rests on -- the contrasts span
%  the conditionally centred functions exactly, the span is base-cell free, and
%  restriction (A) at the contrasts is EQUIVALENT to A1 rather than implied by
%  it -- and it records what extra cells buy (LOCALISATION; round 5 replaced the
%  round-1 "efficiency and overidentification" motivation, A1 being testable at
%  sans-L = 2)
%  and what they do not (identification, which already held at sans-L = 2 and
%  whose hypothesis a finer instrument WEAKENS).  It carries a standing
%  prohibition against explaining the gain through the ill-posed inverse problem.
%
%  THREE NUMBERED DISPLAYS ADDED IN THIS FILE:
%     eq:omegainv     Omega^{-1} in closed form, and Omega^{-1} Wtilde as the
%                     vector of inverse-probability contrasts against the base
%                     (A4 paragraph)
%     eq:omegaprefac  E[Wtilde 1{Delta U <= u} | Delta Z] = Omega d
%                     (A4 paragraph)
%     eq:omeganorm    the normalised contrast Wtilde^Omega (rem:htread).  ADDED
%                     IN ROUND 4; it stood in assumptions_main.tex in rounds 2-3,
%                     where the estimator carried it, and MOVES here with the
%                     reading when the estimator stops carrying it.
%  Four further displays are added in assumptions_main.tex (eq:contrast,
%  eq:omegadef, eq:omegafloor, eq:trunset).  SEVEN IN TOTAL ACROSS THE TWO FILES.
%  *** THE TWO CHANGE LOGS MUST AGREE ON THIS NUMBER.  Rounds 2-3 had this log
%  saying SIX while assumptions_main.tex said EIGHT -- two logs contradicting each
%  other about their own dispatch, which is the defect species this block spends
%  its length policing.  Both now say seven. ***
%
%  ONE NUMBERED ENVIRONMENT IS ADDED IN THIS FILE, IN ROUND 4, AND IT IS THE ONLY
%  COUNTER CHANGE IN THE DISPATCH: rem:htread, at the foot beside the four
%  remarks already there.  This file is position 5 of 8, so every theorem-like
%  number in results_s3_supp.tex, results_supp.tex and proofs.tex ADVANCES BY
%  ONE.  Nothing breaks -- every reference is a \cref -- but printed numbers move.
%  The other four remarks keep their positions and their relative order.
%
%  PARAGRAPHS AMENDED: A2' (the (P1) bound and the Khan-Tamer computation),
%  A2c (ROUND 2: the (m.1) exceptional set is now the FINITE intersection over
%  the p(sans-L - 1) slices, and the (m.2) measurability argument is run at
%  ell^Pi_j(., t, l) rather than at the two-argument object -- round 1 amended
%  A2c's STATEMENT and left its discussion at the old object, so the printed
%  argument did not establish the printed display), A3
%  (a note; THE RETIRED DISPLAY IS DELIBERATELY LEFT AS PRINTED, being an
%  archaeological record), A4 (three derivations prepended; the measurability
%  clause read at the vector), "the sample operator T-hat" (the simplex
%  truncation and what it buys), rem:hatTread(2) and (3), and the Discussion's
%  A4 sentence.
%
%  ROUND-2 CORRECTIONS TO ROUND-1 TEXT IN THIS FILE (against theorist-critic
%  61/100).  Each replaces something round 1 printed; none is an extension.
%   A1 paragraph, third fact:  the equivalence was stated "under eq:omegafloor",
%       i.e. under overlap.  IT NEEDS NO OVERLAP.  The two-case argument via
%       Omega_A now stands in place of the invertibility appeal.
%   A1 paragraph, coarsening:  "V_CI = {varphi_dagger} at the coarsening gives
%       the same at the full instrument" delivered only the INCLUSION.
%       Non-emptiness is A1, and it is now stated as the second half.
%   A4 paragraph, degenerate region:  {det Omega = 0} is NOT the degenerate
%       region at sans-L >= 3 and the restriction is not vacuous there; the
%       region is {Omega = 0} = {max_l varrho_l = 1}.  The withdrawn claim, the
%       sans-L = 3 counterexample, and what lem:bierens does and does not need
%       are printed.
%   A4 paragraph, determinant:  det Omega = prod varrho_l was attributed to
%       eq:omegainv, which presupposes the invertibility that fails exactly where
%       the identity is wanted.  It is now proved by the matrix determinant lemma
%       on the open region plus polynomial continuation.
%   A4 paragraph, WHERE the normalisation goes:  round 1 said which normalisation
%       and never where.  ~~The theta moment (redundant, with the bijection of the
%       weight class proved)~~ -- STRUCK IN ROUND 3; the bijection is false and
%       psi normalises too -- and the varphi operator (adopted, with the
%       sans-L = 2 check and the costs) are now separated.
%   Overidentification:  the "pin from a subset" motivation presumed
%       identification from a SINGLETON contrast, which the next sentence itself
%       said is unavailable.  Replaced by the joint-consistency route.
%
%  ROUND-3 CORRECTIONS IN THIS FILE (against theorist-critic 45/100), TAGGED WITH
%  THE FINDING EACH CLOSES.
%   [AUTHOR DECISION]  A4 paragraph, "where the normalisation is applied": the
%       answer is now BOTH SIDES.  The round-2 claim that Omega^{-1} is redundant
%       on psi, and the bijection-of-the-weight-class argument for it, are STRUCK
%       and replaced by the refutation -- def:weight's clause (a) is a RANGE
%       condition, ell in R(cal-W), and a merely Hoelder Omega(z) takes a
%       real-analytic section out of that range.  The four bound-checks are kept
%       on record as correct-but-insufficient.
%   [C-2]  rem:hatTread(2): "at most two entries are nonzero" is FALSE -- at
%       W_i = 0 all sans-L - 1 entries equal -1/hat-varrho_0, which the same
%       clause's own lower bound says two lines earlier.  The Euclidean bound is
%       now sqrt(sans-L - 1) eps_varrho^{-1}, SHARP.  This was PRINTED text and
%       it understated at sans-L = 4, i.e. at Section 6's design.
%   [M-1]  the lf_panel.R misquotation is struck from the A4 paragraph; a global
%       scalar cannot move a root set and the file no longer suggests it can.
%   [M-2]  "efficiency and conditioning" struck from the A4 paragraph; the ground
%       is interpretability.
%   ~~[M-3]  the cost list in the A4 paragraph is five items and is explicitly
%       not asserted complete~~ -- STRUCK IN ROUND 4: the cost list existed only
%       under the normalisation and went with it.
%   ~~[C-3]  the lem:bierens hand-off is reconciled IN THIS FILE as well: the
%       overlap-free statement is about (a) <=> (b), and the clauses naming M or
%       T carry A4(i) as a well-formedness condition.~~  STRUCK IN ROUND 4 AND
%       LEFT UNSTRUCK BY IT, WHICH ROUND 5 REPAIRS.  With the raw contrast, M and
%       T are defined everywhere and bounded, so NO clause of lem:bierens carries
%       A4(i) and the hand-off is unqualified at every clause.  Reversed in the
%       printed text at the A1 paragraph and at A4's note (b); this log entry was
%       the last site still asserting the carve-out.
%   [m-2]  the four bound-checks now name which norm is spectral and which
%       Frobenius, and cite submultiplicativity across the two.
%   [m-5]  Z_0 is no longer redefined in an unlabelled display; eq:degenreg is
%       cited and the reading fixed.
%   [m-6]  "strictly smaller" is scoped to the simplex, with the note that the
%       z-sets may coincide for a particular DGP.
%   [ROUND 3, no finding]  the A3 retirement note records that lem:momzero(iv)'s
%       vanishing survives the normalisation, since Omega is Delta Z-measurable.
%
%  ROUND-4 CHANGES IN THIS FILE.  THE AUTHOR REVERSED THE NORMALISATION; THE
%  ESTIMATOR CARRIES THE RAW CONTRAST ON BOTH SIDES.
%   [AUTHOR DECISION]  The A4 paragraph's "where is it applied" section now
%       answers "nowhere in the estimator", with the author's ground -- A2, A2',
%       A1c and H1 are stated at A_cond, and mu = s/(2a), eq:frontier, the
%       budget table and the CFF comparison sit downstream of a.  The round-3
%       "repair is to normalise psi as well" paragraph is struck and replaced by
%       what the reversal leaves untouched.  The bijection refutation is KEPT,
%       because the same analytic-versus-Hoelder mechanism is the author's
%       ground.
%   [NEW]  rem:htread, the only numbered environment the dispatch adds.  It
%       carries eq:omeganorm, the Horvitz-Thompson reading, the base-cell-free
%       quadratic form, the two moments, and clause (6) saying why the estimator
%       does not use any of it.  IT SAYS IN ITS FIRST SENTENCE THAT IT IS
%       INTERPRETATION AND NOT SPECIFICATION.
%   [CRITICAL]  the two moments of the normalised contrast are stated correctly
%       at rem:htread(3)-(4): E[Wtilde' Omega^{-1} Wtilde | DZ] = sans-L - 1 is
%       BOUNDED; E[||Omega^{-1} Wtilde||^2 | DZ] = tr(Omega^{-1}) DIVERGES; and
%       E[|Wtilde^Omega_l| | DZ] = 2 exactly, FREE of varrho.  Round 3 computed
%       the first and cited it where the second belonged, and had the
%       varrho-freedom backwards.  Numbers verified at sans-L = 4.
%   [CRITICAL]  "not integrable without A4(i)" struck at both sites (the A1
%       paragraph's lem:bierens qualification and the A4 cost list).  Where a
%       cell vanishes Wtilde^Omega is UNDEFINED; where defined it is integrable
%       with a varrho-free bound.
%   [m-a]  the A1 spanning derivation now carries the degenerate-z count --
%       one less than the number of non-null cells -- so this file no longer
%       leaves a reader concluding that spanning needs A4(i), which the notation
%       block explicitly denies.
%   [M-C]  the display count here now matches assumptions_main.tex at seven.
%   rem:hatTread(2)  the normalised-contrast clause is REMOVED and clause (c) is
%       restated on the ground that survives: hat-varrho must be a probability
%       vector.  The paragraph is explicit that clauses (a) and (b) would hold on
%       the box, so the simplex is not load-bearing for them.
%   "the sample operator T-hat"  same restatement.
%   the A3 retirement note  the round-3 normalisation qualification is marked
%       true-but-idle rather than deleted.
%
%  ROUND-5 CHANGES IN THIS FILE, TAGGED.
%   [C-1]  the round-3 log entry [C-3] above -- "the clauses naming M or T carry
%       A4(i) as a well-formedness condition" -- was left UNSTRUCK by round 4
%       while the printed text reversed it.  Struck.
%   [M-1]  the A4 paragraph's "...what makes the NORMALISED OPERATOR OF THE NEXT
%       PARAGRAPH well defined" pointed in PRINTED TEXT at an object round 4 had
%       removed; the next paragraph is "a claim withdrawn here".  Rewritten to
%       point at rem:htread.
%   [M-2]  Omega-hat_i is DECLARED, at its first printed consumer in the T-hat
%       paragraph.  Round 4 removed the declaration from def:hatT with the
%       normalisation -- correctly -- and kept all four consumers, which is the
%       round-1 orphan defect mirrored.
%   [AUTHOR DECISION 2]  the A1 paragraph's "what extra cells buy" is rewritten
%       on LOCALISATION.  "Converts A1 from untestable" and the sans-L - 2 count
%       are struck: A1 is testable at sans-L = 2, the restriction being a
%       continuum of equations against one function, and Section 6's own probe
%       rejects it there.  The three concepts -- testability, instrument
%       strength, localisation -- are separated, and the Section 6
%       chronic-versus-transient split is named as what the design tests.  Joint
%       consistency stays as the APPARATUS.
%   [m-1]  the tr(Omega^{-1}) example's ratio is 1.79; 2(sans-L - 1) is the worst
%       case and is no longer attached to the example.
%   [m-6]  the spanning uniqueness step cites the paragraph's own standing
%       hypothesis, not eq:omegafloor, which is A4(i)'s global consequence and
%       whose citation there suggested spanning needs A4(i).
%
%  ONE NUMBER IS LEFT STALE ON PURPOSE AND IS FLAGGED AT ITS SITE: the
%  Discussion's 689/300 split and the 18%/9% figures are the two-category
%  reading of the application.  Phase 9 replaces them; this dispatch does not
%  have the four-category split and will not invent it.
% ============================================================================

\section*{Discussion of the assumptions}

\noindent This section carries the material the article does not print: for each
numbered condition of the assumption block, the interpretation, the provenance,
the record of what it does and does not restrict, and -- where a condition was
retired -- the ground on which it was retired. The statements themselves are
printed above, verbatim and unchanged. Headings below name the condition; the
text under each is reproduced verbatim from the consolidated block.

% ------------------------------------------------------------------------
%  MOVED HERE FROM assumptions_main.tex, 2026-08-28, ON THE AUTHOR'S SHORTENING
%  INSTRUCTION.  It stood as a \paragraph immediately under \section*{Assumptions}
%  and printed ~1,700 words IN THE ARTICLE.  It is bookkeeping about how a
%  citation is READ -- register taxonomy, the split test, the three
%  analyst-constant registers, the refutation of the superseded negation test,
%  and the two positions the taxonomy leaves out.  None of it is a condition and
%  none of it tells an applied reader what any condition rules out, which is the
%  test the article-side text is now held to.  WHAT STAYED IN THE ARTICLE is a
%  three-sentence statement of the rule itself, because A0, A4(i), A8, B1(iii)
%  and B1(iv) are written in the split form and prp:frechet, prp:compact and five
%  sites in results_supp.tex cite "the three registers" by name.
%  NO MATHEMATICS MOVED AND NO CLAUSE CHANGED STATUS.  The text is verbatim
%  except that "this block", whose referent was the surrounding assumption block,
%  is spelled "the assumption block" at the four sites where the relocation would
%  otherwise have made it read as the Discussion.
% ------------------------------------------------------------------------
\paragraph{The three registers: how a citation of the assumption block is read.}
A tag is named in more places than three; the three below are the three that have produced mis-citations
and been repaired against. \emph{No claim of exhaustiveness is made, and that is deliberate} --- the
two positions left out are named at the end of this paragraph, and one of them is where the
under-charging rule the assumption block enforces actually lives.
\emph{On a result's hypothesis line} --- inside a ``Suppose \dots\ hold'' --- a citation is a
\textbf{carrier}: naming a tag carries every condition under it, and naming a clause carries that clause
and no other. \emph{Inside a definition} a citation is a \textbf{pointer} \emph{when it locates}:
\cref{def:spaces}'s ``$\Pi=\Pi_u\otimes\Pi_t\otimes\Pi_{\mathsf w}$ as in Assumption~B1(iv)'' names the
object it uses and imposes nothing, so a result whose proof reads B1(iv.d) must carry that clause on its
own line and cannot borrow it from the definition. \emph{It is \textbf{not} a pointer when it is
constitutive}, and two definitions are: A2b \emph{is} ``admits a version satisfying
A2$'$(i.a)--(i.b)'', the citation being its whole content, and \cref{def:hatT} builds
\eqref{eq:trunset} out of the constant A4(i) declares. \emph{Ask what survives deleting the citation}: a
pointer's object is defined elsewhere and survives; a constitutive citation's object does not exist
without it. \emph{The third register is neither.} Several conditions have a \textbf{declarative
half}: the part that introduces a symbol, a type, a carrier set or a construction, as against the part
that restricts it.
\emph{A citation of a declarative half is a \textbf{declaration} and carries no hypothesis}, because a
declarative half is \textbf{standing}: it is in force at every statement in this corpus, named or not,
since without it the statement has no objects to be about. It is therefore written \emph{outside} the
``Suppose \dots\ hold'' list, and never with the word \emph{clause}.

\emph{Where the split falls, so that it is checkable and not a matter of taste.} Cut the fragment at the
point where it stops \textbf{introducing} and starts \textbf{restricting}. What introduces is the
declaration; the test that what remains is a \textbf{clause} is that \emph{there is a spelling of the
assumption block, with the declaration in force, under which every symbol used anywhere in the corpus
still has a referent and the remainder is false}. B1(iii) is the model and is already written this way:
$\mathcal T\subseteq\IR^{p}$ introduces the index set, $\mathrm{Leb}(\mathcal T)>0$ restricts it and can
fail with $\mathcal T$ still named, so the first is a declaration and the second a clause. B1(iv)'s
\textbf{stem} $\Pi=\Pi_u\otimes\Pi_t\otimes\Pi_{\mathsf w}$ is all introduction and leaves no remainder
--- it is what makes \cref{def:spaces}'s three-factor pairing, the masses $\varpi_{\mathsf w}(l)$ and the
third argument of $\ell^\Pi(u,t,l)$ writable at all --- so it is a declaration entire, while B1(iv.a) and
B1(iv.d) are clauses: negate either and \suppref{Proposition}{wellposed} is still a well-formed claim,
with a wrong constant. H1's $H^{s}(\mathcal X)$, the \emph{classical} scale read for non-integer $s$ as
the interpolation space and carrying no boundary condition, is a declaration on the same ground, while
H1's rank, its decay and \eqref{eq:h1match} are clauses.

\emph{Three fragments are split that were not split before, and in each the remainder is a clause
and is carried by the results that read it.}
\textbf{A0} introduces $\mathcal X\subseteq\IR$ as the carrier of the regressor, with
$\underline x:=\inf\mathcal X$ and $\bar x:=\sup\mathcal X$ --- \emph{spelled that way from 2026-08-26,
because as endpoints of an interval the two symbols lose their referents exactly when the restriction is
negated, which is the trap A4(i) fell into} --- and \emph{restricts} it to a compact
interval: a fat Cantor $\mathcal X$ leaves every symbol named and the restriction false, so the
restriction is a clause, and \suppref{Lemma}{intker} accordingly carries it. \textbf{A4(i)} \emph{fixes}
$\varepsilon_\varrho\in(0,1/\mathsf L]$ --- which is what \cref{def:hatT} builds \eqref{eq:trunset} from,
and all that definition reads of A4(i), \eqref{eq:trunset} being non-empty on
$\varepsilon_\varrho\le1/\mathsf L$ alone --- and \emph{imposes}
$\varrho_l(\Delta Z)\ge\varepsilon_\varrho$, which is the overlap condition and is a clause at each of
the results carrying it. \emph{The upper end of that range is \textbf{weak} from 2026-08-27 and was
strict before; why, and why the change is forced by the paragraph after next rather than by taste, is
recorded there.} \textbf{A8} names
$b_1,\dots,b_{L_x}$ and $\beta_1,\dots,\beta_{L_u}$ as the two bases, which is what \cref{def:lfstep}
builds $\mathcal S_n$ from, and requires of them a B-spline of degree exceeding $s-\tfrac12$ on a
quasi-uniform simple-knot mesh inside $n^{1/(2s)}\ll L_x\ll n^{1/2}$: restrictive, refutable at A8's own
$s=1.7$ witness, and correctly carried on \suppref{Lemma}{unifop}'s line.

\emph{One family needs the rule at a finer grain, because there the split falls \textbf{inside} a
numeric range: an interval declared for a constant the \textbf{analyst} chooses.} Three such constants
are printed in this corpus --- A4(i)'s $\varepsilon_\varrho$, \cref{def:lfstep}'s $\varepsilon_\pi$ and
A7's $\tau_{\mathrm{DP}}$ --- in three different registers, which is a defect unless one rule decides all
three, and \emph{is} a defect until the rule is printed. \emph{The rule, and it is the split-first test
applied to a range rather than to a sentence.} A declared range is a declaration exactly as far as its
\textbf{referent boundary}: the widest range on which every symbol used anywhere in this corpus still has
a referent. Whatever the printed range cuts off \emph{strictly inside} that boundary restricts nothing
but the analyst's own hand, is therefore a \textbf{clause} by the test above --- the witness being any
spelling in the gap --- and must either sit where $\tau_{\mathrm{DP}}$'s bound sits, on a hypothesis
line, or be struck. \emph{The three then come out in three registers because the rule puts them there,
not because three conventions are in use.}
\textup{(a)} $\varepsilon_\varrho$. \eqref{eq:trunset} is non-empty \emph{if and only if}
$\varepsilon_\varrho\le1/\mathsf L$, the witness at equality being the single point
$(1/\mathsf L,\dots,1/\mathsf L)$; so the referent boundary is the \textbf{weak} inequality and the strict
one the assumption block printed through 2026-08-27 lay inside it, at the one spelling
$\varepsilon_\varrho=1/\mathsf L$. It is \textbf{struck} rather than relocated, on the block's
retirement discipline: it excludes no law --- a $P$ meeting overlap at some constant meets it at half
that constant --- and no step in this corpus consumes it. \emph{That last is a reading and not a census,
and it is the half of this ruling a later pass should re-measure first.}
\textup{(b)} $\varepsilon_\pi$. \cref{def:lfstep}'s projection target is non-empty \emph{if and only if}
$\varepsilon_\pi\le1$, the witness at equality being the constant function $1$; so there the printed
range \textbf{is} the referent boundary, the remainder is empty, and there is nothing to charge.
\emph{That is why that definition prints a declaration and no clause, and the asymmetry with
\textup{(a)} is the rule working rather than two conventions: one range was strict and one is weak
because one boundary is attained inside the model and the other is not.}
\textup{(c)} $\tau_{\mathrm{DP}}$. No referent anywhere depends on its value ---
\cref{def:lfstep}\textup{(c)} is total at every $\tau_{\mathrm{DP}}>0$ by its emptiness convention ---
so the referent boundary is all of $(0,\infty)$ and the \emph{whole} of $\tau_{\mathrm{DP}}>2$ is a
clause. It is carried on a hypothesis line, \cref{prp:phirate}\textup{(P-e)}, correctly and already.
\emph{The rule is stated for analyst-chosen constants and for no others}; a numeric range on a
\emph{functional of $P$} is a restriction on $P$ and was never in question.

\emph{Why the split, and not the negation test it replaces.} Through 2026-08-25 the printed discriminant
read: negate the fragment, and if statements stop parsing it is a declaration. \emph{That test is
refuted inside this corpus.} A4(i) was written existentially, so negating it left $\varepsilon_\varrho$
with no referent and stopped \eqref{eq:trunset}, $\widehat{\widetilde W}$, \eqref{eq:hatT} and the
results built on them from parsing --- among them \suppref{Proposition}{hatAprop},
\suppref{Proposition}{tccfails}, \suppref{Proposition}{tccrepair}, \suppref{Lemma}{opnoise},
\suppref{Lemma}{aerr}, \suppref{Lemma}{unifop}, \suppref{Proposition}{stopfeas},
\cref{prp:phirate} and the conditions grouped as \textup{(T-5)}, and \emph{no claim of exhaustiveness is
made about that list}. The old test therefore ruled \emph{overlap} a declaration carrying no hypothesis,
while the assumption block carries it as a \emph{condition} on the hypothesis line of every result whose
proof reads it --- the lower half of \suppeq{wbounds} among them --- correctly.
\emph{And the verdict turned on an unstated reading convention}: read ungenerously at
A0 it gave ``declaration'', read generously at A8's basis clause --- which also \emph{introduces} the
$b_j$ that \cref{def:lfstep} needs --- it gave ``carrier'', and the two readings stood three lines apart.
Splitting first removes both defects, because it is the \emph{remainder} that is tested and the
remainder always has its symbols.

\emph{The word \textbf{clause} covers a restrictive fragment whether or not it is numbered.} A8's basis
clause is one and is correctly carried on \suppref{Lemma}{unifop}'s line; B1(iv)'s stem is not one and is
carried nowhere. \emph{And a result whose proof opens only some clauses under a tag cites those clauses:}
the stem is not among them, because the stem was never optional, and a \emph{block} citation instead
would charge the result with clauses its proof never opens --- the ground on which
\suppref{Corollary}{elldiff} is cited by clause, stated once here instead of six times at the sites.

\emph{The two positions this paragraph leaves out, and what each does.} \emph{At a proof step}: naming a
tag inside a proof records where the condition is consumed and charges nothing by itself --- the
hypothesis line does the charging --- so a clause consumed at a step and named on no line is
\emph{under-charged}, and that is the defect repaired at eight results on 2026-08-24 and at both steps of
\suppref{Lemma}{aerr}. This register is not optional to state: it is the one in which the under-charging
rule is defined. \emph{Inside another assumption}: A2$'$ opens ``A2 holds, and \dots'', so a result
citing A2$'$(i.a) carries A2 transitively, and a census run over hypothesis lines alone does not see the
edge.

% ------------------------------------------------------------------------
% ADDED IN DISPATCH 25 (2026-08-08), PHASE 1 OF THE MULTINOMIAL-INSTRUMENT PLAN.
% A1 HAD NO PARAGRAPH HERE, and it did not need one while the instrument was
% binary: the statement was read as self-explanatory and the representation was
% invisible because there was only one of them.  With sans-L cells the
% representation has to be exhibited and shown exhaustive, and that is what this
% paragraph does.  IT PROVES THREE THINGS AND ASSUMES NOTHING NEW.  It is
% placed FIRST because it is read before A1b, and because every later paragraph
% that says "the contrasts" points here for what that means.
\paragraph{A1 (conditional exogeneity).}

A1's statement is cardinality-free, and this paragraph is the reason the passage from a binary to
an $\mathsf L$-category instrument is a change of representation and not of content. A1 says
$\Delta U\upmodels W\mid\Delta Z$, a statement about $\sigma$-fields; nothing in it refers to how many
values $W$ takes. What referred to the cardinality, in every round before the $\mathsf L$-category one,
was the single scalar contrast $\widetilde W=W-\varrho(\Delta Z)$ through which the restriction was
carried. Three facts fix the general representation, and all three are established here rather than
assumed.

% THIS PARAGRAPH POINTS FORWARD TO eq:omegainv AND eq:omegaprefac, WHICH ARE
% DERIVED IN THE A4 PARAGRAPH A FEW PAGES BELOW IN THIS SAME SECTION.  The
% forward pointer is deliberate and is not the defect def:spaces's 2026-08-03
% note records: nothing is DECLARED here and used earlier, and there is no
% circularity -- neither of the two derivations uses anything from this
% paragraph.  A1 is placed first because the Discussion follows the order of the
% assumption block and A1 precedes A1b and A4 there.
First: the contrasts span the conditionally centred functions, exactly. Fix $z$ at which
$\varrho_l(z)>0$ for every $l\in\{0,\dots,\mathsf L-1\}$, which A4(i) supplies for almost every $z$. Let
$a$ be any measurable function of $(W,\Delta Z)$ with $E[|a|\mid\Delta Z=z]<\infty$ and
$E[a\mid\Delta Z=z]=0$, and write $a_l(z):=a(l,z)$, a finite list of numbers at each $z$, since $W$
takes finitely many values, so that
the centring is $\sum_{l=0}^{\mathsf L-1}a_l(z)\varrho_l(z)=0$. Put $c_l(z):=a_l(z)-a_0(z)$ for
$l=1,\dots,\mathsf L-1$. Then
\[
\sum_{l=1}^{\mathsf L-1}c_l(z)\,\un\{W=l\}
\;=\;a(W,z)-a_0(z),
\qquad
\sum_{l=1}^{\mathsf L-1}c_l(z)\,\varrho_l(z)
\;=\;\Big(\sum_{l\ge1}a_l\varrho_l\Big)-a_0\big(1-\varrho_0\big)
\;=\;-a_0(z),
\]
the second using the centring in the form $\sum_{l\ge1}a_l\varrho_l=-a_0\varrho_0$. Subtracting,
$c(z)'\widetilde W=a(W,z)$. So every conditionally centred function of $(W,\Delta Z)$ is a
$\Delta Z$-measurable linear combination of the $\mathsf L-1$ contrasts, and the family is exhaustive
with nothing approximated. The representation is unique: if $c(z)'\widetilde W=0$ almost surely given
$\Delta Z=z$ then $0=\operatorname{Var}(c'\widetilde W\mid\Delta Z=z)=c(z)'\Omega(z)c(z)$, and
$\Omega(z)\succ0$ forces $c(z)=0$, \emph{$\Omega(z)\succ0$ being supplied by \textbf{the standing
hypothesis of this paragraph}, that every cell is non-null at $z$, through the determinant identity
$\det\Omega=\prod_{l}\varrho_l$ below, and \textbf{not} by \eqref{eq:omegafloor}}, which is A4(i)'s
global consequence and is not assumed here. The space therefore has dimension
exactly $\mathsf L-1$ at such a $z$ and \eqref{eq:contrast} is a basis for it.
% ROUND-5 CORRECTION.  The uniqueness step cited eq:omegafloor, which OVER-
% ATTRIBUTES: eq:omegafloor is A4(i)'s consequence and carries the uniform floor
% eps_Omega, whereas all this step needs is that Omega(z) be nonsingular AT THE z
% the paragraph has already fixed.  Citing the global bound here is what would
% make a reader think the spanning statement needs A4(i) -- exactly what the next
% paragraph exists to deny.

At a degenerate $z$ the count is different, and it is stated here because the notation block
promises it and because a reader who stopped at the display above would conclude that spanning needs
A4(i). \emph{It does not.} Let $A_0:=\{l\in\{0,\dots,\mathsf L-1\}:\varrho_l(z)>0\}$ be the non-null
cells. The contrasts indexed by null cells are identically zero at $z$, and the argument above, re-run
with the sums restricted to $A_0$, shows that the conditionally centred functions of $(W,\Delta Z)$ at
$z$ are exactly the $\Delta Z$-measurable combinations of $\{\widetilde W_l:l\in A_0,\,l\ge1\}$, of
dimension $\#A_0-1$ when the base cell is non-null, and $\#A_0-1$ again when it is not, the
constraint $\sum_{l\in A_0}a_l\varrho_l=0$ removing one dimension from an $\#A_0$-dimensional space in
both cases. \emph{So the spanning statement holds at every $z$, with the dimension being one less than
the number of non-null cells}; A4(i) makes $\#A_0=\mathsf L$ everywhere and removes the case
distinction, but it is not a hypothesis of the statement. \emph{This is why no
Bierens-type device is needed in the $W$ direction}: a Bierens family exists to span an
infinite-dimensional space of test functions by a countable or continuum-indexed family, and here the
space is finite-dimensional and a basis is available in closed form. The Bierens family remains where it
was, over $\Delta Z$, and B1 is unchanged apart from B1(iv.d).

Second: nothing above depends on which cell is the base. The derivation used no property of cell
$0$ beyond its being one of the cells, so re-running it with base $k$ produces a second basis of the
\emph{same} space; the two are related by an invertible $\Delta Z$-measurable linear map, and every
statement of the form ``the restriction holds'' or ``the span is exhausted'' is therefore invariant to
the choice. \emph{Objects that are not invariant are flagged where they occur, and the distinction is
finer than it looks}: the individual contrast $\widetilde W_l=\un\{W=l\}-\varrho_l(\Delta Z)$ does not
mention the base cell at all, but \emph{which} $\mathsf L-1$ of the $\mathsf L$ contrasts are collected
into the vector does, the base cell is the omitted one, and $\Omega$, being the covariance of that
collection, moves with it. The quadratic form $\widetilde W'\Omega^{-1}\widetilde W$ does not move: by
\eqref{eq:omegainv} it equals $(1-\varrho_W)/\varrho_W$, which names no cell but the realised one. See
the paragraph \emph{A4}.

\emph{Third: restriction \textup{(A)} at the $\mathsf L-1$ contrasts is equivalent to A1, not merely
implied by it, \textbf{and the equivalence uses no overlap condition.}} For bounded measurable
$\varsigma$ the computation of \eqref{eq:omegaprefac} with
$\varsigma(\Delta U)$ in place of $\un\{\Delta U\le u\}$ gives
$E[\widetilde W\varsigma(\Delta U)\mid\Delta Z=z]=\Omega(z)\,d_\varsigma(z)$ with
$d_{\varsigma,l}(z)=E[\varsigma(\Delta U)\mid W=l,\Delta Z=z]-E[\varsigma(\Delta U)\mid W=0,\Delta Z=z]$,
read at any versions of the conditional expectations, those at null cells being arbitrary. \emph{Fix
$z$ and put $A:=\{l\ge1:\varrho_l(z)>0\}$}; $\Omega(z)$ vanishes off $A\times A$ and its $A$-block is
$\Omega_A=\operatorname{diag}(\varrho_A)-\varrho_A\varrho_A'$ with
$\det\Omega_A=\varrho_0\prod_{l\in A}\varrho_l$. If $\varrho_0>0$ then $\Omega_A$ is nonsingular and
$\Omega d_\varsigma=0$ forces $d_{\varsigma,A}=0$; if $\varrho_0=0$ then
$\ker\Omega_A=\operatorname{span}\{\mathbf 1_A\}$, from $\varrho_l(c_l-\varrho_A'c)=0$ with
$\varrho_l>0$ on $A$, and $d_{\varsigma,A}$ is constant on $A$. \emph{In both cases the restriction
holds for every bounded $\varsigma$ if and only if the conditional law of $\Delta U$ given
$(W,\Delta Z=z)$ is the same across all non-null cells}, which is $\Delta U\upmodels W\mid\Delta Z=z$.
\emph{No condition on $\varrho$ was used, and A4(i) is \textbf{not} a hypothesis of this equivalence}:
what A4(i) is needed for is \eqref{eq:omeganorm}, the normalisation, which is a different matter and is
taken up in the paragraph \emph{A4}. \emph{At $\mathsf L=2$ this is the equivalence the block has
used since round~1}, with the scalar $\varrho(1-\varrho)$ in place of $\Omega$, and the two cases
collapse to the single observation that a null cell makes $W$ conditionally degenerate.

What extra cells buy, and what they do not. They do not buy identification, and no sentence
in this corpus may suggest that they do. Identification of $\vdagger$ is A1b(iii), and A1b(iii) at the
full $\mathsf L$-category $W$ is \emph{implied by} A1b(iii) at any coarsening of $W$ \emph{together with
A1}. The two halves are separate and only the first is about coarsening. \emph{Inclusion}: a coarsening
$g(W)$ imposes fewer restrictions, so $\mathcal V_{CI}(W)\subseteq\mathcal V_{CI}(g(W))$, and
$\mathcal V_{CI}(g(W))=\{\vdagger\}$ gives $\mathcal V_{CI}(W)\subseteq\{\vdagger\}$, \emph{at most}
one element. \emph{Non-emptiness}: $\vdagger\in\mathcal V_{CI}(W)$ is exactly
$\Delta U\upmodels W\mid\Delta Z$, which is A1 and is not delivered by the coarsening. \emph{With both,
$\mathcal V_{CI}(W)=\{\vdagger\}$.} A first draft of this paragraph gave the inclusion alone and called
it the implication, which would leave $\mathcal V_{CI}(W)$ possibly empty. \emph{A finer
instrument therefore makes the identification hypothesis weaker, never stronger}, A1 being maintained
throughout, and identification
already held at $\mathsf L=2$.

What the extra cells \textbf{do} buy is \textbf{localisation}, and three concepts have to be kept
apart to see it.

First, testability, which is \textbf{present at $\mathsf L=2$} and is not what $\mathsf L>2$ adds.
At $\mathsf L=2$ the restriction is $d_1(u\mid z)=0$ \emph{for every $u$ at almost every $z$}, a
continuum of equations against one unknown function. This is a conditional-\textbf{independence}
restriction with a continuously distributed conditioning variable, not the mean-independence case in
which a binary instrument supplies two moments and exhausts them. A1 is therefore testable at
$\mathsf L=2$, and the paper's own application demonstrates it: the Section~6 probe rejects A1 at
$\mathsf L=2$. A statement that more cells make A1 testable is false and must not be written.

Second, instrument strength, which is a property of the data and not of the model. How hard $W$
moves $X$ bears on precision. \emph{Splitting an instrument creates no variation}, and strength within a
sub-population can be examined at $\mathsf L=2$ by subsetting. This is not what $\mathsf L>2$ adds
either.

Third, and this is the claim, \textbf{localisation}. At $\mathsf L=2$ the restriction is the
single object $d_1(u\mid z)=0$, so a violation concentrated in one channel of the instrument and its
absence in another are \emph{averaged into one verdict}: the test rejects, and cannot say through which
channel. At $\mathsf L\ge3$ the restriction is $d_l(u\mid z)=0$ \emph{cell by cell}, so a violation can be
\emph{attributed} to a cell. \emph{That question cannot be posed at $\mathsf L=2$ at all}, and posing it
is what the finer instrument is for.

The formal apparatus is joint consistency, and it is apparatus rather than motivation. Under A1
each of the $\mathsf L-1$ contrast directions carries a valid restriction on the \emph{same} $\vdagger$
, each a whole family, indexed by the test function and by the Bierens index $t$, so the testable
content is that there exists a single $\varphi\in\mathcal V$ at which all $\mathsf L-1$ systems hold, and
a $J$-type statistic for that hypothesis is the vehicle. \emph{This route needs no identification from
any subset}, which is why it is the one taken; and its per-cell decomposition is what delivers
localisation. \emph{A route a first draft took is withdrawn}: it pinned $\vdagger$ from a subset of the
directions and called the remainder testable, which presumes identification from a \emph{singleton}
contrast, and contrast $l$ alone is the binary coarsening $\un\{W=l\}$, so identification from it is
A1b(iii) at that coarsening, \emph{stronger} than A1b(iii) at the full instrument and not implied by it.

The application is what the split exists for, and it makes the three concepts concrete. The
Section~6 instrument divides into a \emph{chronic-level} channel, the annual-mean-only and
both-triggers cells, and a \emph{transient} channel, the bad-day-only cell. \emph{The hypothesis the
split is built to test is that A1 fails through the chronic channel}, where $W$ is close to a function of
$X_1$, the endogenous regressor's own initial value, \emph{and not through the transient one}. At
$\mathsf L=2$ that hypothesis has no expression. \emph{The proposition that states the joint-consistency
restriction and the statistic that localises a violation to a cell are not in this dispatch}; they are
recorded as owed, and Phase~2 builds them off this paragraph.

\emph{\textbf{And consistency of every cell set is \textbf{necessary} for A1, not sufficient, which
governs how a non-rejection may be reported and is stated here because it is stated nowhere else.}}
\Cref{prp:localise}\textup{(ii)} runs in one direction only: A1 implies that every cell set is
consistent, and its contrapositive is that a conflict refutes A1. The converse is never asserted,
and it is false in general. The reason is structural rather than statistical, and survives at infinite
sample size. What the cell system interrogates is whether $\mathcal V_{CI}$ is \emph{empty}; what A1 is
about is whether $\vdagger$ \emph{lies in} $\mathcal V_{CI}$. \emph{Those are different questions.} At a
law where A1 fails at the truth but some other $\varphi_1\in\mathcal V$ satisfies the
conditional-independence restriction, every cell set is satisfied \emph{at $\varphi_1$}, every conflict
set is empty, the diagnostic reports nothing, and the estimator converges to $\varphi_1\ne\vdagger$.
Nothing in this block excludes such a law: A1b(iii) would also fail there, it is the clause that says
$\mathcal V_{CI}=\{\vdagger\}$, and the point is precisely that \emph{the diagnostic cannot see the
failure of the clause it would need to see}.
\emph{So a cell showing no conflict is \textbf{not exonerated}}: the procedure detects
$\mathcal V_{CI}=\varnothing$ and nothing more, which is not an interpretation but
\cref{prp:localise}\textup{(iii)}'s own equivalence, that some consistent $\mathsf S$ has a non-empty
conflict set \emph{if and only if} $\mathcal V_{CI}=\emptyset$, so a violation reconciled by some other
$\varphi\in\mathcal V$ is invisible to it, and the false-negative rate of the attribution is
unpriced by anything in this paper. \emph{Section~6 will attribute a non-rejection to a channel and must report it in exactly those
terms}, as the absence of a detected conflict, not as evidence for A1 in the un-flagged cell.
% ADDED 2026-08-10 (Phase 6a).  THIS IS NOT A COUNTEREXAMPLE and is not recorded as
% one: prp:localise(ii) states only the forward implication, all its hypotheses and
% its conclusion hold, and a probe that constructed such a law (L = 3, varrho_l =
% 1/3, a perturbation orthogonal to the first cell contrast but separating the
% conditional laws at higher moments) found every hypothesis satisfied.  What the
% probe found was that NO PRINTED SENTENCE ANYWHERE IN THE CORPUS says the converse
% is not claimed.  rem:localise(5)'s "What is not claimed" blocks two OTHER
% over-readings -- it is not this one -- and a search for any statement that
% consistency is necessary and not sufficient returned zero hits.  The argument
% above is stated at the level that is verified here (the two questions differ, and
% the diagnostic answers the wrong one), rather than by reciting a construction this
% file does not carry.
% ROUND-5 CORRECTION, AND IT REPLACES THE MOTIVATION AND NOT ONLY THE WORDING.
% Rounds 1-4 said the extra cells "convert A1 from an untestable assumption into
% one with evidence for or against it", and counted "sans-L - 2 independent
% testable restrictions".  BOTH ARE STRUCK.  A1 IS TESTABLE AT sans-L = 2 -- the
% restriction is a continuum of equations against one function, and Section 6's
% own probe rejects it there -- so the sentence was false, and the count was a
% moment-counting reflex imported from the mean-independence case where a binary
% instrument gives two moments.  It does not apply to a conditional-independence
% restriction with a continuous conditioning variable.
% THE THREE CONCEPTS THE FILE MUST NOT CONFLATE ARE TESTABILITY, INSTRUMENT
% STRENGTH AND LOCALISATION.  Only the third is what sans-L > 2 buys.  A later
% dispatch tempted to restore "overidentification" as the headline should note
% that the joint-consistency statistic IS an overidentification statistic -- that
% is not the error -- and that the error was calling its existence, rather than
% its per-cell decomposition, the reason for the design.

% *** AND THE MECHANISM IS NOT THE ILL-POSED INVERSE PROBLEM.  READ THIS BEFORE
% WRITING ANYTHING ABOUT WHAT MORE CELLS BUY. ***
% varphi solves the NONLINEAR operator equation eq:opeq under a conditional
% INDEPENDENCE restriction.  A_cond is the DERIVATIVE of T at the truth, not the
% problem; the degree of ill-posedness a of H1 is a property of A_cond's singular
% values, and the cardinality of W does not act on it in any way this block can
% state, let alone in the way "more instruments improve conditioning" would
% suggest for a LINEAR Fredholm equation.  THE TEMPTATION IS REAL AND THE
% SENTENCE WOULD LOOK PLAUSIBLE: it is the sentence a reader who has the linear
% NPIV literature in mind will supply, and it imports a different problem.  What
% extra cells do is add MOMENT DIRECTIONS -- restrictions on the same varphi --
% which is an efficiency and overidentification statement and lives entirely on
% the moment side.  Recorded here as a standing prohibition because this is the
% paragraph a later dispatch will edit when it writes the overidentification
% result.

% ------------------------------------------------------------------------
\paragraph{A1b (identification).}

\emph{(The corresponding statement for a generic $\varphi\in\mathcal V$ is \cref{def:spaces} and not an
assumption; only membership of the truth is content.)}
Mode of consumption, stated once and binding on every citation of this clause: A1b(i) is consumed
for \textbf{well-posedness} and never for a bound. What it buys is that the truth, and hence the path
$\vdagger+\tau_1h$ for $h\in\mathcal V$, lies in the parameter space. A result whose conclusion
concerns a generic $h\in\mathcal V$, and whose proof needs only $\int h\,d\pi=0$, consumes
\cref{def:spaces} and must cite \cref{def:spaces}; citing A1b(i) there charges the result with a
hypothesis it does not use, because the two statements are typographically identical and semantically
disjoint. The genuine consumers are \cref{lem:secondord} \textup{(Step 1)}, \cref{prp:orthog}\textup{(i)},
\cref{cor:udep} \textup{(clause (a), at $\vdagger$)}, \cref{lem:biasexp} and, transitively,
\cref{lem:solset} and \cref{thm:rootn}.
\emph{\Cref{lem:biasexp}'s use is the archetype}: A1b(i) is what puts
$h_k=\hat\varphi^{(-k)}-\vdagger$ in $\mathcal V$, hence what makes the balancing pairing
$\langle\nu_{g_k}/\pi,h_k\rangle_{L^{2}(\pi)}$ of \eqref{eq:biasdecomp} the object
\cref{lem:nusigma}\textup{(iii)} computes and the object A2 asserts a range condition about, a
\emph{well-posedness} use and not a bound. \Cref{thm:rootn} carries
A1b(i) in group \textup{(T-1)} for the same use, inherited.
\emph{\Cref{lem:prelimtheta} Step~4, \cref{lem:epterm} Step~3 and \cref{lem:sandwich} Steps~3 and~6$'$
are \textbf{not} consumers}, although each applies \cref{lem:nusigma}\textup{(i)} at $h_k$:
that clause is stated for every $h\in L^{2}(\mathcal X,\pi)$, which is all its proof establishes.
% ARCHAEOLOGY.  lem:biasexp was missing from this binding list through dispatch 13,
% which is the dangerous direction: the list was wrong and the citation was right,
% so a reader following the list would have struck a correct citation.  Both the
% assumption-auditor and the regularity-checker reached that conclusion
% independently, the latter classifying it as "a citation gap and not a mathematical
% one".  thm:rootn -- the paper's only theorem -- was missing from the same list
% through dispatch 17, the same defect one level up; it survived the dispatch-14
% repair because that repair ran site-by-site rather than list-by-list.  No other
% clause in this block carries a list declared binding; A1d's "sole consumer
% cor:a1cprim" and A2c's "exactly one consumer lem:reprsmooth" are the only
% comparable claims and both were re-checked and are correct.  The three
% non-consumers would each have needed A1b(i) under lem:nusigma(i)'s round-13
% statement "for every h in V".

% ------------------------------------------------------------------------
%  MOVED HERE FROM assumptions_main.tex, 2026-08-28, ON THE AUTHOR'S SHORTENING
%  INSTRUCTION, ON THE A3 / A6 / C4' PRECEDENT AND ENTIRE.  A1d printed 357 words
%  IN THE ARTICLE -- a struck-through retirement carrying a full proof of its own
%  unsatisfiability -- and a retirement record is working-corpus material.  The
%  tag, the \label ass:instrcomp, the struck banner and every word of the ground
%  are preserved here, so a stale citation still lands on the reason and not on a
%  gap.  NOTHING ELSE MOVED: the assumption environment is unnumbered
%  (\newtheorem*{assumption}), so no counter anywhere shifts.  The \cref to
%  rem:a1cwithdrawn resolves here because the supplement inputs results_supp.tex,
%  which is where that remark now lives.
% ------------------------------------------------------------------------
\paragraph{A1d (RETIRED: instrument completeness given the residual), moved entire.}

\begin{assumption}[A1d: Instrument completeness given the residual]\label{ass:instrcomp}
\textbf{\sout{RETIRED, 2026-08-05.}} It is unsatisfiable in this model. The condition read:
for every $\Lambda\in L^{2}(P_{X_1,X_2})$,
$E[\widetilde W\,\Lambda(X_1,X_2)\mid\Delta U,\Delta Z]=0$ a.s.\ implies
$\Lambda(X_1,X_2)=0$ a.s.\ \sout{[retired]}. \emph{Take $\Lambda\equiv1$}, which lies in
$L^{2}(P_{X_1,X_2})$ because $P$ is a probability measure. Since
$\widetilde W_l=\un\{W=l\}-\varrho_l(\Delta Z)$ by \eqref{eq:contrast} and A1 gives
$\Delta U\upmodels W\mid\Delta Z$,
\[
E\big[\widetilde W_l\,\big|\,\Delta U,\Delta Z\big]\;=\;E\big[\un\{W=l\}\,\big|\,\Delta Z\big]-\varrho_l(\Delta Z)\;=\;0
\quad\text{a.s., for every }l ,
\]
so the antecedent holds while $\Lambda\ne0$. Hence \emph{A1 implies the negation of this condition},
and so does A1b(iii), which delivers A1 at $\vdagger$. More generally, $\widetilde W$ is centred
given $\Delta Z$, so $\Lambda\mapsto E[\widetilde W\Lambda\mid\Delta U,\Delta Z]$ annihilates every
$\Lambda$ agreeing a.s.\ with a $\sigma(\Delta U,\Delta Z)$-measurable function: no completeness
condition on that operator over a class containing the constants can hold here. That is what separates
this design from the conditional-mean one of \citet{feve2014}, whose operator $g\mapsto E[g\mid W]$ does
not annihilate constants, \emph{their hypothesis transplants as a contradiction, not as an excess}.
\suppref{Lemma}{nusigma}\textup{(v)} is the unconditional shadow of the display above and was in this
corpus throughout.

\emph{The status of A1d, written with A1d as its subject.} \textbf{A1d} is not maintained, is not
available for use, and is a hypothesis of no result in this programme; the tag is kept rather than
reused, on the A3 and B1(vi) precedent, so that a stale citation lands on the reason rather than on a
gap. \emph{Its consumer count is zero}: the corollary that derived A1c from it, \texttt{cor:a1cprim},
was withdrawn with it and no other statement takes it as a hypothesis; see \cref{rem:a1cwithdrawn}, which
carries the withdrawal and says so in terms. \emph{And A1d is \textbf{retired}, on A6's ground and not on
A3's}: it is \emph{false} in this model rather than \emph{proved} or merely unconsumed, which is the
distinction the five-clause list of unconsumed-but-live conditions below turns on. \emph{This
block therefore has \textbf{five} retirements and not four}, on \textbf{three} grounds --- A3 and
A$2'$(ii) because they are proved, A6 and A1d because they are false, B1(vi) because it is about a
different estimator --- and the census below counts five.
\end{assumption}

% ------------------------------------------------------------------------
\paragraph{B1 (Bierens family and range space).}

\begin{enumerate}
\item[(vi)] \textbf{\sout{RETIRED.}} It describes an estimator this block
does not implement. The clause read: \emph{the family used on the $\varphi$-side is a finite set
$\mathcal T_n\subset\mathcal T$ with $\#\mathcal T_n\to\infty$ and
$\overline{\bigcup_n\mathcal T_n}\supseteq\mathcal T$} \sout{[retired]}. It is not maintained, is not
available for use, and is a hypothesis of no result in this programme. The clause number is kept
rather than reused, on the A3 and \textup{(P-b)} precedent, so that a stale citation lands
somewhere.
That second claim was false for four dispatches, and what made it false is recorded here rather than
repaired silently. A retirement is complete only when the clause is struck, when no hypothesis list carries
it \emph{including through a block citation of the parent tag}, and when nothing describes it in the present
tense as a live condition. B1(vi) failed the second test at \cref{prp:phirate}, whose list read ``B1'' in full
and now reads B1(i)--(v); and it failed the third at three compiled sites, the Discussion of this block,
\cref{rem:Tcond} in \texttt{results.tex}, and that remark's rendering in
\texttt{paper/sections/theorems/lem-bierens.tex}, each of which called it an identification restriction in
the present tense. All four are closed. \emph{A clause that is struck in one file and live in three is worse
than a clause nobody struck}, because the strike is what stops the next reader from checking.
Why it is retired, and this one is retired for the third of the three reasons, neither
``it is a theorem'' \textup{(A3)} nor ``it is false'' \textup{(A6)} but ``it is about a different
estimator''. The \textbf{growing finite family} $\mathcal T_n$, a finite subset of $\mathcal T$
with $\#\mathcal T_n\to\infty$, is instantiated \textbf{nowhere in this project}, in the compiled text
or in the simulation code. The
estimator does not use a growing finite grid: \cref{def:hatT} integrates against $\Pi_t$ over \emph{all} of
% ARCHAEOLOGY.  The sentence read "cal T_n occurs nowhere in this project outside the
% sentence above" through dispatch 20 and was FALSE on the register in which it was
% written: the varphi-step simulation carries an identifier Tn for a FIXED quadrature
% node set -- a constant number of Halton nodes for a Pi_t uniform on cal T, whose
% own header says in terms that it is B1(iii)--(iv.b)--(iv.c) and NOT B1(vi) -- so
% the STRING occurred and the OBJECT did not.  The claim is restated at the object
% rather than at the string, as of dispatch 21.
$\mathcal T$, and every operator norm in Parts~F--H is an $L^{2}(\Pi)$ norm over the whole index set. The
clause is a leftover from a Dom\'inguez--Lobato-style design with a growing finite family, which this block
considered and did not adopt.
What it was for, and where that work is actually done. Its heading claimed identification of
$\vdagger$. Identification of $\vdagger$ in this block is \textbf{A1b(iii)}, assumed outright, and A1b(iii)'s
own comment records that no primitive route to it is offered. A growing family would be one ingredient of
such a route, \emph{and the counterexample-hunter's sub-probe shows it would be a genuine ingredient and
not decoration}: at a \emph{fixed finite} grid $\Pi_t=\sum_{k\le K}\delta_{t_k}$ with $\Delta Z$ supported on
$K+1$ atoms, the $K\times(K+1)$ matrix $[e^{t_k\Phi(z_j)}]$ has a null vector, and perturbing along it makes
the Bierens family hold while the conditional restriction fails. So the clause is correctly labelled an
identification condition; what it is not is a condition of this estimator. It is recorded here rather than
deleted so that a later dispatch reinstating the growing-grid design finds the statement and the witness
together. \emph{Note also that (vi) was never sufficient for A1b(iii)}: it constrains the index family and
says nothing about $\mathcal V_{CI}$.
\end{enumerate}

% ------------------------------------------------------------------------
\paragraph{A2$'$ (source condition with usable representer).}

% MOVED HERE FROM INSIDE THE STATEMENT OF A2', 2026-08-28, ON THE AUTHOR'S
% SHORTENING INSTRUCTION.  The TYPING -- Euclidean over the p coordinates of
% theta, maximum over the contrast index l -- STAYED in the statement, because it
% fixes what the two budgets mean.  What moved is the ARGUMENT for the asymmetry,
% which is a derivation and not a condition.
\emph{Why the two budgets aggregate in two different ways.} \textbf{Euclidean over the coordinates of
$\theta$} because every consumer of $\Lambda^{(0)}_\star$ and $\Lambda^{(1)}_\star$ reads a Euclidean or
Frobenius norm of the assembled object --- \cref{cor:elldiff}'s two sup-norms and everything downstream
of them --- so a coordinatewise maximum would leave every one of those constants short by $\sqrt p$.
\textbf{Maximal over $l$} because \suppeq{Wdef} treats the contrast index diagonally, each slice
generating its own column and no other, so no downstream bound ever sums across $l$.

\begin{enumerate}
\item[(ii)] \textbf{\sout{RETIRED.}} It is a consequence of \textup{(i)}.
The clause read
$E[\widetilde W^{2}\|\ell_\star\|^{2}]<\infty$ \sout{[retired]} and is not maintained, is not
available for use, and is a hypothesis of no result in this programme. The clause number is kept
rather than reused, on the A3 precedent and for the same reason \textup{(P-b)} is kept: a stale
citation must land somewhere.
Why it is retired, and it is retired because it is a \textbf{theorem}, exactly as A3 was, not
because it is false and not because it is too strong. Under (i) in full, together with B1(i), (ii),
(iv.a), (iv.d), \cref{cor:elldiff} gives $\|\ell_\star\|_\infty\le\|\varpi_u\|_\infty\Lambda^{(0)}_\star<\infty$,
and $\|\widetilde W\|\le\min\{\sqrt2,\sqrt{\mathsf L-1}\}$, in the Euclidean norm on
$\IR^{\mathsf L-1}$, is \textup{(P1b)}, free from \eqref{eq:contrast}, $W$ taking finitely many values and
$\varrho_l\in[0,1]$; hence
$E[\widetilde W^{2}\|\ell_\star\|^{2}]\le2\|\ell_\star\|_\infty^{2}<\infty$ with no further hypothesis.
\emph{The factor $2$ is the square of that bound and is therefore \textup{(P1b)}'s constant and not
\textup{(P1)}'s}: \textup{(P1)} is the
\emph{per-cell} $|\widetilde W_l|\le1$ and does not bound the Euclidean norm of the vector, and a
derivation routed through it here would produce the constant \cref{cor:rangegap}'s proof records as
\textbf{false} at $\mathsf L\ge3$. \emph{The bound is written out rather than called by its symbol},
that symbol being declared at \cref{cor:rangegap}, which this file precedes in reading order.
\emph{This tag was \textup{(P1)} through 2026-08-24, at both sites of
this retirement ground}, while \cref{lem:momzero}'s closing paragraph, which runs the identical
derivation, printed ``the bound on $\|\widetilde W\|$ is the Euclidean one and \textbf{not}
\textup{(P1)}''; \emph{soundness was never at stake, but a retirement ground is a derivation and must
name every clause it uses}. The factor is $1$ at
$\mathsf L=2$; it is a constant, it is uniform in $\mathsf L$, and it changes nothing about the
retirement.
\emph{The attribution is (i) \textbf{in full} and not (i.a) alone}: the only route to
$\|\ell_\star\|_\infty<\infty$ in this block is \cref{cor:elldiff}, that corollary is stated at A2b, and
A2b requires (i.a) \emph{and} (i.b). \emph{Every tag in that derivation is already carried by every list
that could have carried (ii)}, so the retirement changes no conclusion anywhere.
\emph{The clause has \textbf{no consumer}}, as \cref{lem:balance}, \cref{lem:solset},
\cref{lem:secondord}'s own hypothesis line --- ``A$2'$ cited by clause; A$2'$(ii) is consumed by no
part of this lemma'' --- and \cref{lem:momzero}'s closing paragraph each state independently.
\emph{Four attestations, and the number is not maintained}: the draft this sentence replaces said
three and named three, and the fourth had been printing on its own hypothesis line throughout.
What survives, and it is the half the clause was really carrying. At Step~8
% ARCHAEOLOGY.  The round-8 change log of this file recorded "A2'(ii) is consumed by
% none of them and is a Step-8 clause".  A4's comment (d) justified retaining it by
% asserting that "the second-order and bias terms consume" it, which was false, and
% a maintained clause with a false retention ground is what an audit is for.
(\texttt{prp:effbound}), where the variance is computed and where tier~T3 relaxes (i), $\|\ell_\star\|_\infty$
is no longer available and $E[\widetilde W^{2}\|\ell_\star\|^{2}]<\infty$ becomes a genuine restriction, 
it is exactly the Khan--Tamer irregularity boundary, since holding the conditional laws given $W$ fixed and
driving any one cell probability to zero makes
$\|\Omega^{-1}\|\ge\max_{0\le l\le\mathsf L-1}\varrho_l^{-1}$ diverge by
\eqref{eq:omegainv}, read at the diagonal entry
$e_l'\Omega^{-1}e_l=\varrho_l^{-1}+\varrho_0^{-1}$, and with it
$E[\widetilde W^{2}\|\ell_\star\|^{2}]$. \emph{At $\mathsf L=2$ that is
the printed computation}, $\varrho\equiv\varepsilon$ giving
$E[\widetilde W^{2}\|\ell_\star\|^{2}]\asymp\varepsilon^{-1}$; \emph{at $\mathsf L\ge3$ it is
\textbf{one} cell that has to be thin and not all of them}, which is the reading the application needs,
its thin cell being a single trigger category.
\emph{That is a \textbf{forward} obligation and not a present consumer}, and it is recorded here rather than
carried as a maintained clause: Step~8 must state the condition it needs, at the tier where it needs it, and
under T3 it will not be this clause because (i) will not be available to make it redundant. \emph{Flagged
for Step~8}, on the model of \cref{rem:obligations};
\end{enumerate}

\medskip
\noindent
 Status of clause \textup{(i)}, stated here so that the tag names exactly one object. A$2'$(i) is a
\emph{maintained assertion about $\ell^\Pi_\star$, and about $\ell^\Pi_\star$ alone}. The same two
conditions read at a generic $\ell^\Pi\in L^{2}(\Pi)$ are not an assumption but a \emph{property} that a
candidate representer may or may not have; that property is tagged \textbf{A2b} below and is what
\cref{def:weight}(a) names \emph{admissible representer}. \emph{A result that needs the property at a
generic representer cites A2b and never A$2'$(i)}, exactly as a result needing the centring
property of a generic weight cites N1 and never a maintained clause. This is the
treatment N1/N1$'$ received in round 7 and it is applied here for the same reason: one tag, one object.

% ------------------------------------------------------------------------
\paragraph{A2c (source condition on the representer).}


Two things in \eqref{eq:lambdamu} are established here rather than assumed, because neither is
automatic and the second cannot be had from the transform.

\textup{(m.1)} The exceptional set is a Borel set by construction, and stating it that way is a
choice. The weaker phrasing ``for $\Pi_t$-almost every $t$ the slice lies in $H^{\mu''}(\IR;\IR^{p})$''
% ARITY OF THE QUOTATION CORRECTED 2026-08-10 (Phase C).  This quoted A2c's slice
% condition at the SCALAR H^{mu''}(R) while A2c states it at the VECTOR
% H^{mu''}(R;R^p) -- the same species as the p(L-1) miscount M1 closed one
% paragraph later, and the reason (m.1) was coherent but not literally faithful:
% the CONSTRUCTION below does run at the p(L-1) coordinate slices, and the
% paragraph says so, but the clause being weakened is the vector one.  One symbol;
% the argument that follows is arity-neutral and is unchanged.
delivers
the displayed form at no cost: it means that the set of $t$ failing the slice condition is contained in
some $\mathcal N\in\mathfrak B(\mathcal T)$ with $\Pi_t(\mathcal N)=0$, and
$\mathcal T_\ast:=\mathcal T\setminus\mathcal N$ is then Borel, carries full $\Pi_t$-measure, and
contains \emph{only} points at which the slice condition holds. \emph{Nothing anywhere requires
$\mathcal T_\ast$ to be the largest such set}, which need not be Borel. What the choice buys is that
$\IR\times\mathcal T_\ast$ belongs to $\mathfrak B(\IR)\otimes\mathfrak B(\mathcal T)$, the property
\cref{lem:reprsmooth}'s Step 2 needs when it asserts joint measurability there, and the property an
``almost every'' quantifier does not by itself supply.
With $\mathsf L$ cells the construction is run once per slice and the sets are then intersected.
\emph{What A2c asks for is \textbf{one} $\mathcal T_\ast$ serving every $l$}, its slice condition
being membership of the \textbf{vector} slice $\ell^\Pi(\cdot,t,l)$ in $H^{\mu''}(\IR;\IR^{p})$, so that
it has $\mathsf L-1$ slices and not more, \emph{while the construction below runs at the
$p(\mathsf L-1)$ \textbf{coordinate} slices and delivers that requirement rather than a stronger one}:
by the identity A2c prints when it extends \eqref{eq:usobolev} to $\IR^{p}$-valued $g$, membership of
the vector slice in $H^{\mu''}(\IR;\IR^{p})$ \emph{is} membership of all $p$ coordinates in
$H^{\mu''}(\IR)$, so the finer index is bookkeeping and not a further demand. There are
$p(\mathsf L-1)$ coordinate slices $\ell^\Pi_j(\cdot,t,l)$; take $\mathcal N_{j,l}$ for each and set
$\mathcal N:=\bigcup_{j,l}\mathcal N_{j,l}$, a \emph{finite} union, hence Borel and $\Pi_t$-null, and
$\mathcal T_\ast:=\mathcal T\setminus\mathcal N$ then carries full $\Pi_t$-measure and contains only
points at which \emph{every} coordinate slice condition holds, equivalently, by that same identity,
only points at which $\ell^\Pi(\cdot,t,l)\in H^{\mu''}(\IR;\IR^{p})$ for every $l$, which is A2c's
quantifier. \emph{Finiteness of $\mathsf L$ is what makes this
free}: a countable union would serve too, but a union over an index set of unbounded cardinality would
not, which is one more place where A0's ``$\mathsf L$ finite and fixed in $n$'' is doing work rather
than housekeeping. \emph{A per-slice $\mathcal T_{\ast,j,l}$ would not serve} \cref{lem:reprsmooth}'s
Step~2, which needs one product set on which all the slices are simultaneously well behaved.

\textup{(m.2)} The integrand is measurable, by pairing against one fixed countable family of
functions of $u$. There are real $\varkappa_1,\varkappa_2,\ldots\in L^{2}(\IR,du)$ with
\begin{equation}\label{eq:hdual}
\big\|g\big\|_{H^{\mu''}(\IR)}\;=\;\sup_{k\ge1}\Big|\int_{\IR}g(u)\,\varkappa_k(u)\,du\Big|
\qquad\text{for every real }g\in H^{\mu''}(\IR).
\end{equation}
Construction, with no negative-index space and no completion. Call a function of $\xi$
\emph{Hermitian} when its value at $-\xi$ is the complex conjugate of its value at $\xi$, which is what
transforms of real functions are. The Hermitian elements of $L^{2}(\IR,(2\pi)^{-1}d\xi)$ form a closed
real-linear subspace of a separable space, hence a separable real Hilbert space; take a countable subset
of its closed unit ball, dense there, and replace each member by its truncations to
$\{|\xi|\le R,\ |\cdot|\le R\}$ over rational $R$, truncation preserves Hermitian symmetry, because
the truncation set is symmetric, does not increase the norm, and converges in $L^{2}$, so the subset
may be taken to consist of \emph{bounded, compactly supported} Hermitian functions and to remain dense in
the ball. Let $\varkappa_k$ be the inverse Fourier transform of $(1+\xi^{2})^{\mu''/2}$ times the $k$-th
member: that product is bounded with compact support, hence in $L^{1}\cap L^{2}$, so $\varkappa_k$ is a
bounded continuous element of $L^{2}(\IR,du)$, and it is real because the product is Hermitian. For real
$g\in H^{\mu''}(\IR)$, Plancherel in the normalisation of \eqref{eq:usobolev} gives
\[
\int_{\IR}g\,\varkappa_k\,du
\;=\;\frac1{2\pi}\int_{\IR}\big[(1+\xi^{2})^{\mu''/2}(\mathcal Fg)(\xi)\big]\,
\overline{\big[(1+\xi^{2})^{-\mu''/2}(\mathcal F\varkappa_k)(\xi)\big]}\,d\xi ,
\]
whose second factor is the $k$-th member of the dense family and whose first is a Hermitian element of
$L^{2}(\IR,(2\pi)^{-1}d\xi)$ of norm exactly $\|g\|_{H^{\mu''}(\IR)}$, by \eqref{eq:usobolev}. The
supremum over $k$ of the modulus of a continuous linear functional over a dense subset of the closed unit
ball is the functional's norm, which is \eqref{eq:hdual}.

The remaining step is run at a fixed coordinate slice $(j,l)$ and then aggregated exactly as
\eqref{eq:lambdamu} aggregates: Euclidean over the $p$ coordinates, maximal over the $\mathsf L-1$
contrasts. \emph{The second stage is the only place the $\mathsf L$-category form changes anything
here}; the first is present at every $\mathsf L$, $\mathsf L=2$ included. At each fixed $k$ and
each $(j,l)$ the map $t\mapsto\int_{\IR}\ell^\Pi_j(u,t,l)\varkappa_k(u)\,du$ is $\Pi_t$-measurable
on $\mathcal T_\ast$: the product $(u,t)\mapsto\ell^\Pi_j(u,t,l)\varkappa_k(u)$ is
$\mathfrak B(\IR)\otimes\mathfrak B(\mathcal T)$-measurable, \emph{the contrast index $l$ ranges over
a finite set, so a jointly measurable representative of $\ell^\Pi$ on
$\IR\times\mathcal T\times\{1,\dots,\mathsf L-1\}$, with that set given its discrete $\sigma$-algebra, is
exactly a finite list of jointly measurable functions of $(u,t)$, one per $(j,l)$}, Tonelli on the
$\sigma$-finite product
$(\IR,du)\otimes(\mathcal T,\Pi_t)$ makes the $u$-integrals of its positive and negative parts measurable
in $t$, and both are finite at each $t\in\mathcal T_\ast$ because $\ell^\Pi_j(\cdot,t,l)$ and $\varkappa_k$
are both in $L^{2}(\IR,du)$. A countable supremum of measurable functions is measurable, so
$t\mapsto\|\ell^\Pi_j(\cdot,t,l)\|_{H^{\mu''}(\IR)}$ is $\Pi_t$-measurable on $\mathcal T_\ast$ for each
$(j,l)$.
One aggregation then remains, and it is the one over $j$: the integrand of \eqref{eq:lambdamu} is
the \textbf{vector} norm, which the preceding sentence does not by itself deliver. It closes on the
identity A2c prints when it extends \eqref{eq:usobolev} to $\IR^{p}$-valued $g$, read at
$g=\ell^\Pi(\cdot,t,l)$:
\[
\big\|\ell^\Pi(\cdot,t,l)\big\|_{H^{\mu''}(\IR;\IR^{p})}
\;=\;\Big(\sum_{j=1}^{p}\big\|\ell^\Pi_j(\cdot,t,l)\big\|^{2}_{H^{\mu''}(\IR)}\Big)^{1/2},
\qquad t\in\mathcal T_\ast ,
\]
a \emph{finite} sum of squares of $\Pi_t$-measurable functions followed by $z\mapsto z^{1/2}$, which is
continuous on $[0,\infty)$ and hence Borel; so
$t\mapsto\|\ell^\Pi(\cdot,t,l)\|_{H^{\mu''}(\IR;\IR^{p})}$ is $\Pi_t$-measurable on $\mathcal T_\ast$,
now for each $l$. Every summand is finite there, A2c placing the slice in $H^{\mu''}(\IR;\IR^{p})$ at
every $t\in\mathcal T_\ast$ and that membership being, by the same identity, membership of all $p$
coordinates. \emph{The step is one line because $p$ is a fixed finite dimension, exactly as
\textup{(m.1)}'s intersection was free because $\mathsf L$ is finite}; nothing here is a Fubini
statement in $j$.
Each of the $\mathsf L-1$ integrals in \eqref{eq:lambdamu}, \emph{one per contrast $l$, and
\textbf{not} one per $(j,l)$: the $p$ coordinates are aggregated inside the integrand, by the Euclidean
norm, so only the contrast index survives as an index of integrals}, is therefore well defined in
$[0,\infty]$, and so is their maximum, \emph{a maximum over a finite index set requiring no further
argument}. That is the quantity of which A2c asserts finiteness.
\emph{This is the same countable-supremum device \cref{lem:reprsmooth} uses at its Step 3 and
\cref{def:dirclass} at \eqref{eq:shiftnormone}, applied one level up}, and it is used \emph{instead}
of a Fubini statement about $(\xi,t)\mapsto(\mathcal F\ell^\Pi(\cdot,t))(\xi)$, which is not available,
for the reason \cref{lem:reprsmooth}'s Step 2 prints: on $L^{2}(\IR,du)$ the transform is an
$L^{2}$-limit and not an integral. \eqref{eq:hdual} pairs the slice against a family fixed \emph{once},
in the variable $u$ in which the slice is jointly measurable, so no such statement is needed.


% ------------------------------------------------------------------------
\paragraph{H1 (Hilbert-scale link).}

\emph{Two notational points.} \emph{First}, bare sans-serif $\mathsf c$ is this project's \emph{generic
finite positive constant}: bound by the quantifier of the clause in which it appears, free to take a
different value in every such clause, and never carried from one result to another. This clause binds it,
in the words ``there is a finite $\mathsf c$, not depending on $\zeta$'', and no site downstream reads its
value.
\emph{Second, the pre-image is written $h$ and \textbf{not} $\varsigma$}, $h$ being the
\emph{domain}-side bound variable \cref{rem:spacesread}\textup{(1)} declares, the counterpart of
$\Xi,\Xi'$ on the range side, and the pre-image being exactly that: a generic element of $L^{2}(\pi)$, existentially
bound per $\zeta$ inside this clause. \emph{The object is the spectral source pre-image}, the same species
as in (P-a.3) and \cref{lem:scaleinv}\textup{(iii)}, read here at the \emph{population} operator on all of
$\mathcal V$.
\emph{$H^{s}$ is the \textbf{classical} Sobolev space and \textbf{not} a scale generated by
$\mathcal D$}: for a Hilbert scale generated by a
differential operator on a bounded interval the second scale step imports the generator's boundary
conditions, so the space the generator produces at order $s'$ coincides with $H^{s'}(\mathcal X)$
only for
$s'<\tfrac32$ (Neumann-type) or $s'<\tfrac12$ (Dirichlet-type), \emph{$s'$ throughout this assumption
and A8's derivation denoting a \textbf{generic} order of the scale, bound where it appears
and used nowhere else in this project}, while \eqref{eq:lxwindow} forces
% ARCHAEOLOGY.  Through dispatch 21 the constant line read "the constant is the
% Sec.1.6 sans-serif slot and carries no name of its own"; that pointer was false --
% Sec.1.6 scopes bare sans-c to one named PROOF site (prp:nophifree Step 3) -- and
% the glossary was amended rather than the pointer patched (C88).  varsigma is a
% Sec.1.6 LOCAL and a local may not be made binding (C25); it had been made binding
% here, in a maintained assumption's statement.  eq:h1match is the matching written
% as a condition on one named pair of objects rather than as machinery, which is
% dispatch 21's change; through dispatch 20 the clause read "the Sobolev scale H^s
% of order s GENERATED BY D", and that phrase is struck.  The generator's space at
% order s' was written X_{s'} through dispatch 23 and is written out: italic X is
% the BINDING regressor X_1, X_2, subscripted by a period index, and a scale space
% subscripted by an order on the same letter is separated from it by nothing but
% what the subscript happens to be; swept d24, X_{s'} -> A 1 / R 0 / P 0.  s' was
% written sigma and tau through dispatch 21, against the four-way sigma rule and
% against the reserved tau_DP.  Glossary C87, C106.
$s>\tfrac{27}{16}$ at $\alpha=\tfrac38$, and it is \eqref{eq:lxwindow} and not \eqref{eq:hxwindow} that
binds there, the latter asking only $s>1.4818$; at every $\alpha<\approx0.359$ the two exchange roles
and \eqref{eq:hxwindow} is the higher floor.
% CORRECTED 2026-08-10 (Phase 6a), AND IT WAS FALSE BEFORE THE (P-g) REPAIR TOO.
% This read "and eq:hxwindow forces more still".  At alpha = 3/8, eq:lxwindow gives
% s > 27/16 = 1.6875 while eq:hxwindow gives 1/2 + 16(3/8)/(8 - 9/8) = 1.3727 on the
% old (P-g) and 1/2 + 18(3/8)/(8 - 9/8) = 1.4818 on the corrected one -- BELOW 27/16
% in both cases, so eq:hxwindow forced nothing more at the alpha named and the
% sentence was self-falsifying as printed.  Which of the two floors binds switches at
% the root of 9 alpha/(8(1 - 2 alpha)) = 1/2 + 18 alpha/(8 - 3 alpha), ~0.35914, so
% eq:hxwindow is the binding one BELOW that and eq:lxwindow above it, and alpha = 3/8
% sits above.  The paragraph's own point -- that s is forced above 3/2, so the
% literal generated reading of the scale would be false at every admissible s -- is
% carried by 27/16 alone and is unaffected.
\emph{A B-spline on a
quasi-uniform mesh satisfies no boundary condition}, so under the literal generated reading
\eqref{eq:sieveinhs} would be \textbf{false} at every admissible $s$ and the four conditions it underwrites
, A5(vi), A5(vi$'$), \cref{prp:phirate}(P-a.1) and \eqref{eq:foldrates}'s
second display, would again be events on the \emph{empty} set, which is the failure the amendment of
A8 closed and which the generated reading reopened verbatim.
\emph{What held it together was a directive and not a condition}: one sentence in a scope note in
\texttt{proofs.tex} instructing that this scale be read on the same scale as \eqref{eq:jackson}, i.e.\
classically, in a different file, and cited at neither this assumption, \cref{def:srcball}, nor
\eqref{eq:sieveinhs}. \emph{It is settled here instead, in the clause, once}: the space is the classical
one, and the comparison the generator was there to produce is \eqref{eq:h1match}, assumed outright.
\emph{Nothing downstream loses anything}: only the single index $s$ is ever used, no site reads a second
member of the scale, and \eqref{eq:h1match} is exactly what the generated construction was invoked to
deliver. \emph{Consequently the source index in A2 is}
\[
\mu_{\text{\textup{H1}}} \;:=\; \frac{s}{2a}.
\]
The subscripted name is used \textbf{only} where the two indices must be told apart, which is the
disposition \cref{rem:pasrc}\textup{(5)} makes binding. Bare $\mu$ in this project is (P-a.3)'s
\emph{primitive} index at the finite-rank sieve compression of the \emph{realised} operator at the anchor
$\theta_\circ$; $\mu_{\text{\textup{H1}}}$ is the index of A2 at the \emph{population}
operator $\mathcal A_{\mathrm{cond}}$ on all of $\mathcal V$ at $\theta_\dagger$, \emph{read into the
whole three-factor range space and not contrast by contrast, so that $a$ is the decay exponent of the
\textbf{stacked} operator at every $\mathsf L$}, and $s$ and $a$ in the
display are objects at $\theta_\dagger$ as well, against the $\theta_\circ$-specific $a$ of
\cref{def:crossfit}(b). \emph{They are equal only under the interpolation proviso this file states and
declines}, and the sites that computed one and read it into a condition on the other are qualified
in place: the two daggered lines of the smoothness-budget table, the sentence deducing
$\mu>\tfrac56$ from $s>a+\tfrac12$, \cref{rem:rootnscope}\textup{(1)}, and clause \textup{(2)} of
\cref{rem:h6bandwidth}, which \cref{rem:rootnscope}\textup{(1)} cites while qualifying its own copy
and which carried no proviso of its own until it was given one. \emph{No claim of exhaustiveness is
made, and that is deliberate}: the sentence this one replaces said ``the three sites'', then named
four things, and the fourth site was unqualified while the sentence asserted that all of them were.
\emph{\textbf{The same identity is \textbf{not} asserted for \cref{prp:phirate}.}}
\cref{prp:phirate}'s hypothesis (P-a.3) states a
% ARCHAEOLOGY.  The change of scope is the round-12 amendment to this assumption.
\textbf{spectral} source condition, at the sieve compression of the \emph{realised} operator, with $\mu$
and its radius $\mathsf R_\ast$ \emph{primitive}. What H1 supplies to that route is the \emph{scale}
$H^{s}$, which \cref{def:srcball}, A5(vi$'$), (P-a.1), (P-c) and (P-g) are written in, 
and nothing more; the passage from a Hilbert-scale source condition plus a link condition to a spectral
one is operator interpolation at exponent $2\mu$, licensed only at $2\mu\le1$, and
\cref{prp:phirate} lives at $2\mu>\tfrac53$. \cref{rem:pasrc} carries the full disposition, including
the commutativity proviso under which the identity above \emph{would} extend to the sample operator, and
what that proviso costs.

What $a$ is an exponent \textbf{of}, at $\mathsf L\ge3$. H1 is a hypothesis about the
\textbf{stacked} operator, $\mathcal A_{\mathrm{cond}}$ read into
$L^{2}(\Pi_u\otimes\Pi_t\otimes\Pi_{\mathsf w})$ entire, as \cref{def:spaces} fixes it, and not about
any single contrast. \emph{At $\mathsf L\ge3$ that is a genuinely different operator from the binary one},
and what makes it different is the \textbf{range enlargement}: the domain $\mathcal V$, the scale $H^{s}$
and \eqref{eq:h1match}'s form are untouched, and the range gains a third factor. \emph{It is not the
normalisation}, which was reversed, the operator carries the raw contrast \eqref{eq:contrast} on both
sides at every $\mathsf L$, as the paragraph \emph{A4} sets out.

And $a$ does not depend on $\Pi_{\mathsf w}$, which is worth saying because $\Pi_{\mathsf w}$ is the
analyst's choice. B1(iv.d) puts strictly positive masses on a finite index set, so any two admissible
$\Pi_{\mathsf w}$, $\Pi'_{\mathsf w}$ give \emph{equivalent} range norms,
\[
\min_{l}\frac{\varpi'_{\mathsf w}(l)}{\varpi_{\mathsf w}(l)}\;\|\Xi\|^{2}_{L^{2}(\Pi)}
\;\le\;\|\Xi\|^{2}_{L^{2}(\Pi')}\;\le\;
\max_{l}\frac{\varpi'_{\mathsf w}(l)}{\varpi_{\mathsf w}(l)}\;\|\Xi\|^{2}_{L^{2}(\Pi)}
\qquad\text{for every }\Xi ,
\]
with both ratios finite and strictly positive because the index set is finite and B1(iv.d)'s masses are;
and the \emph{domain} norm is untouched, $\Pi_{\mathsf w}$ appearing only on the range side. Singular
values of a compact operator are its approximation numbers, infima of operator-norm distances to
operators of rank below $j$, and an operator norm inherits those same two constants, so the singular
values under the two weightings differ by at most their square roots \emph{at every $j$ alike}. \emph{The
constants are $j$-free, so the exponent is identical and only the constant in the $\asymp$ moves}; and
$\max_l\varpi_{\mathsf w}(l)^{-1/2}$, the same factor that prices the passage between the
$L^{2}(\Pi)$ norm and the per-contrast norms, is what bounds the move. So the choice of range
weights is not a choice of ill-posedness.

One reading of $a$ in terms of the contrasts is available, and it is conditional. Write
$\mathcal A^{(l)}h:=\mathcal A_{\mathrm{cond}}h(\cdot,\cdot,l)$, so that
$\mathcal A^{(l)}:\mathcal V\to L^{2}(\Pi_u\otimes\Pi_t)$ is the $l$-th contrast operator, Hilbert--Schmidt
by \cref{prp:compact} read at fixed $l$, its kernel being $\mathsf k(\cdot,\cdot,l;\cdot)$ of
\eqref{eq:kernel}, and let $\mathcal A^{(l)\ast}\mathcal A^{(l)}\succeq0$ be its Gram operator on
$\mathcal V$, the adjoint taken in the \emph{unweighted} range pairing. Since
$\|\mathcal A_{\mathrm{cond}}h\|^{2}_{L^{2}(\Pi)}
=\sum_{l}\varpi_{\mathsf w}(l)\,\|\mathcal A^{(l)}h\|^{2}$ by \cref{def:spaces}'s pairing, and since two
bounded self-adjoint operators with the same quadratic form are equal,
\[
\mathcal A^{\ast}_{\mathrm{cond}}\mathcal A_{\mathrm{cond}}
\;=\;\sum_{l=1}^{\mathsf L-1}\varpi_{\mathsf w}(l)\;\mathcal A^{(l)\ast}\mathcal A^{(l)}
\qquad\text{on }\mathcal V .
\]
Two comparisons follow, both from the max--min characterisation of the eigenvalues of a compact positive
operator. \emph{Discarding all but one summand} leaves that sum $\succeq
\varpi_{\mathsf w}(l_0)\,\mathcal A^{(l_0)\ast}\mathcal A^{(l_0)}$ in the L\"owner order for each $l_0$,
and the max--min is monotone in that order, so
$\sigma_j(\mathcal A_{\mathrm{cond}})\ge\max_l\varpi_{\mathsf w}(l)^{1/2}\,\sigma_j(\mathcal A^{(l)})$ at
every $j$. \emph{In the other direction}, Weyl's inequality for compact positive operators, the
$(i+i'-1)$-th eigenvalue of a sum of two is at most the $i$-th of the first plus the $i'$-th of the
second, which is the same max--min evaluated on the intersection of two finite-codimensional subspaces
, iterated over the $\mathsf L-1$ summands gives
\[
\lambda_{\sum_lj_l-(\mathsf L-2)}\big(\mathcal A^{\ast}_{\mathrm{cond}}\mathcal A_{\mathrm{cond}}\big)
\;\le\;\sum_{l=1}^{\mathsf L-1}\varpi_{\mathsf w}(l)\,\sigma_{j_l}\big(\mathcal A^{(l)}\big)^{2}
\qquad\text{for any }j_1,\dots,j_{\mathsf L-1}\ge1 .
\]
Taking $j_l=\lceil j/(\mathsf L-1)\rceil$, for which $\sum_lj_l-(\mathsf L-2)\le j$ at every $j\ge1$, and
using that $\lambda$ is non-increasing in its index,
\[
\sigma_j(\mathcal A_{\mathrm{cond}})^{2}\;\le\;
\sum_{l=1}^{\mathsf L-1}\varpi_{\mathsf w}(l)\,\sigma_{\lceil j/(\mathsf L-1)\rceil}\big(\mathcal A^{(l)}\big)^{2} .
\]
\emph{So if the contrast operators have polynomial exponents, $\sigma_j(\mathcal A^{(l)})\asymp j^{-a_l}$,
the two bounds pinch and}
\[
\sigma_j(\mathcal A_{\mathrm{cond}})\;\asymp\;j^{-a},
\qquad
a\;=\;\min_{1\le l\le\mathsf L-1}a_l ,
\]
the constants depending on $\mathsf L$ and on $\varpi_{\mathsf w}$ and \emph{not} on $j$. \emph{Read the
lower bound at any $l$ attaining the minimum} and it is $\gtrsim j^{-\min_la_l}$; \emph{read the upper
bound through $j^{-2a_l}\le j^{-2\min_la_l}$ for $j\ge1$}, and, since
$\lceil j/(\mathsf L-1)\rceil\asymp_{\mathsf L}j$, it is $\lesssim j^{-\min_la_l}$ as well.
\emph{At $\mathsf L=2$ every step is an identity and
not a special case}: there is one contrast, $\Pi_{\mathsf w}=\delta_1$,
$\mathcal A^{(1)}=\mathcal A_{\mathrm{cond}}$, both comparisons hold with constant $1$ and with $j_1=j$
and offset $\mathsf L-2=0$, and $a=a_1$ is the binary $a$ of every earlier round.

Two things do \textbf{not} follow, and each is a way the display above invites being over-read.
First, the implication runs one way. Per-contrast polynomial decay is \textbf{sufficient} for
H1's decay clause and is not necessary, a design in which one contrast operator has finite rank and
another decays polynomially satisfies that clause and has no $a_l$ at the first contrast, so H1 stays
a hypothesis about the stacked operator, assumed and not derived. And it reaches only the decay
clause. \eqref{eq:h1match} compares $H^{s}$ with the powers of the \emph{summed} Gram operator
$\sum_l\varpi_{\mathsf w}(l)\,\mathcal A^{(l)\ast}\mathcal A^{(l)}$;
per-contrast matchings do not aggregate to it, and nothing above supplies it. Second,
$\min_la_l$ is not a comparison with ``the binary $a$'' unless the binarisation is named.
$\mathcal A^{(l)}$ is contrast $l$ against the \emph{base} cell in the $\mathsf L$-category
parameterisation. Merging cells, $W\ge1$ against $W=0$, say, produces the contrast
$\sum_{l\ge1}\widetilde W_l$, hence by linearity of \eqref{eq:Mdef} in the contrast the moment
$\sum_{l\ge1}\mathcal M(\cdot,\cdot,l)$ and the linearisation $\sum_{l\ge1}\mathcal A^{(l)}$, which is a
\textbf{sum} of contrast operators and is none of them; a sum admits cancellation that no hypothesis on
its summands excludes, so no exponent transfers to it in either direction. \emph{And merging is not a
free re-parameterisation in any case}, for the reason the \emph{Discussion} paragraph gives at its
coarsening warning: it adopts the \emph{stronger} identification hypothesis, since
$\mathcal V_{CI}(W)\subseteq\mathcal V_{CI}(g(W))$ for every coarsening $g$, and it forfeits the
localisation. A reader who takes $a=\min_la_l$ to
mean ``the exponent improves on the binary design'' is reading more than this gives. \emph{What it does
license is the comparison between the full stack and a sub-collection of its own contrasts, inside one
fixed $\mathsf L$-category design}: $\varpi_{\mathsf w}(l)>0$ makes
$\ker\mathcal A_{\mathrm{cond}}=\bigcap_l\ker\mathcal A^{(l)}$, so the full stack has the weakly smaller
kernel and the weakly larger rank, and H1's \emph{infinite-rank} clause is to that extent weakly easier
at the stack than at any one contrast. \emph{The \textbf{level} of $\sigma_j$ carries no such
monotonicity}: $\Pi_{\mathsf w}$ is a probability measure, so
$\|\mathcal A_{\mathrm{cond}}h\|^{2}$ is an \emph{average} over contrasts and not a sum, and adding a
contrast that carries little signal can lower it.

% ==========================================================================
%  THIS BLOCK CLOSES THE DEFERRAL THE DISPATCH-25 CENSUS RECORDED AT
%  assumptions_main.tex, "RE-READING THE BUDGET TABLE AT sans-L >= 3 IS NOT DONE
%  HERE AND IS FLAGGED FOR PHASE 2", and plan rev-2 Sec.4 item 2.  What was
%  deferred was a QUESTION -- do the per-contrast operators share a decay
%  exponent, and if not what is `a' -- and the answer is the pinch above.
%
%  WHAT WAS CHECKED, AND ONE THING IN THE HAND-OFF WAS WRONG.
%   * The Gram decomposition A*_cond A_cond = sum_l varpi_w(l) A^{(l)*} A^{(l)}
%     is exact: the quadratic forms agree by def:spaces's pairing, and two bounded
%     self-adjoint operators with equal quadratic forms are equal (polarisation).
%     NO SYMBOL IS INTRODUCED FOR THE PER-CONTRAST GRAM OPERATOR.  B_l is what the
%     hand-off called it and is what a draft of this block used; calligraphic B is
%     ALREADY the second-order form operator of suppeq:Bpi (declared at
%     def:spaces's "one declared exception"), and calligraphic C, E, G, H, U and R
%     are claimed too (results_supp.tex:4127).  Writing A^{(l)*} A^{(l)} out is
%     five characters longer and collides with nothing.
%   * THE INDEX IN THE WEYL STEP.  The hand-off proposed
%     j_l = ceil((j + sans-L - 2)/(sans-L - 1)).  THAT IS THE WRONG DIRECTION:
%     it makes m := sum_l j_l - (sans-L - 2) >= j, and Weyl then bounds
%     lambda_m, not lambda_j.  What is needed is m <= j.  The correct choice is
%     j_l = ceil(j/(sans-L - 1)) = floor((j + sans-L - 2)/(sans-L - 1)) -- the
%     two are the same integer, and it is the CEILING OF j/(sans-L - 1), not of
%     the shifted quotient.  Check: (sans-L - 1) ceil(j/(sans-L - 1))
%     <= j + sans-L - 2 for every j >= 1 (write j = q(sans-L - 1) + r; at r = 0
%     the left side is j, at r >= 1 it is j - r + sans-L - 1).  The CONCLUSION is
%     unaffected -- the shifted version still gives the same exponent after
%     re-indexing j -> j - (sans-L - 1) -- but the inequality as handed off does
%     not bound sigma_j, so it is printed in the corrected form.
%   * sans-L = 2 REDUCES WITH NO EXCEPTION: one summand, Weyl is the identity,
%     the offset sans-L - 2 is 0, varpi_w(1) = 1, and both bounds are equalities.
%
%  WHAT IS *NOT* CLAIMED, AND THE TWO LIMITS ARE PRINTED ABOVE RATHER THAN
%  RECORDED HERE BECAUSE A READER WILL REACH FOR EXACTLY THESE TWO.
%   (1) The stacked decay does NOT imply per-contrast decay.  Only the forward
%       implication holds.  And the MATCHING half of H1 is not reached at all.
%   (2) min_l a_l is NOT "better than the binary a".  The merged operator is
%       sum_l A^{(l)}, not any A^{(l)}; sums cancel.  The only monotone
%       comparison available is the KERNEL one, inside a fixed sans-L-category
%       design, and it is about RANK and not about the exponent.
%
%  AND NOTE WHAT DOES NOT MOVE.  `a' does not enter the OPERATIVE lines of the
%  smoothness-budget table at all -- only the two DAGGERED interpretation lines
%  carry it, mu having been primitive at (P-a.3) since round 12 -- so the pinch
%  changes no condition this block imposes.  The table's own paragraph says so.
% ==========================================================================

% ------------------------------------------------------------------------
\paragraph{A3 (RETIRED: local rank), moved entire.}

\begin{assumption}[A3: RETIRED, local rank, now \cref{lem:momzero}\textup{(iv)}]\label{ass:rank}
\textbf{\sout{RETIRED.}} This assumption is not maintained, it is not available for use, and no
result in this programme takes it as a hypothesis. It is retired because it is a \textbf{theorem}, not
because it is false, as A6 was, and not because it is too strong. It read: \emph{suppose
A2 holds, so that $\ell_\star$ exists; then}
% ARCHAEOLOGY.  It was maintained through dispatch 12.
\[
J \;:=\; E[\Delta Z\,\Delta Z'] \;+\; E\big[\widetilde W\,\ell_\star(\Delta U,\Delta Z)\,\Delta Z'\big]
\quad\text{\emph{is nonsingular.}}
\qquad\text{\sout{[retired]}}
\]
\emph{Why it is retired.} \Cref{lem:momzero}\textup{(iv)} proves that under A1's \emph{first} clause the
second summand is the zero matrix, $\ell_\star(\Delta U,\Delta Z)\Delta Z'$ is
$\sigma(\Delta U,\Delta Z)$-measurable and $E[\widetilde W\mid\Delta U,\Delta Z]=0$, so
$J=E[\Delta Z\Delta Z']$ whenever $J$ is defined at all, and nonsingularity of that matrix is A1b(iv).
% DISPATCH-25 NOTE, AND THE DISPLAY ABOVE IS DELIBERATELY LEFT AS IT WAS PRINTED.
% The retired statement is an ARCHAEOLOGICAL RECORD of what A3 read; rewriting it
% would misreport the history the retirement exists to preserve.  At sans-L >= 3
% the same object is E[ell_star(Delta U, Delta Z) Wtilde Delta Z'], with the
% p x (sans-L - 1) matrix ell_star acting on the (sans-L - 1)-vector Wtilde to
% give a vector in R^p and then an outer product with Delta Z' -- which is what
% the notation block's type declaration produces, and which is the display above
% at sans-L = 2 because Wtilde is then a scalar and commutes.  THE REFUTATION IS
% UNCHANGED AND IS CARDINALITY-FREE: conditioning on (Delta U, Delta Z) leaves
% ell_star E[Wtilde | Delta U, Delta Z] Delta Z', and E[Wtilde | Delta U, Delta Z]
% is the ZERO VECTOR under A1 for exactly the reason it was the zero scalar --
% each Wtilde_l is conditionally centred by construction of eq:contrast.
% ROUND 4: psi CARRIES THE RAW CONTRAST, SO THIS RECORD IS READ AT THE OBJECT IT
% WAS ALWAYS ABOUT AND NEEDS NO QUALIFICATION.  (Round 3 added one, psi having
% been normalised then; it noted that the vanishing survives at Wtilde^Omega too,
% since Omega is sigma(Delta Z)-measurable and
% E[Wtilde^Omega | Delta U, Delta Z] = Omega^{-1} E[Wtilde | Delta U, Delta Z] = 0.
% That remains TRUE and is now IDLE, the normalisation having been reversed.)
The
proof consumes A1, A1b(iv), A$2'$(i), A$2'$(iii) and B1(i), (ii), (iv.a), (iv.d), \emph{A1 for the vanishing and
A1b(iv) for the nonsingularity; A$2'$(i), A$2'$(iii) and the three B1 clauses only to make the second
summand a finite matrix in the first place, through \cref{cor:elldiff}}. \emph{Every one of those tags was
already carried by every list
that carried A3}, so striking the tag is a strict weakening of four hypothesis lists with no change to any
conclusion: \cref{lem:sandwich}\textup{(i)} and \textup{(i$'$)}, \cref{thm:rootn}'s group (T-2), and
\cref{cor:ratefloor} through (T-2).

What survives, and it is the half the conditional statement was invented to carry. A3 was stated
conditional on A2 because $\ell_\star$ exists only under A2, so the display above is a well-formed object
only under A2 and ``A3 holds while A2 fails'' had no truth value. \emph{That remains exactly true of the
display} and \cref{rem:rootnscope}\textup{(4)} still says so, in the sharper form that what fails under
$\neg$A2 is the \emph{existence} of $\ell_\star$, $L_\star$ and $\psi_\star$ and not a rank condition. It
also remains true that A3 was never what identifies $\theta_\dagger$: that is A1b(iv). What is
withdrawn is the further claim, made at A1b(iv)'s comment and at three sites in
\texttt{results.tex}, that the two are \textbf{different} matrices. They are the same matrix wherever
both are defined.
% ARCHAEOLOGY.  The round-2 correction (M-3) was right about the conditionality and
% wrong about the identity, and the identity was never computed.

\emph{What is \textbf{not} conceded.} No conclusion weakens, no rate moves, and no other tag changes
status. In particular A$2'$(i.b) is \textbf{not} freed, see \cref{lem:momzero}\textup{(iv.1)}, and
A2 is not freed, being consumed at \cref{lem:biasexp}, \eqref{eq:biasb1}, for the first-order bias, which
is where the source condition does its work. \emph{And the counterexample probe is redirected as it always
was}: the object to probe is A2.
\end{assumption}

% ------------------------------------------------------------------------
\paragraph{A4 (overlap and propensity regularity).}

% ADDED IN DISPATCH 25.  A4(i) NOW CONTAINS THE WORD "CONSEQUENTLY", AND THE A8
% PRECEDENT IN THIS FILE IS EXPLICIT THAT A DERIVATION BEHIND THAT WORD IS
% PRINTED RATHER THAN ASSERTED.  Three things are derived here: the floor
% eq:omegafloor, the prefactor identity eq:omegaprefac, and the inverse
% eq:omegainv on which the normalisation decision rests.  NONE OF THE THREE IS
% AN ASSUMPTION and none adds a hypothesis to any list.
\emph{The word ``consequently'' in A4(i) is a derivation, and it is printed here rather than asserted}:
\eqref{eq:omegafloor} is what every statement that divides by $\Omega$ points at. Write
$\varrho=\varrho(z)$ and $\Omega=\Omega(z)$ throughout, and recall $\varrho_0=1-\mathbf1'\varrho$.

The inverse in closed form, and it is where everything else comes from. For any
$\varrho$ with $\varrho_l>0$ at every $l\in\{0,\dots,\mathsf L-1\}$,
\begin{equation}\label{eq:omegainv}
\Omega^{-1}\;=\;\operatorname{diag}(\varrho)^{-1}+\varrho_0^{-1}\mathbf1\mathbf1' ,
\qquad\text{and}\qquad
\big(\Omega^{-1}\widetilde W\big)_l\;=\;\frac{\un\{W=l\}}{\varrho_l}-\frac{\un\{W=0\}}{\varrho_0} .
\end{equation}
\emph{A notational convention, stated here because this is the identity map's first site in compiled
order: the identity map, on any space and in any dimension, is written $\mathrm{Id}$ throughout this
project and never $I$, italic $I$ being C3's boundary index \eqref{eq:rescaled} --- a closed interval,
not an operator.} For the first identity, multiply out:
$\big(\operatorname{diag}(\varrho)-\varrho\varrho'\big)\big(\operatorname{diag}(\varrho)^{-1}+\varrho_0^{-1}\mathbf1\mathbf1'\big)
=\mathrm{Id}+\varrho_0^{-1}\varrho\mathbf1'-\varrho\mathbf1'-\varrho_0^{-1}(1-\varrho_0)\varrho\mathbf1'
=\mathrm{Id}+\varrho\mathbf1'\big\{\varrho_0^{-1}-1-\varrho_0^{-1}+1\big\}=\mathrm{Id}$,
using $\operatorname{diag}(\varrho)\mathbf1=\varrho$, $\varrho'\operatorname{diag}(\varrho)^{-1}=\mathbf1'$
and $\varrho'\mathbf1=1-\varrho_0$; every one of the three needs $\varrho_l>0$ at some point, and
$\varrho_0>0$ is what makes the second summand finite. For the
second identity, $\widetilde W_l/\varrho_l=\un\{W=l\}/\varrho_l-1$ and
$\mathbf1'\widetilde W=(1-\un\{W=0\})-(1-\varrho_0)=\varrho_0-\un\{W=0\}$, so the two added terms are
$\un\{W=l\}/\varrho_l-1$ and $1-\un\{W=0\}/\varrho_0$ and the $\pm1$ cancel.

\emph{The floor.} \emph{Matrix norms in this paragraph are spectral}, so that
$\lambda_{\min}(\Omega)=\|\Omega^{-1}\|^{-1}$ for the symmetric positive definite $\Omega$. From
\eqref{eq:omegainv} and the triangle inequality,
$\|\Omega^{-1}\|\le\max_{l\ge1}\varrho_l^{-1}+\varrho_0^{-1}\|\mathbf1\mathbf1'\|
=\max_{l\ge1}\varrho_l^{-1}+(\mathsf L-1)\varrho_0^{-1}$, since the only nonzero eigenvalue of
$\mathbf1\mathbf1'$ on $\IR^{\mathsf L-1}$ is $\mathbf1'\mathbf1=\mathsf L-1$. Under \textbf{A4(i)'s
condition} --- the sentence beginning ``The condition is''; this pointer read ``A4(i)'s \emph{first
sentence}'' through 2026-08-25, and A4(i)'s first sentence is now its \emph{declaration}, which fixes
$\varepsilon_\varrho$ and asserts no inequality --- each of the
$\mathsf L$ cell probabilities is at least $\varepsilon_\varrho$, so
$\|\Omega^{-1}\|\le\mathsf L/\varepsilon_\varrho$ and
$\lambda_{\min}(\Omega)=\|\Omega^{-1}\|^{-1}\ge\varepsilon_\varrho/\mathsf L$, which is
\eqref{eq:omegafloor}. \emph{The bound is valid and not sharp}: at $\mathsf L=2$ it asserts
$\varrho(1-\varrho)\ge\varepsilon_\varrho/2$ where the exact minimum over
$[\varepsilon_\varrho,1-\varepsilon_\varrho]$ is $\varepsilon_\varrho(1-\varepsilon_\varrho)$. Nothing in
this corpus is stated \emph{at} the value of the floor, only under it, so the slack costs nothing;
\emph{a dispatch that begins to track the constant should re-derive it rather than sharpen this one in
place}. \emph{The base cell is in A4(i)'s quantifier for this reason and not for symmetry}:
$\varrho_0\ge\varepsilon_\varrho$ is what bounds the second summand, and if $\varrho_0\to0$ then
$\|\Omega^{-1}\|\to\infty$ however far the non-base cells stay from zero. \emph{A version of A4(i)
quantified only over $l\in\{1,\dots,\mathsf L-1\}$ would therefore not deliver
\eqref{eq:omegafloor}}, and at $\mathsf L=2$ it would be the one-sided condition
$\varrho\ge\varepsilon_\varrho$, which is strictly weaker than the two-sided interval every earlier
round assumed.

The prefactor, which is what A4(i) cancels. Write $F_m(u\mid z):=P(\Delta U\le u\mid W=m,\Delta Z=z)$
and $d(u\mid z):=\big(F_1-F_0,\dots,F_{\mathsf L-1}-F_0\big)'(u\mid z)$. Then, at $\sigma(\Delta Z)$ and
only there,
\begin{equation}\label{eq:omegaprefac}
E\big[\widetilde W\,\un\{\Delta U\le u\}\,\big|\,\Delta Z=z\big]\;=\;\Omega(z)\,d(u\mid z) .
\end{equation}
Component $l$ of the left side is
$\varrho_lF_l-\varrho_l\sum_{m}\varrho_mF_m=\varrho_l\big(F_l-\sum_m\varrho_mF_m\big)$; writing
$F_l-F_m=d_l-d_m$ with $d_0:=0$ and using $\sum_m\varrho_m=1$ gives
$F_l-\sum_m\varrho_mF_m=d_l-\varrho'd$, so the vector is
$\operatorname{diag}(\varrho)d-\varrho\varrho'd=\Omega d$. At $\mathsf L=2$ this is
$\varrho(z)(1-\varrho(z))\{F_1(u\mid z)-F_0(u\mid z)\}$, the display the block has carried since
round~1. Under \eqref{eq:omegafloor} the map $\Omega(z)$ is invertible and the restriction can be
\emph{divided through}. \emph{That division is what makes the reading of \suppref{Remark}{htread}
available; it is \textbf{not} performed by the estimator}, whose operator and moment both carry the raw
contrast.
% ROUND-5 CORRECTION.  This sentence ended "...and hence what makes the
% NORMALISED OPERATOR OF THE NEXT PARAGRAPH well defined."  THERE IS NO
% NORMALISED OPERATOR IN THE NEXT PARAGRAPH -- round 4 replaced it by "a claim
% withdrawn here" -- so the sentence pointed, IN THE PDF, at an object the
% reversal had removed, and did so at exactly the place a reader checks the
% reversal.  An orphaned forward pointer in printed text is the defect species
% this block spends its length policing.

What that division is \textbf{not} needed for, and a claim withdrawn here. It is not what
makes \cref{lem:bierens} an equivalence. A first version of this paragraph said it was, and located
the failure on $\{z:\det\Omega(z)=0\}$; \emph{both halves are wrong at $\mathsf L\ge3$}. Take
$\mathsf L=3$ and $\varrho(z)=(q,0)'$ with $q\in(0,1)$: cell~$2$ is null, $\det\Omega(z)=0$, and yet $W$
is not conditionally degenerate, $\widetilde W\ne0$, and the restriction still forces $F_1=F_0$. The
equivalence $\Omega d_\varsigma=0\Leftrightarrow$ conditional independence holds at \emph{every} $z$ and
under \emph{no} overlap condition, by the two-case argument given in the paragraph \emph{A1} above, 
$\Omega_A$ nonsingular when $\varrho_0>0$, $\ker\Omega_A=\operatorname{span}\{\mathbf1_A\}$ when
$\varrho_0=0$, and in both cases the non-null cells are forced to share one conditional law.
\emph{The only set on which the restriction is vacuous is where $W$ is conditionally degenerate}, that
is, where $\Omega(z)=0$, equivalently where $\max_{0\le l\le\mathsf L-1}\varrho_l(z)=1$. \emph{This is
the $\mathsf L$-category reading of $\mathcal Z_0$ and it is \textbf{not} a new definition}: $\mathcal Z_0$
is defined at \eqref{eq:degenreg} in exactly this form, which recovers the binary
$\{z:\varrho(z)(1-\varrho(z))=0\}$ verbatim at $\mathsf L=2$.
% AMENDED 2026-08-09 BY THE ORCHESTRATOR, closing a blocker Phase 5a raised against a file it did
% not own.  This sentence read: "Z_0 is defined at eq:degenreg IN ITS BINARY FORM
% {z:varrho(z)(1-varrho(z))=0}, which {z:Omega(z)=0} recovers verbatim at L=2; AMENDING THAT
% DEFINITION IS THE BUSINESS OF THE FILE THAT CARRIES IT, and this paragraph fixes the reading
% rather than the symbol."  5a amended eq:degenreg -- it was ill-typed at L>=3, varrho being a
% vector there -- so BOTH halves became false: the definition is no longer binary, and the amending
% has happened.  That is this project's dominant defect species (a claim about the artifact that the
% artifact contradicts) and it was in PRINTED text.  eq:degenreg -- cited by LABEL, because the line
% number this note carried ("proofs.tex:2173") was stale by the time it was read; the display sits in
% lem:quantu's proof, after the eq:binaryquant paragraph, and is at proofs.tex:2244 as of 2026-08-10 --
% now reads
% Z_0 := {z: Omega(z)=0} = {z: max_{0<=l<=L-1} varrho_l(z)=1}, agreeing character-for-character with
% :1063 above and with results_s2_identification.tex:281.  [POINTER REPAIRED 2026-08-10, Phase C.  A
% stale line number INSIDE A DISCHARGE NOTE is the worst place for one: the next reader follows it to
% check the discharge.  Anchor on the label and on the surrounding text, not on the number.]
% The round-3 correction note below still
% describes a real hazard and is kept: there must remain exactly ONE definition of Z_0, at
% eq:degenreg, and this paragraph must go on citing it rather than restating it.
% ROUND-3 CORRECTION.  Round 2 wrote the sans-L-category set out as an UNLABELLED
% DISPLAY here, which put a second definition of Z_0 in the corpus while
% eq:degenreg carries the binary one in proofs.tex.  Two definitions of one
% symbol, in two files, is the collision this block spends most of its length
% policing.  The display is withdrawn; the set is described in words and
% eq:degenreg is cited.  ~~Amending eq:degenreg itself is Phase 3's, and it is
% listed with the other cross-file obligations.~~  STRUCK 2026-08-10 (Phase C):
% PHASE 3 CLOSED WITHOUT IT -- the fourth instance of the stranded-item failure
% mode, since assumptions_main.tex:4384 assigned it to the same phase -- and the
% amendment was performed in the 5a round instead, as the note above records.
% The obligation is discharged; only the routing was wrong.
\emph{As a condition on the propensity vector, degeneracy is strictly stronger than singularity of
$\Omega$ once $\mathsf L\ge3$}, $\varrho=(q,0)'$ is singular and not degenerate, \emph{though for a
particular data-generating process the two $z$-sets may of course coincide}, and nothing here asserts
that they do not. \emph{\Cref{lem:bierens} states its equivalence under B1(i)--(iv) alone and its
clauses (a) and (b) are
correct as stated at every $\mathsf L$}; what does not generalise is the binary \emph{ground} recorded
beside it, that the scalar prefactor vanishes on the degenerate region. A dispatch that
generalises the printed reason instead of the printed result will import an overlap hypothesis
(a)$\Leftrightarrow$(b) neither has nor needs.

The statement is unqualified, and a qualification an earlier draft attached to it is withdrawn.
That draft said \cref{lem:bierens}'s clauses naming $\mathcal M$ and $T$ acquire A4(i), because those
objects would not exist at a normalised contrast. The estimator carries the raw contrast, so
$\mathcal M$ and $T$ are defined everywhere and bounded, and no clause of \cref{lem:bierens} needs an
overlap condition. \emph{The draft's reason was also wrong on its own terms}: $\Omega^{-1}\widetilde W$
is not ``non-integrable'' where a cell vanishes, it is \emph{undefined} there, $\Omega$ having no
inverse, and where it is defined it is integrable with a $\varrho$-free bound, as
\suppref{Remark}{htread}\textup{(4)} computes. \emph{Undefinedness and non-integrability are different
failures} and the corpus must not conflate them.

The determinant, since it was used above and its usual derivation presumes what fails there. For
$\varrho$ in the open region where every $\varrho_l>0$, the matrix determinant lemma applied to
\eqref{eq:omegainv} gives
$\det\Omega=\det\operatorname{diag}(\varrho)\,(1-\varrho'\operatorname{diag}(\varrho)^{-1}\varrho)
=\big(\prod_{l\ge1}\varrho_l\big)\varrho_0$. \emph{Both sides are polynomials in
$(\varrho_1,\dots,\varrho_{\mathsf L-1})$}, they agree on a non-empty open subset of
$\IR^{\mathsf L-1}$, and two polynomials agreeing on a non-empty open set agree identically; hence
$\det\Omega=\prod_{l=0}^{\mathsf L-1}\varrho_l$ \emph{everywhere}, including at the boundary where
$\Omega$ is singular and the lemma does not apply. \emph{That is exactly where the identity is wanted},
which is why the continuation step is printed rather than the lemma cited alone.

\medskip
\noindent
 If the prefactor is inverted at all, which the estimator does not do, but which the reading of
\suppref{Remark}{htread} does, the inverse is $\Omega^{-1}$, and that is settled by an identity rather
than by a preference.
The scalar $\varrho(1-\varrho)$ admits exactly one componentwise generalisation, the diagonal
matrix with entries $\varrho_l(1-\varrho_l)$, and it is wrong. \eqref{eq:omegaprefac} says the map
carrying the cell-versus-base contrasts $d$ into the observable moment is $\Omega$, whose off-diagonal
entries $-\varrho_l\varrho_m$ are the negative cross-cell correlation that a multinomial has and a
collection of separate binaries does not. Dividing component $l$ by $\varrho_l(1-\varrho_l)$ therefore
does not return $d_l$, and what it returns is not a contrast between cell $l$ and the base cell at all.
Two further facts make $\Omega^{-1}$ the right inverse to take. By \eqref{eq:omegainv},
$\Omega^{-1}\widetilde W$ is the vector of \emph{inverse-probability} contrasts of each cell against the
base, so $\Omega^{-1}$ is the only inverse under which the contrast vector is the
Horvitz--Thompson object the binary design already was, and at $\mathsf L=2$ it reads
$(W-\varrho)/\{\varrho(1-\varrho)\}=\un\{W=1\}/\varrho-\un\{W=0\}/(1-\varrho)$, the identity a binary
design writes down without comment. And the quadratic form is base-free:
$\widetilde W'\Omega^{-1}\widetilde W=(1-\varrho_W)/\varrho_W$, where $\varrho_W$ is the conditional
probability of the \emph{realised} cell, a quantity that names no base cell, is symmetric in all
$\mathsf L$ cells, and satisfies $E[\widetilde W'\Omega^{-1}\widetilde W\mid\Delta Z]=\mathsf L-1$. The
componentwise weighting has neither property. So \textbf{which} inverse to take, if one is taken at
all, is settled by an identity and not by a preference.

\medskip
\noindent
 \textbf{Whether} to take it is a second question, and the answer is \textbf{no}: the estimator
carries the raw contrast \eqref{eq:contrast} on both sides. The operator equation and $\psi$ of
\eqref{eq:psidef} are both written at $\widetilde W$, and $\Omega^{-1}$ appears nowhere in either.
The reading it supports is retained, as \suppref{Remark}{htread}, and is interpretation rather than
specification.

What normalising would have changed, so that the trade is legible. By the tower property with
\eqref{eq:omegaprefac} the two candidate forward residuals are
\[
T(u,t,\cdot)\;=\;E\big[\Omega(\Delta Z)\,d(u\mid\Delta Z)\,w(t'\Phi(\Delta Z))\big],
\qquad
T^{\Omega}(u,t,\cdot)\;=\;E\big[d(u\mid\Delta Z)\,w(t'\Phi(\Delta Z))\big] ,
\]
which differ in where $\Omega$ sits relative to the \emph{outer} integration over $\Delta Z$, a
difference no choice of norm on the range space can undo. The gain would be that the conditional
forward residual is the cell difference $d$ itself rather than a matrix mixture of it.

The price is that $T^{\Omega}$ has a different linearisation, and the whole smoothness budget is
stated at the one the paper has. A2's range condition, A$2'$, A1c and H1, the last through
$\sigma_j(\mathcal A_{\mathrm{cond}})\asymp j^{-a}$, the exponent of the \emph{stacked} operator on the
whole three-factor range space at every $\mathsf L$, are all stated at $\mathcal A_{\mathrm{cond}}$,
and $\mu_{\text{\textup{H1}}}=s/(2a)$, the two daggered lines of the smoothness-budget table below and
the comparison to the cross-sectional companion sit downstream of $a$. \emph{\eqref{eq:frontier} does
not}, its $\mu$ being (P-a.3)'s primitive index, and no other line of that table imposes a condition
involving $a$; \emph{that
narrows what a re-derivation would have to cover and does not remove the objection}, since the scale
$H^{s}$ that \cref{def:srcball}, A5(vi$'$), (P-c) and (P-g) are written in, and A2's Picard form, would
both have to be re-established at the normalised operator. \emph{A4's note (c) already records that the two
operators must not be assumed to share range or kernel}, and the paragraph on Khan--Tamer irregularity
scopes range-invariance to a \emph{constant} $\varrho$, a $z$-varying $\Omega$ being exactly what that
scoping excludes. \emph{An interpretational gain does not buy a re-derivation of the
ill-posedness budget}, and the author declines the trade. \emph{Nothing else in the dispatch depends on
it}: the $\mathsf L$-category instrument, the contrast vector, $\Omega$ as covariance and prefactor,
A4(i), \eqref{eq:trunset} and \eqref{eq:omegafloor} all stand.

One consequence of the reversal is worth stating positively. At $\mathsf L=2$ the operator, $\psi$
and every assumption clause reduce to the printed binary objects \emph{verbatim, with no exceptions}, and
the binary estimator is unchanged. Under the normalisation they did not: \eqref{eq:hatT} would have
carried a factor the pre-dispatch display did not have.

And one argument made along the way is false and is recorded because it is the first repair a
re-inventor will reach for. It ran: $\Omega$ is $\sigma(\Delta Z)$-measurable,
so $L\widetilde W=(L\Omega)\widetilde W^{\Omega}$, and $L\mapsto L\Omega$ is a bijection of
\cref{def:weight}'s admissible class onto itself, whence a normalised operator could be paired with a
$\psi$ printed at $\widetilde W$
without loss. \emph{Four bounds were checked and each of the four is correct}, \emph{the norms being spectral on
$\Omega$ and Frobenius on $L$, $\ell$ and $b$, as the notation block fixes, which is why the two are
named here rather than left to context}, 
$L^{\,b}_{\ell}\Omega=L^{\,b\Omega}_{\ell\Omega}$ since $\Omega$ does not depend on $u$;
$\|\Omega\|\le1$, because for a unit $c$, $c'\Omega c=\operatorname{Var}(c'\widetilde W\mid\Delta Z)\le
E[(c'\widetilde W+c'\varrho)^{2}\mid\Delta Z]=\sum_lc_l^{2}\varrho_l\le1$, so
$\|\ell\Omega\|_\infty\le\|\ell\|_\infty$ and $E\|b\Omega\|^{2}\le E\|b\|^{2}$, \emph{a Frobenius
norm is submultiplicative against a spectral one, $\|AB\|_F\le\|A\|_F\|B\|$}; the inverse map with
$\|\Omega^{-1}\|\le\mathsf L/\varepsilon_\varrho$; and $E[L\Omega\mid\Delta Z]=E[L\mid\Delta Z]\Omega$,
so N1 transfers. They are jointly insufficient, because \cref{def:weight} does not define its
admissible class by bounds. Its clause (a) binds $\ell:=\mathcal W\ell^\Pi$, so membership requires the
generator to lie in $\mathcal R(\mathcal W)$; the $z$-sections of $\mathcal W$ are Laplace transforms
against $w=\exp$ on a compact index set under B1(i)--(ii) and (iv.c), hence \emph{real-analytic} in $z$,
while A4(ii) makes $\varrho$, and therefore $\Omega$, only H\"older of order $r$.
\emph{Right-multiplication by a non-analytic $\Omega(z)$ leaves $\mathcal R(\mathcal W)$}, and the
failure occurs already at $\mathsf L=2$ with $p=1$ and a $\varrho$ that is H\"older but non-analytic at a
point. So the two parameterisations do not have the same admissible weight class, and no bijection
is claimed anywhere in this corpus. \emph{The four bound-checks are left on record rather than deleted},
because the lesson is that verifying the bounds of a class is not verifying membership of it when the
class is defined by a range condition, \emph{and because the same analytic-versus-H\"older mechanism
is what makes the author's ground for keeping the raw contrast bite}: a normalised
$\mathcal A_{\mathrm{cond}}$ cannot be assumed to inherit the range and decay statements A2, A$2'$, A1c
and H1 are made at.

\emph{Nothing in the assumption block changes on account of any of this}, because the estimator was not
re-written. The bounds the block does use are the free ones, $\|\widetilde W\|_\infty\le1$ and
$\|\widetilde W\|_2\le\min\{\sqrt2,\sqrt{\mathsf L-1}\}$, uniform in $\mathsf L$ and free of $\varrho$;
A4(i) remains idle in \cref{prp:wellposed}, \cref{prp:frechet} and \cref{prp:compact};
$\hat T$ remains affine in $\hat\varrho$, so the propensity branch of the operator error keeps its
exactly-zero second-order term; and no clause of \cref{lem:bierens} needs an overlap condition.

% ARITHMETIC CHECK OF THE QUADRATIC FORM, kept because it is the one line a
% reader will want to verify and it is two cases.  At W = k >= 1:
% Wtilde_l = 1{l = k} - varrho_l and 1{W = 0} = 0, so the sum over l >= 1 of
% Wtilde_l (1{W = l}/varrho_l) is (1 - varrho_k)/varrho_k.  At W = 0:
% Wtilde_l = -varrho_l and 1{W = 0}/varrho_0 = 1/varrho_0, so the sum is
% sum_{l>=1} varrho_l / varrho_0 = (1 - varrho_0)/varrho_0.  Both are
% (1 - varrho_W)/varrho_W.  The conditional mean is sum_k varrho_k
% (1 - varrho_k)/varrho_k = sans-L - 1 = tr(Omega^{-1} Omega).  At sans-L = 2 the
% mean is 1, which is the binary case's E[(W - varrho)^2]/{varrho(1-varrho)}.

\medskip
\noindent
 The lower bound $d_Z\ge1$ is load-bearing. At $d_Z=0$ the exponent $\alpha=r/(2r+d_Z)$ is
$\tfrac12$ \emph{exactly}, and:
% ARCHAEOLOGY.  The lower bound is stated in dispatch 14; it was assumed silently
% before it, at four sites and in one place substantively.
\begin{itemize}
\item One site genuinely breaks, and it is the $\varphi$-step. (P-d)'s sieve-entropy condition is
$L_x\ln n=o(n\hat\delta_n^{2})$. At $\alpha=\tfrac12$, A5(i)'s $\hat\delta_n=\max\{n^{-1/2}\sqrt{\ln n},r_n\}$
is $n^{-1/2}\sqrt{\ln n}$, \emph{an equality, and it is the declared floor of clause \textup{(i)} that
supplies it, not a two-sided reading of an $O_P$}, so $n\hat\delta_n^{2}=\ln n$ exactly and the
condition reads $L_x=o(1)$, 
\textbf{no integer $L_x\ge1$ satisfies it}. Equivalently, the window \eqref{eq:lxwindow} is empty at every
$s$, its right-hand side $9\alpha/\big(8(1-2\alpha)\big)$ diverging as $\alpha\uparrow\tfrac12$.
\cref{prp:phirate} therefore has no admissible sieve dimension, group (T-5) of \cref{thm:rootn} is
unsatisfiable, and the theorem is \emph{vacuous} at $d_Z=0$ rather than false.
\item \emph{Three further sites print ``$\alpha<\tfrac12$'' as a \textbf{reason} and are false as printed,
although each conclusion happens to survive at the boundary}: \cref{lem:prelimtheta}'s statement, which
writes $\alpha=r/(2r+d_Z)<\tfrac12$; that lemma's Step~6, whose collapse argument gives
$\kappa\alpha<\tfrac12$ from ``$\kappa<1$ and $\alpha<\tfrac12$'', true at $\alpha=\tfrac12$ from
$\kappa<1$ alone; and \cref{rem:h6bandwidth}'s ``impossible because $\alpha<\tfrac12$ always''. \emph{A
true conclusion resting on a false printed reason is the harder defect to find}, which is why they are
listed.
\end{itemize}
\emph{$d_Z=0$ is not a pathology and it is not being excluded as one}: it
says $\varrho$ has no non-degenerate argument, hence is a constant, hence is estimable at the parametric
rate, a genuinely easier propensity problem. \emph{What fails there is not the propensity but the
$\varphi$-step's entropy budget}, because $\hat\delta_n$ then sits at $n^{-1/2}$ and leaves no room between
the noise level and the sieve dimension. \emph{A block that assumes a strict inequality should say which
primitive delivers it}, and this is that primitive.

\medskip
\noindent
 The upper bound $d_Z\le p-1$ is a \textbf{modelling} restriction, and read with A1b(iv) it is
sharper than it looks. It is stated on one line and, unlike its companion above, has been discussed
nowhere; the asymmetry is not warranted, because the two bounds do different work. The lower bound
protects a rate. \emph{The upper bound restricts the specification of $\Delta Z$}, as follows.
\begin{itemize}
\item[--] $d_Z$ counts the \emph{non-degenerate} coordinates of $\Delta Z$, degenerate meaning almost
surely equal to a constant, which is what makes a coordinate carry no information for the propensity
regression, so $d_Z\le p-1$ says that \emph{at least one} coordinate of $\Delta Z$ is a.s.\ constant.
\item[--] \emph{That constant cannot be zero.} If coordinate $j$ is a.s.\ $0$ then row and column $j$ of
$E[\Delta Z\Delta Z']$ vanish and the matrix is singular, contradicting A1b(iv).
\item[--] And there cannot be two of them. If coordinates $j\ne k$ are a.s.\ the non-zero
constants $c_j$ and $c_k$, then rows $j$ and $k$ of $E[\Delta Z\Delta Z']$ are $c_jE[\Delta Z]'$ and
$c_kE[\Delta Z]'$, proportional, and the matrix is singular again.
\end{itemize}
\emph{So A4(ii) and A1b(iv) together force $d_Z=p-1$ \textbf{exactly}, and force $\Delta Z$ to contain
exactly one constant coordinate, non-zero: \textbf{the differenced equation must carry a period effect.}}
That is a restriction on how the analyst writes the model and not a regularity condition, and it is
substantive: a specification of $\Delta Z$ in which every coordinate varies is \emph{excluded} by A4(ii),
and one in which two do not is excluded by A1b(iv). \emph{What is \textbf{not} claimed is that the
period effect is non-zero in the population}, the coefficient on the constant coordinate is an
ordinary component of $\theta_\dagger$ and may be zero; what is forced is that the coordinate is
\emph{present}, so that $\Delta Z$ spans an intercept in differences. \emph{Nor is $d_Z=p-1$ a
strengthening of anything assumed}: it is a consequence of two clauses the block already maintains, which
is precisely why it had to be derived somewhere rather than left for a reader to notice.
% ADDED 2026-08-10 (Phase 6a).  The upper bound d_Z <= p-1 appeared at exactly one
% line in A4(ii) and at one glossary row ("because Delta Z's constant coordinate
% carries no information"), and its consequence -- read against A1b(iv) -- was
% printed nowhere.  It is not a counterexample to anything: no result is false, and
% the pair of clauses is satisfiable.  It is an UNDISCUSSED SCOPE RESTRICTION, and
% the Section 6 author has to know it before choosing Delta Z, because a
% specification with no constant coordinate does not satisfy A4(ii) as written.

\medskip
\noindent
\emph{The $\sigma$-algebra of the fold is written out and given no symbol here}, because
\cref{def:hatT}, which names it $\sigma(\mathcal I_a)$ by adopting \cref{lem:balance}'s convention, 
comes several hundred lines later in this file, and a maintained clause must not name an undefined object;
this is the write-it-out precedent C12/C24/C78, applied for the same reason A7's own statement forced
\cref{def:hatT} into Part~V.
\emph{The joint-measurability half is load-bearing}: two things this clause is asked to support are not
well posed without it.
\emph{First, its own display.} ``$O_P$'' is a statement about the distribution of
% ARCHAEOLOGY.  Through dispatch 15 this clause carried NO measurability statement
% of any kind, and the gap was recorded as OPEN; the half is new in dispatch 16.
$\omega\mapsto\|\hat\varrho-\varrho\|_{L^2}$, which is an \emph{integral of a random function}; joint
measurability makes it a measurable function of the sample, by Tonelli's theorem applied to the
non-negative jointly measurable integrand $(\omega,z)\mapsto\|\hat\varrho(\omega)(z)-\varrho(z)\|^{2}$
against the law of $\Delta Z$. Per-realisation Borel measurability of $z\mapsto\hat\varrho(\omega)(z)$
does not supply it. \emph{The integrand is a sum of $\mathsf L-1$ squares and is jointly measurable
exactly when each component is}, so the clause is read componentwise and the argument is the one it
always was; \emph{and the two norms in the display are different objects}, the inner
$\|\cdot\|$ is Euclidean on $\IR^{\mathsf L-1}$ and the outer $\|\cdot\|_{L^2}$ integrates it against
the law of $\Delta Z$, which is what A4(iii)'s ``$:=$'' spells out.
\emph{Second, \cref{lem:opnoise} Step~4}, which freezes $\hat\varrho$ on the auxiliary fold and applies
\textup{(P8)}'s freezing lemma. That lemma's hypothesis is joint measurability of the \emph{composite},
and the composite there is $\|\hat\varrho(\omega)(\cdot)-\varrho(\cdot)\|$ read at the main fold's
$\Delta Z$-draws, which the clause above supplies by composition with the Borel coordinate maps.
\emph{The shape is \cref{def:crossfit}\textup{(c$'$.1)}'s, one Part later}: a requirement on the
% ARCHAEOLOGY.  Through dispatch 15 lem:opnoise Step 4 was the one consumer of the
% freezing lemma whose hypothesis no clause of this block supplied; the repair was
% located here rather than papered over there.
estimator's own pointwise function and \textbf{not} a container declaration. That matters and is not a
stylistic preference, no equivalence class is formed and no representative is chosen, so
\cref{def:hatT}'s truncation, being a pointwise operation, continues to make
$\varepsilon_\varrho\le\big|\widehat{\widetilde W}_{i,l}\big|\le1$ hold \emph{surely} at every $i$ and
every $l$, and
\cref{def:crossfit}\textup{(a)}'s reading of that bound at the folds is exact rather than almost sure.
\emph{And it costs nothing about $P$}: it is a construction requirement on an estimator, satisfied by any
$\hat\varrho$ given by a formula jointly measurable in the auxiliary sample and the evaluation point, 
which every kernel and every series construction is, by inspection of the formula. \emph{The
Carath\'eodory route \cref{def:crossfit}\textup{(c$'$.1)} records for \eqref{eq:prodmeas} is available
here too but is \textbf{not} the same statement}: it delivers joint measurability from continuity in one
argument and measurability in the others, so here it would run on continuity of
$z\mapsto\hat\varrho(\omega)(z)$ at each realisation, true of a kernel or series construction and not
of a measurable selection, and in any case a sufficient condition and not the requirement, which is the
display above.

% ------------------------------------------------------------------------
\paragraph{A5 (nuisance rates).}

% EXTRACTED FROM THE STATEMENT OF A5(i) on 2026-08-04, together with the repair
% that made the extraction necessary.  The clause's two branches were written
% with O_P and its closing sentence was an unquantified identity; both are
% conditions-only violations of a different kind from the 2026-08-02 pass -- the
% first stated an upper bound where the corpus consumed a floor, the second
% stated an interpretation where the corpus consumed a margin.
\emph{Clause \textup{(i)} \textbf{declares} its two branches and \textbf{assumes} one thing, and the
two must not be run together.} What the clause assumes about the data-generating process is that the
operator error is $O_P(\hat\delta_n)$, the last equality of \eqref{eq:opnoise}, which
\cref{lem:opnoise} delivers rather than requires. What it declares is the level itself: $\hat\delta_n$ is
the noise level \emph{fed to} the discrepancy principle, and the stagewise schedule is computed from it
before the data are seen, so it is a deterministic sequence and the two superscripts are component tags
naming which summand of the maximum is meant, not estimators. \emph{A deterministic sequence does not
carry an $O_P$}, and this block's own preamble already says the clauses \emph{define} these rates.
% ARCHAEOLOGY.  Through 2026-08-03 the branches were written
% delta-hat_n^{(T)} = O_P(n^{-1/2} sqrt(ln n)) and delta-hat_n^{(varrho)} =
% O_P(r_n) -- UPPER bounds -- while FIVE sites read them two-sidedly to obtain a
% LOWER bound on delta-hat_n: prp:stopfeas(i)'s statement and its proof,
% lem:prelimtheta's Step 6, rem:lfnonlinear(3)'s mu = 1/2 disposition, and A4's
% d_Z >= 1 note in this file.  FROM X = O_P(a_n) ONE CANNOT INFER a_n <~ X.  A
% max of two O_P upper bounds is still an upper bound, which is exactly what let
% an implementation keep the propensity branch alone and stay formally
% consistent with the clause: rate_delta_n in app_tsp_common.R did so and
% printed a wrong number in the shipped d_Z = 0 feasibility row (corrected
% 2026-08-03).  The code is fixed; this clause is why the error was available.

\medskip
\noindent
 Nothing about the data-generating process is strengthened by the repair, and the floor is the
thing gained. The $O_P(\hat\delta_n)$ half is upper-only, as it was, and is still discharged by
\cref{lem:opnoise}; every ``A5(i) is discharged, not assumed'' downstream survives verbatim. What
changes is that
\[
\hat\delta_n\;\ge\;\hat\delta_n^{(T)}\;=\;n^{-1/2}\sqrt{\ln n}
\]
is now \textbf{arithmetic}, a property of a maximum whose branches are named, where before it was a
two-sided reading of an $O_P$. \emph{\Cref{lem:opnoise}\textup{(i)} cannot supply it and is no longer
cited for it}: that clause is a Chebyshev tail bound and bounds the empirical component from above.
The floor is what \cref{prp:stopfeas}\textup{(i)} runs Chebyshev against, and it is the reason a
statement proved at $\tfrac12\mathsf c\,n^{-1/2}\sqrt{\ln n}$ may be printed at $\mathsf c\,\hat\delta_n$.

\medskip
\noindent
 The clause's closing sentence is struck, and what replaces it is a quantifier. It read ``in every
configuration of interest $\hat\delta_n=\hat\delta_n^{(\varrho)}=O_P(r_n)$'', an identity with no
``eventually'' in it, and \emph{false as printed} over an admissible range. The propensity branch attains
the maximum exactly when $n^{1-2\alpha}>\ln n$, so the crossover sits near $n=5{,}500$ at
$\alpha=\tfrac38$ and \emph{diverges} as $\alpha\uparrow\tfrac12$; at $n=2{,}000$, the sample size
\cref{lem:scaleinv} prints an inflation factor at, the \emph{parametric} branch is the maximum, and the
shipped headroom ledger carries a branch column that records which branch binds cell by cell.
\emph{What survives is an eventual statement and it is a consequence of A4(ii), not a configuration
label}: $d_Z\ge1$ gives $\alpha=r/(2r+d_Z)<\tfrac12$ strictly, hence for each admissible configuration
there is a finite $n_\star$ with $\hat\delta_n=\hat\delta_n^{(\varrho)}=r_n$ at every $n\ge n_\star$.
\emph{That is precisely the licence clause \textup{(iii)} needs} in order to write
$\hat\delta_n=n^{-r/(2r+d_Z)}$ inside an asymptotic condition, and clause \textup{(iii)} now names it.

\medskip
\noindent
 It is \textbf{not} the licence the two sites that consumed the struck sentence as strict dominance
needed, and neither of them cites it any longer. \Cref{rem:l19cost}\textup{(2)} needs
$n\hat\delta_n^{2}$ to diverge \emph{polynomially}, which the floor does not give, the floor gives
$n\hat\delta_n^{2}\ge\ln n$, and it now carries $\alpha<\tfrac12$-with-a-margin as its own scope
condition on Step~6. \Cref{rem:h6bandwidth}\textup{(2)} needs only $\hat\delta_n\gg n^{-1/2}$, which the
floor delivers outright, in \emph{every} configuration and not merely where the propensity branch binds;
that site is strengthened by the repair rather than weakened.

\medskip
\noindent
\emph{The factor $\ln n$ is not optional}: \cref{prp:phirate} delivers \eqref{eq:phirate} at
$(\hat\delta_n\ln n)^{2\mu/(2\mu+1)}$, and (P-f),
\cref{rem:h6bandwidth}\textup{(0)} and \cref{prp:stopfeas}\textup{(iv)} all read that log-inflated
$\delta_n$. A log-free reading of this clause is off by $(\ln n)^{2\mu/(2\mu+1)}$.
\emph{The \textbf{rate} is what that proposition delivers and the \textbf{object} is not the object of
this clause}, a distinction that costs nothing here and is load-bearing two paragraphs over:
\eqref{eq:phirate} is stated at the infeasible reading $\hat\varphi^{\pi}$ of \cref{def:lfstep} from
2026-08-27, while this clause is stated at $\hat\varphi$, so the clause's first half is
\textbf{assumed} and not discharged --- and $\delta_n$ is the same deterministic sequence either way,
which is all this paragraph needs.
\Cref{rem:rootnlog} carries the full disposition, including which of the four rate products
% ARCHAEOLOGY.  The factor is new in dispatch 10 and is a CORRECTION, not a
% strengthening: rounds 1--9 wrote this clause at delta-hat_n^{2mu/(2mu+1)}, so the
% RATE the block STATED was not the rate the proposition returns, and the
% block was already inconsistent with itself, not merely with the proposition.
% ["DISCHARGES" -> "returns", 2026-08-27.  prp:phirate no longer discharges this
% clause at all: it is scoped to hat-varphi^pi and A5(ii) is stated at
% hat-varphi.  The RATE arithmetic in this paragraph is untouched by that.]
\cref{thm:rootn} needs are affected. In one line: the second half of this clause and clause (iii) are
\textbf{unchanged}, because \eqref{eq:frontier} is a strict inequality between polynomial exponents and its
strictness absorbs the logarithm; clause (iv) becomes \textbf{stronger} by the factor $(\ln n)^{2\mu/(2\mu+1)}$,
and that is the only place the logarithm bites; clause (v) becomes \textbf{weaker}, its right-hand side being
inflated; and \eqref{eq:muhalf} does not move, having been computed at the inflated $\delta_n$ already.

\medskip
\noindent
% EXTRACTED FROM THE STATEMENT OF A5 on 2026-08-02.  Clause (iv) carried the
% attribution of its own typing and clause (vi') carried the well-posedness
% argument for its own display, both inside the printed condition.  The TYPING
% stays in clause (iv) -- it is what makes the display a probability statement --
% and it is the ATTRIBUTION and the ARGUMENT that moved here.
\emph{Clause \textup{(iv)}'s fold and measurability are fixed by \cref{lem:balance}'s sample-splitting
convention, which also names what that measurability is \emph{of}.} \emph{The naming is not decoration.}
The
first display is an $O_P$ statement about $\omega\mapsto\|\hat\ell-\ell_\star\|_{L^{2}(P)}$, an integral of
a random function; joint measurability of the fixed representative makes it a measurable function of the
sample by Tonelli's theorem, and \emph{membership plus an unspecified ``measurability'' does not}.
\emph{Nothing here is strengthened}: the joint measurability
% ARCHAEOLOGY.  Through dispatch 15 the convention this clause points at stated
% exactly membership and a measurability whose target was not named, so the clause
% was well posed only by reference to something that did not state it -- the same
% defect as A4(iii)'s and lem:balance's own, written a third time and at one remove.
% The repair is at lem:balance and the naming half is dispatch 16's.
is scoped to what \cref{lem:balance} Steps~0 and~2 use, which is strictly less than
\cref{def:crossfit}\textup{(c$'$.1)}'s \eqref{eq:prodmeas}, because \cref{lem:balance} evaluates $\hat\ell$
at an undistorted $(\Delta U,\Delta Z)$ draw. \emph{In Part~I it is \eqref{eq:prodmeas} that governs and
not this convention}: what \cref{def:crossfit}\textup{(d)} reads at the folds is this clause's
\emph{display}, and the measurability of $\hat\ell^{(-k)}$ is \cref{def:crossfit}\textup{(c$'$.1)}'s own.
What this clause does and does not say about $\hat L$. It prices
$\hat\ell$, and nothing else. The identity
$\|\partial_u(\hat L-L_\star)\|_{L^{2}(P)}=\|\hat\ell-\ell_\star\|_{L^{2}(P)}$ is \textbf{an assertion,
not notation}, and is \emph{not} made here: it holds if and only if $\hat L$ is differentiable in $u$ at
every point with $\partial_u\hat L=\hat\ell$, that is, if and only if $\hat L=L^{\,\hat b}_{\hat\ell}$ in
the sense of \cref{def:weight}(b), which is exactly the construction obligation \cref{rem:dirmember}
records as owed by Step 6 or 7 and \emph{undischarged in this block}. \cref{lem:balance}(i)--(ii) consume
the clause in the form printed above, at $\hat\ell$ itself, and need nothing further;
\cref{lem:secondord}(iv) consumes it at $\|\partial_uD\|_{L^{2}(P)}$ with $D=\hat L-L_\star$, and may
therefore be instantiated at the estimated direction only \emph{after} that obligation is discharged.
% ARCHAEOLOGY.  Rounds 1--7 printed the identity as the first of two equalities
% inside the clause, which asserted the construction silently -- a forward
% obligation migrated into the block unstated.  Round-8 correction.

\medskip
\noindent
\emph{The clause is a statement about a
finite quantity, and what makes it one is \eqref{eq:sieveinhs} together with H1}:
$h=\hat\varphi-\vdagger$ with $\hat\varphi\in\mathcal S_n$ at the configuration
\cref{lem:pathcont} states, so $\hat\varphi\in H^{s}$ by \eqref{eq:sieveinhs} and $\vdagger\in H^{s}$
by H1's source set; without the first of those $\|h\|_{H^{s}}=+\infty$ surely and the
$O_P(1)$ statement would hold on the \emph{empty} event rather than with probability tending to one.

\medskip
\noindent
\emph{Clause \textup{(vi$'$)}'s display is a probability statement about a finite quantity only because
\textbf{both} memberships hold, and both are needed.} Every stage
$\hat\varphi^{(0)},\dots,\hat\varphi^{(\bar k)}$ lies in $\mathcal S_n$ once $\hat\varphi^{(0)}$ does, by
\cref{def:lfstep}\textup{(b)},\textup{(d)}, and hence in $H^{s}$ by \eqref{eq:sieveinhs}, \emph{while
$\vdagger\in H^{s}$ by H1's source set}. The display being a norm of a \emph{difference}, either
membership failing makes it $+\infty$ at every realisation and the event \textbf{empty} rather than
merely improbable.

\medskip
\noindent
\emph{That same $\mathsf R$ serves at every site \cref{def:srcball} enumerates, and \textbf{that} is
where the enumeration is maintained}: \cref{def:srcball} carries the list, and this clause,
\cref{lem:pathcont} and \cref{rem:paenvelope} point at it and maintain no copy.
\emph{``No second copy is maintained'' is what this sentence used to say and it was false}:
\cref{def:crossfit}'s clause-\textup{(d)} reading carries a three-entry instance list, which now
declares itself an instance list and points here --- the very failure the archaeology below
condemns, committed by the sentence recording it.
\emph{The quantifier is an existential over
% ARCHAEOLOGY.  Four copies of one list were maintained through dispatch 20, of
% lengths seven, five, five and four, so that a reader assembling the consumer set
% from any one of them got a different answer on a constant def:srcball declares
% "the radius that is not free".
one radius meeting all of those demands at once}, which is what \cref{lem:pathcont} consumes and what
\cref{rem:paenvelope} explains; an $O_P(1)$ reading does not serve, because \cref{def:srcball}'s ball must
be deterministic and an $O_P(1)$ bound is a family indexed by tolerance rather than a number.

% ------------------------------------------------------------------------
\paragraph{A8 (sieve window).}


\emph{The word ``consequently'' is a derivation, and it is printed here rather than asserted}: ten live
citations rest on it. Write $p_b$ for the degree, so $p_b>s-\tfrac12$; and write $s'$ for a
\emph{generic} order of the Sobolev scale, bound in this derivation and in H1's
boundary-condition paragraph and used nowhere else.
% ARCHAEOLOGY.  Until dispatch 22 the containment carried neither a derivation nor a
% pin while the weaker statement beside it carried two pins and was called
% load-bearing.  Both letters are dispatch-22 replacements and neither changes
% anything derived: the degree was bare italic d, one glyph from the dimensions
% d_Z, d_K and from the differentials dx, du, in the FIRST file of the compiled
% reading order; the generic order was written tau at two sites and sigma at three,
% against the reserved discrepancy constant and against the four-way sigma rule.
% Glossary C87 records both.
A polynomial spline of
degree $p_b$ on a mesh with simple interior knots lies in $C^{p_b-1}(\mathcal X)$
\citep{schumaker2007}; its $p_b$th derivative is constant on each mesh cell, hence a finite step function;
and a step function on a bounded interval lies in $H^{s'}$ for every $s'<\tfrac12$, and, unless it
is constant, i.e.\ unless the spline is a single polynomial, in no $H^{s'}$ with $s'\ge\tfrac12$.
Hence $\mathcal S_n\subset H^{s'}(\mathcal X)$ for every $s'<p_b+\tfrac12$, and the exponent
$p_b+\tfrac12$ cannot be raised uniformly over the span. Since $p_b>s-\tfrac12$ gives $s<p_b+\tfrac12$,
this is
\eqref{eq:sieveinhs}. Three ingredients of that derivation are load-bearing and each is now stated
rather than read in. The interior knots must be \emph{simple}: at an interior knot of multiplicity two
the spline is only $C^{p_b-2}$ there, its $(p_b-1)$st derivative already jumps, and membership drops by a
full order to $s'<p_b-\tfrac12$, so the containment fails at the stated degree.
% ARCHAEOLOGY.  The clause above did not say "simple interior knots" through
% dispatch 20, where the derivation was in a comment and the hypothesis it needs was
% in neither the clause nor the comment.
The space must be the \emph{classical} one, for the
reason H1 now records. And $p_b>s-\tfrac12$ is \emph{strict}, $s'=p_b+\tfrac12$ being
exactly where membership fails.

% ------------------------------------------------------------------------
\paragraph{the sample operator $\hat T$.}

Why the truncation is in the maintained block and not in a proof. Without it, $\hat\varrho$
taking the value of the realised dummy vector $\big(\un\{W=1\},\dots,\un\{W=\mathsf L-1\}\big)'$ has
positive probability, hence so does $\widehat{\widetilde W}_i=0$ in $\IR^{\mathsf L-1}$; on the
event that this holds at every $i\in\mathcal I_m$ the sample operator is identically zero and every
statement about it is vacuous. A4(i) constrains $\varrho$ and says nothing about
$\hat\varrho$, so the degeneracy cannot be excluded by assumption, it is excluded by construction,
which is where it belongs.

And the truncation set is the $\varepsilon_\varrho$-interior of the simplex \eqref{eq:trunset} and
not a box, for a second and separate reason. On \eqref{eq:trunset} the implied base cell
$\hat\varrho_0=1-\sum_{l\ge1}\hat\varrho_l$ is itself at least $\varepsilon_\varrho$, so $\hat\varrho$ is
a probability vector bounded away from every face; \emph{on the box
$[\varepsilon_\varrho,1-\varepsilon_\varrho]^{\mathsf L-1}$ it need not be one at all}, the implied base
cell being possibly negative, in which case $\hat\varrho_0$ names nothing, the \emph{estimated contrast
covariance}
\[
\hat\Omega_i\;:=\;\operatorname{diag}\big(\hat\varrho(\Delta Z_i)\big)-\hat\varrho(\Delta Z_i)\hat\varrho(\Delta Z_i)'
\]
need not be
positive semi-definite, and the overlap diagnostic $\min_l\hat\varrho_l$ the application reports is not
what it is taken for. \emph{$\hat\Omega_i$ is declared here, at its first use, and not in
\cref{def:hatT}}: the estimator does not evaluate it, \eqref{eq:hatT} carries the raw contrast, so a
declaration inside the definition would name an object that definition never uses. \emph{It is used by
the diagnostics and by the two clauses below}, which is why it is declared at all.
% ROUND-5 REPAIR.  Rounds 2-3 declared Omega-hat_i in def:hatT's PRINTED text,
% where the normalised eq:hatT consumed it.  Round 4 removed the normalisation
% and the declaration WITH it -- correctly, since an orphaned printed symbol in a
% definition is the round-1 defect -- but LEFT THE FOUR CONSUMERS, so the
% supplement used Omega-hat in printed text with no printed definition anywhere.
% That is the same defect mirrored.  The declaration returns HERE, at the first
% of the four consumers in file order, where it is used rather than merely named.
% TEXT RECOVERED 2026-08-10 (Phase 6a).  The sentence below had been SWALLOWED BY
% THIS COMMENT BLOCK -- spliced onto the end of the line above, so it did not
% compile and the paragraph began, in the PDF, with "\suppref{Remark}{hatTread}(2)
% hold on the box as well", a sentence with no subject.  Found by scanning for
% comment lines carrying printed markup; not assigned by any audit.
\emph{Be exact about what does and does not need this}: the two surely-bounds of
\suppref{Remark}{hatTread}\textup{(2)} hold on the box as well, since each uses only
$\hat\varrho_l\in[\varepsilon_\varrho,1-\varepsilon_\varrho]$. \emph{What needs \eqref{eq:trunset} is
that the fitted object be a probability vector}, which is a requirement on the first stage rather than a
bound any proof invokes. At $\mathsf L=2$ the box and
\eqref{eq:trunset} coincide, which is why no earlier round had to distinguish them.
% ROUND-4 CORRECTION.  Rounds 2-3 grounded this paragraph on the NORMALISATION:
% without the simplex, Omega-hat^{-1} would be undefined on an event, and eq:hatT
% was then written at the normalised contrast.  THE NORMALISATION IS REVERSED and
% that ground is gone.  The discovery stands; its reason is restated at the
% object the first stage returns.

% ------------------------------------------------------------------------
\paragraph{the smoothed sample linearisation.}

This fixes a reading, and the reading is substantive rather than notational.
$\hat{\mathcal A}_\varphi$ is a \textbf{separately smoothed operator} and \emph{not} the Fr\'echet
derivative of $\hat T$ in $\varphi$; the hat marks a sample analogue of $\mathcal A_\varphi$, not a
derivative of $\hat T$. \emph{The bandwidth is what forces the choice}: C3 attaches
$C_{h_x}$ to $\hat{\mathcal A}_{\mathrm{cond}}$ and $\hat{\mathcal A}^\ast_{\mathrm{cond}}$, and a
derivative of $\hat T$ would carry no bandwidth of its own. \emph{And the rejected reading is now known
to be unavailable rather than merely different}: by \cref{prp:hatAprop}(v), $\hat T(\cdot,\theta)$ is
nowhere Fr\'echet differentiable on $\mathcal V$, so ``the derivative of $\hat T$'' names no operator.

Forward obligation, DISCHARGED at \cref{def:hatAcon} and \cref{prp:hatAprop}, with one negative
answer. This definition names a construction; it does not perform one. Step~6 was to exhibit
$\hat{\mathcal A}_\varphi$ explicitly, in particular the smoothed indicator and the two kernel
estimates, and to verify the two conditions then maintained for the object exhibited.
\eqref{eq:hatAcon} exhibits it; A7 is verified in the amended form
\eqref{eq:a7amend} by \cref{prp:stopfeas}; and \emph{the tangential cone condition A6 is not verified
but refuted} (\cref{prp:tccfails}), which is why the block below carries it as retired.
Recorded in \cref{rem:obligations} as obligation (O4).

% ------------------------------------------------------------------------
\paragraph{A6 (RETIRED: tangential cone condition), moved entire.}

\begin{assumption}[A6: RETIRED, tangential cone condition for the \emph{random} operator]\label{ass:tcc}
\textbf{\sout{RETIRED.}} This assumption is not maintained, it is not available for use, and no
result in this programme takes it as a hypothesis. It read: \emph{there exist
$\eta_{\mathrm{TCC}}<\tfrac12$, not depending on $n$, and a ball
% ARCHAEOLOGY.  It was maintained through round 9.
$\mathcal V_n\subset\mathcal V$ around $\vdagger$ containing the Landweber--Fridman path, such that, with
$\hat{\mathcal A}_{\tilde\varphi}$ the smoothed sample linearisation of \cref{def:hatA},}
\[
P\Big(\big\|\hat T(\varphi,\theta)-\hat T(\tilde\varphi,\theta)-\hat{\mathcal A}_{\tilde\varphi}(\varphi-\tilde\varphi)\big\|_{L^2(\Pi)}
\;\le\;\eta_{\mathrm{TCC}}\,\big\|\hat T(\varphi,\theta)-\hat T(\tilde\varphi,\theta)\big\|_{L^2(\Pi)}
\ \ \forall\,\varphi,\tilde\varphi\in\mathcal V_n\Big)\;\longrightarrow\;1 .
\qquad\text{\sout{[retired]}}
\]
Why it is retired, and it is not a matter of strength. \Cref{prp:tccfails} proves that the
displayed event has probability zero, at every $n$, at every $\eta_{\mathrm{TCC}}<1$, at every ball
of positive $L^{2}(\pi)$-radius, and \textbf{for every bounded operator $\mathcal V\to L^{2}(\Pi)$ put
in place of $\hat{\mathcal A}_{\tilde\varphi}$}. An assumption whose event is null cannot be maintained
and cannot be repaired by lowering a constant or shrinking a ball.
\emph{What it should have said is \cref{def:a6b} (A6b)}, which pairs $\hat{\mathcal A}_\varphi$ with the
operator it actually differentiates, $\hat T_{h_x}$; A6b is a \emph{property}, on the A2b/A2c model,
never asserted and consumed by no theorem. \emph{Scope of the refutation, stated because it is
narrower than the retirement}: \cref{prp:tccfails} is a statement about $L^{2}(\pi)$-balls in
$\mathcal V$. On the spline sieve $\mathcal S_n$ the spike it uses does not exist, so a cone condition
\emph{restricted to $\mathcal V_n\cap\mathcal S_n$} is untouched by it, and would still not deliver
the inherited route, for the independent reason recorded at \cref{rem:tccwhy}\textup{(3)}.
\emph{Nothing downstream is affected}: \cref{prp:phirate} uses neither A6 nor A6b.
\end{assumption}

% ------------------------------------------------------------------------
\paragraph{C2$'$ (smoothness of the residual law).}

Two readings of that display are fixed here, \textbf{no letters are attached to them}, so that
nothing below can be cited as a sub-clause of C2$'$(i), because without them neither $E[\bar f_k]$ nor
$\|\bar f_k\|_\infty$ is a well-defined quantity.

First, the version, and why each $\bar f_k$ is measurable. The density is read at a version that
is \textbf{jointly Borel} in $(u,x_1,x_2,z,\mathsf w)$ and, at \emph{every} conditioning point rather than
almost every one, $p_K-1$ times continuously differentiable in $u$, which is what ``continuously
differentiable in its first argument'' is read to mean above. \emph{The supplier is one Part away and is now
cited rather than left implicit}: \textup{(P6a)} of \cref{def:pregver} fixes exactly such a version for the
conditional law of $\Delta Y_\theta-\Delta\varphi$ given $\mathcal F_{XZW}=\sigma(X_1,X_2,\Delta Z,W)$, for
every $\omega$, at $(\theta,\varphi)=(0,0)$ that law \emph{is} the one this clause states, and every other
is a translate of it by the shift identity of \textup{(P6)}, and the Doob--Dynkin lemma transports an
$\mathcal F_{XZW}$-measurable function of $\omega$ to
a Borel function of the conditioning \emph{value}, which is the argument $\bar f_k$ carries.
Those two supply the version \textbf{one variable at a time}, and the declaration above is a
\textbf{joint} one, so a third step is named rather than elided. \textup{(P6a)} gives joint measurability
in $(\omega,u)$; Doob--Dynkin runs at each \emph{fixed} $u$ and returns, for that $u$, a Borel function of
the conditioning value; neither, nor the two together, is joint Borel measurability in
$(u,x_1,x_2,z,\mathsf w)$. \emph{What delivers it is the rationals-plus-continuity step this paragraph
already performs below for the envelope}: at every conditioning point the map $u\mapsto\partial^{k}_uf$ is
continuous, and a
function that is Borel in one argument at each fixed value of the other and continuous in the other at each
fixed value of the first is jointly Borel, the Carath\'eodory criterion, reached here by taking $u$ along
the rationals and passing to the limit, which is the same device \cref{def:crossfit}\textup{(c$'$.1)}
runs for \eqref{eq:prodmeas} and the proof of \cref{prp:frechet} runs for $G_\varphi$.
\emph{Nothing that consumes this clause is affected}: every site reads $\bar f_k$ or
$\|\bar f_k\|_\infty$, and the sentences below build both from the version fixed above at each separate
conditioning point. \emph{The shortfall was in the \textbf{declaration} and not in the consumption}, and it is
repaired here because a declaration a reader cannot derive from its own stated suppliers is a gap
wherever it sits. \emph{It is also
the version the corpus already runs}: the Carath\'eodory argument opening the proof of \cref{prp:frechet}
establishes joint measurability of $G_\varphi$ and of its first two $u$-derivatives $G'_\varphi$ and
$G''_\varphi$, which \emph{are} $\partial_u^{k}f$ at $k=0$ and $k=1$, from measurability in $\omega$ and
continuity in $u$, and says there that it is used silently below. \emph{Extending it to $k\le p_K-1$ is the same argument and no more}: each $\partial_u^{k}f$ is the
pointwise limit of the difference quotients of $\partial_u^{k-1}f$ in $u$ taken along the reciprocals of the
positive integers, and a pointwise limit of countably many jointly Borel functions is jointly Borel, the
device \cref{def:crossfit}\textup{(c$'$.1)} uses for $\partial_u\hat\ell^{(-k)}$. \emph{Then the envelope is a
countable supremum}: continuity in $u$ makes $\sup_{u\in\IR}=\sup_{u\in\IQ}$ at every conditioning point, so
each $\bar f_k$ is Borel as a map into $[0,+\infty]$, $E[\bar f_k]$ is well defined in $[0,+\infty]$, and the
display above is the requirement that it be finite. \emph{Without this the display is an expectation of a
function not known to be measurable}, the ground on which \cref{lem:reprsmooth} Step~3 refuses the same
elision, and the same species as the defect repaired at A4(iii).

Second, the outer norm is an \textbf{essential} supremum, and the inner one is genuine.
$\|\bar f_k\|_\infty:=\operatorname{ess\,sup}\bar f_k$ under the law of $(X_1,X_2,\Delta Z,W)$, the supremum
in $u$ inside $\bar f_k$ remaining a genuine supremum, \emph{the same split
\cref{def:crossfit}\textup{(c$'$.2)} settles for \eqref{eq:supnormconv}, and settled the same way and for the
same reason}. \emph{The essential reading is what every consumer needs and is the weaker of the two}: in the
proof of \cref{prp:frechet} at Steps~3--5, in that of \cref{prp:compact}, in that of \cref{lem:adjfubini} and
at \cref{lem:aerr} Step~4,
$\|\bar f_0\|_\infty$ and $\|\bar f_1\|_\infty$ appear only on the right of an inequality \emph{inside} an
expectation or a conditional expectation, so an almost-sure bound is exact there and a sure bound would be
surplus strength. \emph{One thing does \textbf{not} carry over from \eqref{eq:supnormconv}}: there the outer
essential supremum needed a measurability argument of its own because the object inside it was random; here
$\bar f_k$ is a \emph{deterministic} function and $\|\bar f_k\|_\infty$ is a number, so no measurability
question arises for the outer norm and none is settled. What had to be settled was only which of the two
readings it is.

% ------------------------------------------------------------------------
\paragraph{C3 (kernel and boundary adjustment).}

\emph{The normalisation $\int\Theta_1=1$ is printed rather than left inside the words ``of order
$p_K$''}: it is consumed at \cref{lem:intker}(i), (iii) and (iv).
% ARCHAEOLOGY.  Rounds 1--11 left it implicit.

\medskip
\noindent
\emph{At $I=\IR$ clause (ii) \textbf{is} clause (i)}, so the interior kernel is $\Theta_1$ and there is
no second convention. \emph{At a point near the \textbf{right} endpoint of a compact interval no
separate clause and no reflection convention are needed}: the realised interval $I$ is then truncated
on the right rather than the left, $\Theta_I$ is the member of the family that interval indexes, and
\eqref{eq:truncmom} holds at it by hypothesis. \emph{The reflection is a
\textbf{theorem} rather than a convention} --- $\Theta_{-I}(\varsigma)=\Theta_I(-\varsigma)$,
\cref{lem:bdryfam}(f) --- \emph{and, being a theorem, it is not needed}: the interval already records
which side is truncated. \emph{The case a scalar index could not express, and which the old
\cref{lem:intker}(iv) excluded by hypothesis, is now covered by the same clause}: both endpoints close,
i.e.\ $I$ bounded on both sides.
% ARCHAEOLOGY.  Through dispatch 25 this paragraph read "At s_0 = 1 clause (ii) IS
% clause (i)" and then stated the reflection AS A CONVENTION, with the reflected
% kernel supported in [-c_C, c_C s_0].  Both halves are superseded by the interval
% indexing of dispatch 26; the reflection survives as lem:bdryfam(f), proved rather
% than declared, and the both-endpoints-close case, which no value of a SCALAR s_0
% could name, is now inside the family.

% ------------------------------------------------------------------------
%  MOVED HERE FROM THE ARTICLE-SIDE "Notes on C3" IN assumptions_main.tex,
%  2026-08-28, ON THE AUTHOR'S SHORTENING INSTRUCTION.  Provenance of constants.
%  What stayed in the article: that the envelope is not unique, the reported pair
%  (c_2 = 34.0, c_C = 0.155), and the two article-side displays it feeds.
% ------------------------------------------------------------------------

\medskip
\noindent
\emph{The envelope pair, and how it moves.} A smaller envelope returns a smaller
$\big(\mathsf c_2,\|\bar\Theta\|_{L^{1}}\big)$, which is the monotone object --- $c_C$ is a ratio and is
\textbf{not} monotone under pointwise domination. Beyond the relaxation parameter \eqref{eq:normconst}
and the cap \eqref{eq:cap}, the two constants enter $\mathsf c_3:=1\vee2c_C\mathsf c_2$ of
\cref{lem:intker}\textup{(i)}, increasingly; and they enter the certified radius of
\cref{prp:tccrepair}\textup{(ii)} \emph{reciprocally}, so that a loose envelope does not merely inflate a
printed number there --- it shrinks a sufficient condition towards vacuity. \emph{This list is
\cref{lem:bdryfam}'s, at its note to clause \textup{(g)}, and it is offered as exhaustive neither there
nor here}: the sentence it replaces read ``both enter three bounds'', named three, and was short by
$\mathsf c_3$. \emph{The Bernstein exponent $\mathsf c_{18}$ of \cref{prp:phirate} Step~1 is
\textbf{not} a second such case, and an earlier draft made it one}: it is a ratio of a mean bound to a
range bound, both quadratic in $\|\bar\Theta\|_{L^{1}}$, so the envelope cancels and it is invariant.

\emph{Why the envelope is taken minimal and not in the closed form that proves it finite.} The explicit
majorant inside \cref{lem:bdryfam}\textup{(d)} is admissible too, and returns
$\mathsf c_2=4.50\times10^{8}$ with $c_C=1.82$ --- the same mathematics at a cap larger by
$3.6\times10^{29}$ and a certified radius smaller by $4.5\times10^{7}$. \emph{There is a third envelope
between the two, and it is the only one whose pair is \textbf{proved} rather than measured}:
\cref{lem:bdryfam}\textup{(g)} exhibits a closed-form admissible envelope with
$\mathsf c_2=2^{1/4}e^{1/2}p_K^{2}=125.5$ and $c_C=\sqrt\pi\big(1+\operatorname{erf}(2^{-3/2})\big)=2.451$,
so $2c_C\mathsf c_2=615.2$ --- six orders inside the majorant on the product the downstream displays
read, \emph{their number is deliberately not printed}, no criterion for membership having been stated
--- and a factor $58$ outside the minimal envelope. \emph{Which pair a given reading should use is
decided by whether the inequality it feeds must be \textbf{proved} or may be \textbf{measured}}, and
that is \textup{(g)}'s subject.

% ------------------------------------------------------------------------
\paragraph{C4$'$ (residual truncation), moved entire.}

\begin{assumption}[C4$'$: Residual truncation]\label{ass:app4}
Let $\ell_n$ be a real sequence that is either bounded or diverges slowly to $\infty$ with $n$. The
marginal density $f_{\Delta U}$ exists, is continuous, and is strictly positive on $[-\ell_n,\ell_n]$
for every $n$; set $\kappa_n:=\inf_{|u|\le\ell_n}f_{\Delta U}(u)$, so that $\kappa_n>0$. If $\ell_n$ is
bounded then $\inf_n\kappa_n>0$.
\end{assumption}

% ------------------------------------------------------------------------
\clearpage

\paragraph{Discussion.}
A0 and A1 are the sampling and exogeneity primitives; A1 is the draft's own conditional-exogeneity primitive and is
the weaker of the two candidate exogeneity conditions, since it permits the instrument to be assigned on
characteristics that also enter $\Delta Z$. A1b(iii) rules out a second structural function generating
the same conditional law of the residual given $(W,\Delta Z)$; A1b(iv) is what identifies
$\theta_\dagger$, and it involves no source condition and no debiasing weight; A1c is the linearisation
of A1b(iii) and governs the local behaviour of the regularised inverse. B1 is a set of conditions on
objects the analyst \emph{chooses}, not on the data-generating process, with the exception of
B1(iii)--(iv), which restrict how coarse an index family may be used. \emph{B1(vi) is \textbf{retired} and is
not part of that set}: it was an identification restriction, it is not a condition of this estimator, and it
is not available as a hypothesis, see the struck clause and the reason recorded with it. A2 is the
Severini--Tripathi condition and is the binding
assumption for $\sqrt n$ \emph{through this moment function}: in the linear model its necessity for
$\sqrt n$-estimability is known \citep{severini2006}, but in the present nonlinear problem necessity
remains conjectural, and two claims must be kept apart, necessity for this moment, and necessity for
$\sqrt n$-estimability by \emph{some} estimator. A2$'$ adds what is needed for the second-order term to
be a finite number. A3 is retired, and it never ruled out ``a design in which the
debiased moment has no local information about $\theta_\dagger$''. By \cref{lem:momzero}\textup{(iv)} the debiasing
% ARCHAEOLOGY.  A3 was described here through dispatch 12 in those terms, which was
% an overstatement of what it could rule out.
term contributes nothing to the $\theta$-Jacobian at all, $E[\psi(\theta,\zeta_\star)]$ is affine with
gradient $-E[\Delta Z\Delta Z']$, so \emph{all} of the moment's information about $\theta_\dagger$
comes from restriction (B), which is what \cref{lem:shiftinv} and \cref{rem:R1} said in words from Step~1
onward and what A1b(iv) prices. \emph{The debiasing term buys orthogonality in $\varphi$, not information
about $\theta$}, and \eqref{eq:sigsplit} shows it is paid for in variance. A4 rules out a design in which
the instrument is degenerate on part of the support of the controls.
% THE THREE NUMBERS IN THE NEXT SENTENCE ARE THE TWO-CATEGORY READING OF THE
% APPLICATION AND ARE LEFT AS THEY STAND IN DISPATCH 25.  The plan's Sec.5
% replaces the binary trigger by a four-category one, at which point 689/300,
% varrho ~ 0.303 and the 18%/9% figures are computed for a design the paper no
% longer runs, and the correct replacement is the realised min_l hat-varrho_l and
% lambda_min(hat-Omega) named in the sentence after them.  THEY ARE NOT AMENDED
% HERE because the new split is an EMPIRICAL quantity this dispatch does not
% have -- inventing it would be worse than leaving a number a later phase
% replaces -- and because the arithmetic as printed is correct for the design it
% describes.  Flagged for Phase 9.
The application's $689/300$ split is not itself a difficulty: at $\varrho\approx0.303$ the
factor $\varrho(1-\varrho)\approx0.211$ costs $0.25/0.211\approx1.18$ relative to a balanced design, 
about $18\%$ in variance and about $9\%$ in standard error. A4(i) binds only if some
$\varrho_l$ approaches $0$
somewhere on the support of $\Delta Z$, which is a question about the conditional propensity and
not about the marginal split. \emph{At $\mathsf L\ge3$ that sentence is the one to read twice, because
the marginal split can look comfortable while a single cell is thin}: the binding quantity is
$\min_{0\le l\le\mathsf L-1}\varrho_l(\Delta Z)$ and, through \eqref{eq:omegainv}, one thin cell alone
sends $\|\Omega^{-1}\|$ to infinity whatever the others do. \emph{The realised
$\min_l\hat\varrho_l$ and $\lambda_{\min}(\hat\Omega)$ are therefore reportable diagnostics and not
optional ones}, and a design whose thinnest cell will not support a conditional multinomial fit against
$d_Z$ controls should be coarsened by merging cells rather than pushed through, but \emph{merging is not free,
and its cost runs in the direction opposite to the one a reader expects}. Since
$\mathcal V_{CI}(W)\subseteq\mathcal V_{CI}(g(W))$ for any coarsening $g$, identification at the
merged instrument \emph{implies} identification at the full one: merging therefore \textbf{adopts the
stronger} of the two hypotheses, and it forfeits the localisation the separated cells were carrying.
Both are to be stated whenever a design is coarsened, and neither is a formality. N1 is a \emph{property} of a debiasing weight, not a normalisation built
into it: it can always be arranged by re-anchoring, at no cost to $\partial_uL$ or to $E[\psi]$, and
\cref{prp:orthog}(iii) shows that it is \emph{necessary and sufficient} for exact orthogonality in the
propensity. Stating it as a property rather than as a construction is what leaves that equivalence with
content. \emph{At $\mathsf L\ge3$ the count is the thing to carry away, because a construction can
satisfy it at two cells and fail it at four without any display changing}: N1 is
$p(\mathsf L-1)$ scalar conditions and the re-anchoring that arranges it is $p(\mathsf L-1)$
conditional expectations, computed slice by slice, so centring \emph{one contraction} of the
$\mathsf L-1$ columns discharges N1 at $\mathsf L=2$ and leaves $p(\mathsf L-2)$ of it undischarged
above, at a cost of $r_n$ before the $\sqrt n$ scaling and $n^{1/2-\alpha}$ after it
(\cref{cor:centarity}). \emph{The realised
$\max_{j,l}\big\|E\big[\hat L^{(-k)}_{jl}\mid\Delta Z\big]\big\|_{L^{2}(P)}$, the largest of the
$p(\mathsf L-1)$ entrywise centring residuals of the fitted weight --- the object \textup{(O2)} writes
$E[\hat L-L_\star\mid\Delta Z]$, read entry by entry and estimated by a regression on $\Delta Z$ ---
is therefore a reportable diagnostic too}, on exactly the footing of the $\min_l\hat\varrho_l$ and
$\lambda_{\min}(\hat\Omega)$ named just above: all three are quantities a design either produces or
does not, and none of them is visible in a $\mathsf L=2$ rehearsal of the same estimator. A5 collects what the regularised first step must deliver; clauses (i)--(iii) are
discharged, not assumed, by \cref{prp:phirate}, and clause (iii) is the paper's real feasibility
frontier. \emph{Clauses (vi) and (vi$'$) are not discharged}: they are unproven inputs, and from Step~6
onward they are consumed on \emph{both} routes, (vi) by \cref{lem:balance}(iii) on the balancing
route, (vi$'$) by \cref{lem:pathcont} on the rate route, which needs the Sobolev bound at every stage
and not only at the final iterate.
\emph{A7 is the condition under which a discrepancy-principle stopping rule is a theorem rather than a
heuristic when the noise sits in the operator, and it is now the only one}: A6 is retired, having been
refuted for the operator it named (\cref{prp:tccfails}), and the stopping theorem
\cref{prp:phirate} runs on A7 in the amended form \eqref{eq:a7amend} together with the
stage-adaptive schedule the estimator computes, with \emph{no} cone condition anywhere. A8 fixes the
sieve window and, with it, the basis; from Step~6 onward its \emph{basis clause} and its \emph{upper
end} are consumed by the $\varphi$-step as well as by the balancing solve, through
\cref{lem:srctransfer}(iii) and \cref{lem:unifop}(3d), \emph{but its lower end is not}, and
\cref{prp:phirate} cites A8 by clause for that reason; (P-c) is the rate route's own lower end. C1--C5 are
regularity conditions for uniform consistency of the kernel components, carried over from the draft's
appendix and restated for the panel operator.

% A1b(ii) JOINED THIS CATEGORY 2026-08-05 and is recorded here rather than in the
% list below, which is scoped to Part VI.  Its only route into a result was
% lem:nusigma(ii), whose only consumer was the corollary withdrawn at
% rem:a1cwithdrawn.  Withdrawing a result orphans whatever entered only through
% it, and the consumer analysis has to be re-run when one is withdrawn -- it was
% not, and this clause is what that omission left behind.
\paragraph{Two clauses outside Part VI in the same position: A1b(ii) and A1c.}
Both are, as of the withdrawal recorded at \cref{rem:a1cwithdrawn}, clauses no theorem rests on;
both are maintained and printed; and neither is retired. They are collected here because the withdrawal
that orphaned them is recorded in the \emph{results} file while every comparable status note in this
programme sits in the assumption block, so a reader checking the block alone would not meet either.
\begin{itemize}
\item[--] \textbf{A1b(ii)} (measure separability) reached results only through
\cref{lem:nusigma}\textup{(ii)}, and that clause's only consumer was the withdrawn corollary. It is
\emph{not} retired, it is a substantive and checkable restriction on the joint law of $(X_1,X_2)$, and
it is the natural hypothesis of any future primitive for A1c, but nothing in the present programme
consumes it. \Cref{lem:nusigma}\textup{(ii)} carries the same note at its own site.
\item[--] \textbf{A1c} (the linearisation of A1b(iii), governing the local behaviour of the regularised
inverse) \emph{is a hypothesis of no result in this paper}, and its non-declaration occurrences are
\emph{mentions} and not consumptions --- the lists pairing it with A2, A$2'$ and H1 as statements
posed at $\mathcal A_{\mathrm{cond}}$, of which there are \textbf{three} and not two; the Discussion
sentence naming what it is; \cref{rem:a1cwithdrawn}, which withdraws its primitive;
\cref{rem:spacesread}, which names it as the completeness condition alongside the singular system;
and the closing line of \cref{prp:compact}'s proof, which records that the $e_j$ span a proper
subspace ``unless A1c holds''. \emph{No claim of exhaustiveness is
made, and that is deliberate}: the draft this sentence replaces counted three and the count is at
least six. It is \emph{not}
retired, and the distinction is the same one the five Part~VI clauses turn on: A1c is available as a
hypothesis and is a real restriction, it simply has no consumer. \emph{What it lacks is a primitive}:
\cref{rem:a1cwithdrawn} withdraws the only one offered, because that route's hypothesis set is empty in
this model, and records that A1b(iii) is therefore assumed and not derived. That is the status
note, and this is where a reader of the assumption block will look for it.
\end{itemize}
% A1c ADDED TO THIS REGISTER 2026-08-10 (Phase 6a), AND THE HEADING WAS RETYPED WITH
% IT.  The paragraph read "A clause outside Part VI in the same position: A1b(ii)" --
% singular -- so appending a second clause without moving the heading would have left
% a count contradicting its own list, which is the defect class this block has been
% bitten by four times.  The remedy for A1c already existed at rem:a1cwithdrawn; what
% did not exist was any pointer to it from the assumption block, and the block's own
% stated model (A5(v), and the five Part VI clauses) is that a maintained clause with
% no theorem behind it gets a printed status note WHERE THE CLAUSES ARE.

\paragraph{Five clauses of Part VI that no theorem rests on, and the status of each.}
These five are collected here because a maintained clause with no theorem behind it is exactly what a
reader cannot check by grepping, and because the block already does this for A5(v)
(\cref{rem:rootnscope}\textup{(4)}), that is the model, and the treatment below is uniform with
it. \emph{None of the five is retired, and the distinction matters}: this block has \textbf{five}
retirements resting on three grounds, A3 and A$2'$(ii) because they are \emph{proved}, A6 \textbf{and
A1d} because they are \emph{false}, B1(vi) because it is \emph{about a different estimator}, and none of
the five below meets any
of the three.
% COUNT CORRECTED 2026-08-26: "four retirements" -> FIVE.  A1d has carried a printed
% "RETIRED, 2026-08-05" banner in the assumption block since that date -- it is
% unsatisfiable in this model, Wtilde being conditionally centred given Delta Z, so
% the completeness condition fails at Lambda == 1 at every law A1 admits -- and this
% census never acquired it.  A1d sits on A6's ground, FALSE, not on A3's, PROVED, so
% the number of GROUNDS is unchanged at three.  The omission is the standing defect
% species here (a closure claim short by one), and it is the reason a reader checking
% A1d's standing against the block's own registers found it in none of them. Unconsumed is not retired, and only a retired clause is unavailable as a hypothesis. What each
of the five is instead:
\begin{itemize}
\item[--] \textbf{C1(i)} has exactly one consumer in the corpus, Step~2 of \cref{lem:aerr}, where it supplies
the $p_K$ derivatives of the two products $E[\breve g_{u,t}\mid X_j=\cdot]f_{X_j}$, \emph{the breve
marks the \textbf{infeasible} version of a hatted object, the one built at the true $\varrho$ rather
than at $\hat\varrho$; it is declared at \cref{lem:opnoise} Step~1, which is where the two stems that
carry it are defined, and this is its first occurrence in compiled order}, behind that step's
$O(h_x^{p_K})$ $x$-smoothing bias. \emph{That consumer is itself consumed by nothing}, \cref{lem:aerr}
opens ``Consumers: none'' and is retained for the two reasons its reading~\textup{(iv)} gives, so C1(i) is
a hypothesis of one lemma and, transitively, of no theorem. It is kept because that lemma's proof needs it.
\emph{\Cref{prp:phirate} does \textbf{not} consume C1(i)}.
% ARCHAEOLOGY.  A comment in this file asserted that it did, through dispatch 17.
\item[--] \textbf{C4$'$.} $\ell_n$ and $\kappa_n$ occur nowhere in this project outside C4$'$ itself, and no
object anywhere in these files is truncated to $[-\ell_n,\ell_n]$. The clause prices a residual truncation
the estimator does not perform. \emph{It is inert}: no bound stated anywhere degrades if it is deleted. It is
nevertheless not retired, because \emph{inert} and \emph{unavailable} are different registers, a truncated
variant of the residual direction would need exactly this clause, and \cref{def:a6b} is carried in this block
on the same footing, being ``never asserted and consumed by no theorem''.
\item[--] \textbf{C5$'$(ii).} No hypothesis list consumes it, \emph{and that statement is true only from
this dispatch}: \cref{lem:aerr}, \cref{prp:phirate} and \cref{rem:rootnhyp}\textup{(3)} block-cited C5$'$ and
therefore carried this clause, on the register rule stated with the retirements above, and all three now cite
C5$'$(i). \Cref{lem:aerr} reading~\textup{(i)} shows that the one operator-error bound stated at this pair of
bandwidths needs only $(nh_xh_u)^{-1}\ln n\to0$, which is what the clause now states.
\emph{What it does do is fix the lower end of the admissible $h_x$ window}, and two printed computations turn
on that end: \cref{rem:tccwhy}'s third reading,
which collects C5$'$(ii) with A8, \eqref{eq:entropysieve}, (P-c) and \cref{lem:unifop} Step~4 and concludes
that $h_xL_x$ is free in both directions, and \cref{lem:scaleinv}'s cap inflation, which reads $\bar m_n$
across the $h_x$ window that C5$'$(ii) and (P-g) leave. Both are computed at the strengthened form of the
clause and at $s=3$, where \eqref{eq:hxwindow} holds with room, so both stand. So C5$'$(ii) is consumed
\emph{arithmetically} and not as a hypothesis, and the two readings must be kept apart: an auditor grepping
hypothesis lists will find nothing, and deleting the clause would nonetheless leave two displays without a
lower limit. \emph{And a lower bound stated more strongly than any proof needs is \textbf{not} free here},
which is the correction \eqref{eq:hxwindow} records: it must intersect the upper bound (P-g) imports, and the
strengthened form and (P-g) have no common solution at any $s\le\tfrac12+18\alpha/(8-3\alpha)$, 
a band that is non-empty at every $\alpha\in(\tfrac14,\approx0.359)$, and \eqref{eq:frontier} forces
$\alpha>\tfrac14$.
% BOTH NUMBERS RECOMPUTED WITH (P-g)'s h_u.  The threshold's alpha-term scales by
% 9/8, so 16 -> 18.  The BAND is the set of s clearing eq:lxwindow and failing
% eq:hxwindow, non-empty iff 9 alpha/(8(1 - 2 alpha)) < 1/2 + 18 alpha/(8 - 3 alpha),
% whose root moves from ~0.3494 to ~0.35914; both are printed rounded DOWN, as the
% old number was, so that "non-empty at every alpha in the interval" stays true.
\item[--] \textbf{C5$'$(iii)} is a \emph{cross-reference} and not a restriction. It repeats A4(iii) and
\eqref{eq:frontier}, and every result that needs the propensity rate cites those two directly. Nothing cites
C5$'$(iii) and nothing should; it is here so that the smoothing-parameter block can be read on its own.
\item[--] \textbf{C5$'$(iv)} is a \emph{statement of scope} about the estimator, not a restriction on
$P$. It is
discharged by construction, \cref{def:hatTexp} and \cref{def:hatAcon} introduce $h_u$ and $h_x$ and no
bandwidth in $t$, in $\Delta Z$ or in $W$, so there is nothing for a proof to cite it for.
\end{itemize}
% RENAMED 2026-08-26, AND THE RENAME IS THE CHEAPER SIDE OF A COLLISION.  These two
% items read "C5'(iii) is a POINTER and not a restriction" and "C5'(iv) is a SCOPE
% DECLARATION about the estimator".  POINTER and DECLARATION are both TECHNICAL
% TERMS of the assumption block's register paragraph, where they mean, respectively,
% a locating citation inside a definition and a citation of a declarative half; here
% they meant a clause that restates another clause and a clause that describes the
% estimator rather than P.  TWO WORDS, FOUR MEANINGS.  The register pair is cited at
% roughly nine sites and is the load-bearing one, so THIS pair was renamed:
% "cross-reference" and "statement of scope".  Three sites moved -- these two and
% lem:aerr's statement in results_supp.tex, which reproduced both words.  DO NOT
% REVERT: "pointer" and "declaration" are reserved to the register paragraph.
Why a status sentence and not a retirement, applied to all five alike. A retirement is a claim that the
clause is unavailable as a hypothesis, and for C1(i) and C5$'$(ii) that claim would be false, the first is a
live hypothesis of \cref{lem:aerr} and the second is what puts the lower end on the $h_x$ window the two
printed computations above read. Retiring the remaining three alone would file clauses of one species under two
registers and leave a reader inferring the difference. The uniform remedy is the one A5(v) already
carries: state the standing, in one place, in the printed text.

\paragraph{The smoothness budget, in one place.}
Four smoothness parameters appear above and they do not compete for the same budget once H1 is imposed:
\begin{center}
% Wrapping columns, not `llll'.  At natural width this table runs 232pt past the
% text block; every cell below is prose or a display, so it wraps rather than
% needing a smaller body font alone.
\footnotesize
\begin{tabular}{@{}p{0.23\textwidth}p{0.29\textwidth}p{0.10\textwidth}p{0.28\textwidth}@{}}
\toprule
Condition & In terms of $(s,a,r)$ & Route & Source \\
\midrule
Mild ill-posedness\textsuperscript{\dag} & $s>a+\tfrac12$ & both & plan \S2.6 scope \\
$\mu_{\text{\textup{H1}}}>\tfrac12$\textsuperscript{\dag} & $s>a$ & both & implied by the line above, with room \\
Sieve window non-empty & $s>1$ (splines); $s>\tfrac32$ (polynomials) & \emph{balancing only} & A8 \\
Sieve conditioning ($L_x\ll n^{1/2}$) &, & both & A8 \\
Propensity frontier & $\dfrac{2\mu}{2\mu+1}\cdot\dfrac{r}{2r+d_Z}>\dfrac14$ & both & A5(iii) \\
Kernel bias & $u\mapsto\mathcal M(u,t)$ in $C^{p_K}$ uniformly in $t$, $p_K=8$ & both & C2$'$(ii) \\
Representer & source index $\mu''>\tfrac32$ for $\ell^\Pi_\star$, independent of $\mu$ & both & A2c, \cref{lem:reprsmooth} \\
\emph{Linearisation remainder} & $\mu>\tfrac56$ \emph{strict}, at C5$'$(i) as amended & \emph{rate only} & (P-f), \cref{rem:h6bandwidth} \\
\emph{Sieve entropy} & $L_x\ln n=o\big(n\hat\delta_n^{2}\big)$ & \emph{rate only} & (P-d), \eqref{eq:entropysieve} \\
\emph{$x$-smoothing bias} & $\mathsf R\,h_x^{s-1/2}=O\big(h_u\hat\delta_n\big)$ (an \emph{upper} end for $h_x$; the $h_u$ is what Step~4 charges in operator norm) & \emph{rate only} & (P-g), \cref{lem:unifop} Step 4 \\
\emph{Sobolev stability, stagewise} & $P\big(\max_{k\le\bar k}\|\hat\varphi^{(k)}-\vdagger\|_{H^{s}}\le\mathsf R\big)\to1$ & \emph{both} & A5(vi$'$) \\
\midrule
\emph{$L_x$ window, rate route} & $s>\dfrac{9\alpha}{8(1-2\alpha)}$ & \emph{rate only} & (P-c)$\,\wedge\,$(P-d), \eqref{eq:lxwindow} \\
\emph{$L_x$ window, both routes at once} & $s>\dfrac{1}{2(1-2\alpha)}$ & \emph{both, jointly} & A8 lower end $\wedge$ (P-d) \\
\midrule
\emph{Kernel order vs.\ Sobolev order} & $p_K\ge s-\tfrac12$, i.e.\ $s\le\tfrac{17}{2}$ at $p_K=8$; against \eqref{eq:lxwindow} this restricts $\alpha<\tfrac{68}{145}$, i.e.\ $9r<68\,d_Z$ & \emph{rate only} & \cref{lem:unifop} Step 4, gates (P-g), (N1) \\
\emph{Source index, now primitive} & $\mu>\tfrac56$ at (P-a.3), \textbf{not} derived from $s/(2a)$ & \emph{rate only} & (P-a.3), \cref{rem:pasrc} \\
\bottomrule
\end{tabular}
\end{center}
\noindent\textsuperscript{\dag}These two lines carry $\mu_{\text{\textup{H1}}}=s/(2a)$, H1's index at
the \textbf{population} operator $\mathcal A_{\mathrm{cond}}$ on all of $\mathcal V$ at $\theta_\dagger$, 
\textbf{not} the primitive $\mu$ of (P-a.3), which is what every other $\mu$ in this table and in
\eqref{eq:frontier} means. They are \emph{interpretations} and transfer to the primitive index only under
the commutativity proviso \cref{rem:pasrc}\textup{(4)} states and declines. \cref{rem:pasrc}\textup{(5)}
carries the binding reading.

\noindent The table is read at the \textbf{stacked} operator, and that is what makes it exact at every
$\mathsf L$ rather than merely safe. $a$ is the exponent H1 assumes for $\mathcal A_{\mathrm{cond}}$ on
the whole three-factor range space, so at $\mathsf L\ge3$ the $a$ in the two daggered lines is already the
right object and \emph{no line of this table needs re-deriving}; what that operator is, and what it is
not, is the Discussion paragraph \emph{H1}. Those two lines are also the only ones that condition on
$a$ at all. Every other row restricts some combination of $s$, $\mu$, $\mu''$, $L_x$, $h_x$, $p_K$,
$\alpha$ and $r$, and none of them $a$; and $\mu$ has been
\emph{primitive} at (P-a.3) since the identity was struck, so \eqref{eq:frontier}, \eqref{eq:muhalf}
and the $\mu>\tfrac56$ line do not move with $a$ in either direction, and the last row names $s/(2a)$ only
to deny that $\mu$ is derived from it. \emph{The operative content of this table is therefore free of
$\mathsf L$}, and the flagship configuration $(\mu,s,a,\alpha)=(2,3,\tfrac34,0.32)$ below is a statement
about four numbers that does not care how many contrasts produced $a=\tfrac34$.
\emph{There is one reading a reader can be misled by}, and naming it is the whole of what
$\mathsf L\ge3$ changes here: someone who imports an exponent measured on a \emph{single} contrast, $a_l$
for cell $l$ against the base, and reads $s>a+\tfrac12$ and $s>a$ at that $a_l$, is reading them
\emph{conservatively}, under per-contrast polynomial decay the stacked exponent is
$\min_la_l\le a_l$, and both lines weakly relax as $a$ falls. \emph{That is a statement about the reading
and not about the table}; it is conditional on per-contrast exponents existing at all, which H1 does not
assume; and it says nothing about a \emph{merged} binary design, whose operator is a \emph{sum} of
contrast operators and none of them. The subsumption argument below inherits exactly this and nothing
more: it turns on $a\le\tfrac34$, which a smaller $a$ makes weakly easier to meet, and it already runs
under the proviso \cref{rem:pasrc}\textup{(4)} declines.

\noindent Two readings of the table are needed rather than one. \emph{The ``Route'' column is not decoration:} by
\cref{lem:balance} the first-order bias is controlled by \emph{either} A5(ii)+A5(iv) \emph{or}
A8+\eqref{eq:nuhat}+A5(vi), not by both, so the $s>1$ line is charged on the balancing route only. On the
rate route the $s$ budget faces the first two lines, the frontier, the representer condition and
\emph{the last two rows of the table}, and \emph{nothing from A8's lower end}, which is why
\cref{prp:phirate} cites A8 by clause, taking its basis and its upper end and \emph{not} its lower end.
The upper end of A8, which is the conditioning constraint on the
sieve solve, is charged on both routes and carries no $s$ requirement. The trade-off between
$\mu$ and $r$ in the frontier should be displayed rather than hidden: at
$(\mu,d_Z)=(1,3)$ the frontier reads $r>4.5$, at $(1,5)$ it reads $r>7.5$, and at $(2,3)$ it reads
$r>2.5$.

The representer line is the one that changed at Step 4, and the change is a discharge rather than a
budget. Through round 8 it read ``source index $\mu''$ for $\ell^\Pi_\star$'' and pointed at
A2$'$(i), a \emph{primitive}, with no route to it. It now points at A2c and at
\cref{lem:reprsmooth}: the index is pinned at $\mu''>\tfrac32$, which is the one-dimensional
$C^{1}$-embedding threshold and not a free choice, and the line records a \emph{sufficient condition}
rather than an assumption. It still competes for no part of the $s$ budget, because $\mu''$ indexes
smoothness of $\ell^\Pi_\star$ in the residual and $s$ indexes smoothness of $\vdagger$ in $x$.

The four lines added at Step 6 should be read as one addition and three bookkeeping entries.
\emph{The addition is $\mu>\tfrac56$}, and it \emph{subsumes} the $\mu_{\text{\textup{H1}}}>\tfrac12$ line,
which is therefore now slack rather than binding: under H1 and the scope condition $s>a+\tfrac12$ one has
$\mu_{\text{\textup{H1}}}>\tfrac12+\tfrac1{4a}$, so $\mu>\tfrac56$ is \emph{implied} whenever
$a\le\tfrac34$, \emph{under the interpolation proviso of \cref{rem:pasrc}\textup{(4)}, without which the
computed $\mu_{\text{\textup{H1}}}$ is not the $\mu$ of \eqref{eq:muhalf}}; see the dagger note above, and is a
genuine extra restriction otherwise. \emph{It is strict}, and the strictness is not decoration: at the
\emph{boundary} configuration $(\mu,d_Z,r)=(1,3,4.5)$ the underlying condition (P-f) holds in powers of
$n$ with no margin and \emph{fails} once the logarithms A7's amended threshold introduces are counted
, see \cref{rem:h6bandwidth}, clauses (1) and (5). \emph{The three bookkeeping entries} are the
sieve-entropy line (which replaces the stronger Sobolev-entropy condition \eqref{eq:entropy}, not
consumed by any result), the $x$-smoothing line (which is the block's \emph{only} upper bound on $h_x$
and is currently carried as a proposition hypothesis rather than as a clause of C5$'$), and the
stagewise Sobolev line (which is A5(vi$'$), an input and not a discharge). \emph{None of those four
competes with $s$ against the frontier}; they compete with $\mu$, with $L_x$ and with $h_x$
respectively.

Two further lines \textbf{do} compete with $s$. (P-c) and (P-d) are the two ends of one window for
$L_x$, and reading
% ARCHAEOLOGY.  Through round 10 the text asserted that "the s budget is unchanged";
% it is not.  The two lines were added in the dispatch-8 pass.
them together at $h_u\asymp\hat\delta_n^{1/p_K}$ gives $n^{9\alpha/(8s)}\ll L_x=o(n^{1-2\alpha}/\ln n)$,
which is non-empty if and only if \eqref{eq:lxwindow}, $s>9\alpha/\big(8(1-2\alpha)\big)$. That is
$s>\tfrac{27}{16}=1.6875$ at $\alpha=\tfrac38$ and $s>1$ at $\alpha=0.32$: satisfiable, but at
$\alpha=\tfrac38$ it consumes most of the available $L_x$ margin, the admissible window there being
only $\big(n^{0.2411},\,n^{0.25}/\ln n\big)$. \emph{And if A8's \textbf{lower} end is carried as well
, which it must be by any configuration on which both the rate route and the balancing route are to
be available, the requirement is instead $s>\big(2(1-2\alpha)\big)^{-1}$, i.e.\ $s>2$ at
$\alpha=\tfrac38$}: at $\alpha<\tfrac49$ the two thresholds differ by
$(4-9\alpha)/\big(8(1-2\alpha)\big)>0$ and A8's is strictly the larger.
\emph{The practical consequence, and it is the reason the line is recorded rather than absorbed}: the
boundary configuration $(1,3,4.5)$ was priced at $s>1.75$ in earlier rounds, which clears
\eqref{eq:lxwindow} and does \emph{not} clear A8 in full. It is therefore not an admissible design for
a paper that keeps both routes open, independently of the fact that it already fails (P-f) by
$(\ln n)^{1/3}$. The frontier-clearing configuration $(\mu,s,a,\alpha)=(2,3,\tfrac34,0.32)$ clears
every line in this table. \emph{Read that configuration as \textbf{four} numbers and not as three plus a
derivation}: $\mu=2$ is (P-a.3)'s primitive index, assumed at the compression of the realised operator at
$\theta_\circ$, while $(s,a)=(3,\tfrac34)$ fix the scale and the population decay at $\theta_\dagger$; that
$s/(2a)=2$ \emph{coincides} with the assumed $\mu$ is a property of the numbers chosen and is not an
identity this block can invoke, the bridge being the proviso \cref{rem:pasrc}\textup{(4)} declines. \emph{The
coincidence is deliberate and is worth keeping}, a configuration on which the two indices agree is the
one a reader will find least surprising, but it is not what makes the configuration admissible.
\emph{Mechanised at \texttt{scripts/symbolic/specs/phirate-sieve-joint-window.yaml}.}

Two lines were added in round 12, and the second changes how the whole table is read.
% CORRECTED 2026-08-05.  This sentence read "so it binds at no configuration this
% block contemplates", which rem:phirhyp(2) refutes in the same corpus with an
% admissible witness.  Both statements were compiled.  The cap is not slack.
\emph{The first is a budget line that binds}: $p_K\ge s-\tfrac12$, which \cref{lem:unifop}
Step~4's Taylor expansion consumes and which no earlier round stated. At $p_K=8$ it caps
$s\le\tfrac{17}{2}$, and \emph{against \eqref{eq:lxwindow}'s floor that cap empties the hypothesis set
at every $\alpha\ge\tfrac{68}{145}$}, as \cref{rem:phirhyp}\textup{(2)} works out with an admissible
witness; it is a condition,
it gates (P-g) and (N1), and it was previously supplied by an appeal to
$H^{s}\hookrightarrow C^{0,s-1/2}$ that is \textbf{vacuous} for every $s>\tfrac32$ and hence at every
admissible configuration. The rate $\mathsf Rh_x^{s-1/2}$ survives; its justification did not.

The second is that $\mu$ is no longer a function of $(s,a)$ on the rate route. (P-a.3) states the
source condition \textbf{spectrally}, at the compression \eqref{eq:compress}, with $\mu$ and
$\mathsf R_\ast$ primitive; the identity $\mu=s/(2a)$ is \emph{struck} from \cref{prp:phirate}'s
hypothesis line and re-appears in \cref{rem:pasrc}(4) as an interpretation under an explicitly stated
interpolation proviso. Three consequences for reading the table. (i) The frontier line and the
$\mu>\tfrac56$ line are now conditions on a \emph{primitive} $\mu$, so they no longer translate
automatically into conditions on $s$, the entry ``$\mu=s/(2a)$'' has been removed from the frontier
row. (ii) \emph{H1 is not retired and the rest of the table is unchanged}: H1 still fixes the scale
$H^{s}$ that \cref{def:srcball}, A5(vi$'$), (P-c) and (P-g) are written in, still supplies A2's Picard
form, and still delivers the mild-ill-posedness and $\mu>\tfrac12$ lines \emph{as an interpretation}.
(iii) \emph{The honest reading of the collapse}: H1's matching is a statement about
$\mathcal A_{\mathrm{cond}}$ on all of $\mathcal V$, while (P-a.3) is about a finite-rank, sieve-compressed
operator at a random base point, and \cref{lem:srctransfer}(iii.a) records that a link condition
\emph{cannot} hold on $\mathcal V$. So H1 never delivered $\mu=s/(2a)$ for the operator (P-a) is about,
and round 12 makes visible an assumption the block was already carrying rather than adding one. See
\cref{rem:pasrc}, all four clauses.

% ---------------------------------------------------------------------------
%  THE REMARKS BELOW WERE EXTRACTED FROM ASSUMPTION STATEMENTS ON 2026-08-03,
%  under the house style (.claude/references/theorem-environment-style.md): an
%  ASSUMPTION CARRIES CONDITIONS ONLY, and motivation, worked arithmetic, scope
%  notes and "what this rules out" go after it.  NOTHING WAS DELETED: every
%  sentence removed from a statement is below, verbatim except for the connectives
%  and the splice-point punctuation needed to make it stand outside the
%  enumeration it used to sit in.
%
%  THEY ARE AT THE FOOT OF THE FILE ON PURPOSE.  Nothing printed follows them
%  here, so they renumber nothing in this file; and this file is input after the
%  four shared files, so they shift no number the article and the supplement have
%  to agree on.  They are the first numbered environments in this file, which is
%  why the placement is stated rather than assumed.
%
%  ASSUMPTIONS ARE STILL CITED BY TAG.  \begin{assumption} is \newtheorem*, so a
%  bare \ref{ass:...} prints nothing; the remarks below name C5', A2c, C4' and C3
%  in prose and never \ref them.
% ---------------------------------------------------------------------------

\begin{remark}[How the smoothing-parameter block C5$'$ is read]\label{rem:app5read}
Two clauses, both of which stood inside the statement of C5$'$. Neither is a condition, and the
block restricts exactly the same bandwidth sequences without them.
\begin{enumerate}
\item[(1)] \emph{What clause \textup{(i)}'s display buys.} It is chosen so that $h_u\to0$ and the smoothing
bias $h_u^{p_K}$ is $\asymp\hat\delta_n$. No lower bound on $h_u$ is required for this component.
\item[(2)] \emph{Clause \textup{(ii)}'s displayed form, the stronger inherited form, and the window they
leave.} The displayed form is the one \cref{lem:aerr} consumes. \emph{The stronger inherited form
$\big(nh_x^{2}h_u^{3}\big)^{-1}\ln n\to0$ remains available, and is what the two arithmetic readings
recorded elsewhere in this Discussion are computed at, but it is imposed \textbf{only} where it is
compatible with an upper bound on $h_x$}, because clause (ii) is a \emph{lower} bound and the block
carries no upper bound at
all: the upper bound is \cref{prp:phirate}'s hypothesis (P-g),
$\mathsf R h_x^{s-1/2}=O\big(h_u\hat\delta_n\big)$.
At the tuning of clause (i), $h_u\asymp\hat\delta_n^{1/p_K}=n^{-\alpha/8}$, $\hat\delta_n=n^{-\alpha}$,
$p_K=8$, so that $h_u\hat\delta_n=n^{-9\alpha/8}$, the stronger form reads
$h_x\gg n^{-(1-3\alpha/8)/2}$ and (P-g) reads
$h_x=O\big(n^{-9\alpha/(8(s-1/2))}\big)$, so the two leave a non-empty window \emph{if and only if}
\begin{equation}\label{eq:hxwindow}
s\;>\;\tfrac12\;+\;\frac{18\,\alpha}{8-3\alpha} ,
\end{equation}
whereas the form displayed in the statement reads $h_x\gg n^{-(1-\alpha/8)}$ and is compatible with (P-g)
as soon as
$s>\tfrac12+9\alpha/(8-\alpha)$, \emph{the \textbf{displayed-form threshold}, which is what the block
actually imposes}, hence, since $\alpha<\tfrac12$ by A4(ii), at \emph{every} $s>\tfrac{11}{10}$,
uniformly in $\alpha$.
% BOTH THRESHOLDS RECOMPUTED WHEN (P-g) ACQUIRED ITS h_u, AND THE ARITHMETIC IS
% MECHANICAL: (P-g) is R h_x^{s-1/2} = O(h_u hat-delta_n) = O(n^{-9 alpha/8}), so
% h_x = O(n^{-9 alpha/(8(s-1/2))}) rather than O(n^{-alpha/(s-1/2)}), and BOTH
% non-emptiness thresholds have their alpha-dependent term multiplied by 9/8:
% 16 alpha/(8 - 3 alpha) -> 18 alpha/(8 - 3 alpha) and 8 alpha/(8 - alpha) ->
% 9 alpha/(8 - alpha).  The uniform bound is sup_{alpha < 1/2} of the second,
% attained in the limit: 9(1/2)/(8 - 1/2) = 4.5/7.5 = 3/5, so 1/2 + 3/5 = 11/10.
% It was 8(1/2)/7.5 = 8/15 and 1/2 + 8/15 = 31/30.
And the strengthening costs this block no configuration, which is worth stating because it is not
obvious from the thresholds alone. \Cref{rem:h6bandwidth}\textup{(5)}'s two running configurations are
the frontier-clearing $\alpha=0.32$ and the boundary $\alpha=\tfrac38$. At those, the displayed-form
threshold reads $s>0.8750$ and $s>0.9427$, against $s>0.8333$ and $s>0.8934$ before, 
and \eqref{eq:hxwindow} reads $s>1.3182$ and $s>1.4818$, against $s>1.2273$ and $s>1.3727$. \emph{The
frontier-clearing configuration fixes $s=3$ and clears both, at $+2.13$ against the displayed-form
threshold and $+1.68$ against \eqref{eq:hxwindow}}; the boundary configuration fixes only $(\mu,d_Z,r)$ and
leaves $s$ free, so it is quoted at the same $s=3$, where it clears at $+2.06$ and $+1.52$, and it still
clears at $s=2$, the value below which \cref{rem:h6bandwidth}\textup{(5)} records its sieve window as
empty, with margins $+1.06$ and $+0.52$. So no conclusion of this corpus is lost to the
correction, and the reader who wants the one-line version has it: at $s=3$ nothing in this block moves.
\emph{One printed witness did have to move with it}, \cref{rem:rootngroups}\textup{(5)}'s, which is
chosen to sit \emph{below} \eqref{eq:re2window} and had to be raised to stay \emph{above}
\eqref{eq:hxwindow}, and it is the only one.
\end{enumerate}
\end{remark}

\begin{remark}[What A2c is, and what it is for]\label{rem:reprsourceread}
One clause, which stood inside the statement of A2c. It is a status note and not a condition, and
A2c restricts exactly the same representers without it.
\begin{enumerate}
\item[(1)] \emph{A2c is NOT MAINTAINED.} It is a hypothesis \cref{lem:reprsmooth} imposes on a representer it
quantifies over, on the model of A1d, whose sole consumer was the withdrawn corollary of \cref{rem:a1cwithdrawn}. Its only role is to
\emph{imply} A2$'$(i) at $\ell^\Pi=\ell^\Pi_\star$, where the budget is written
$\Lambda^{(\mu'')}_\star$, and it is never asserted. \emph{A2c is stated in the same
$\bar w$-weighted $L^{1}(\Pi_t)$ geometry as A2$'$(i.a)--(i.b) themselves}, which is what
makes \cref{lem:reprsmooth}(iii) a one-line integration of a slicewise inequality rather than an
interchange of a supremum with an integral.
\end{enumerate}
\end{remark}

\begin{remark}[How \cref{def:hatT}'s two construction conventions are read]\label{rem:hatTread}
Four clauses, each of which stood inside the definition. None of it is construction:
\eqref{eq:foldbal}, the truncation, the definition of $\widehat{\widetilde W}_i$ and \eqref{eq:hatT} are
unchanged by its removal.
\begin{enumerate}
\item[(1)] What the balanced split costs, and where it is consumed.
\emph{This is a construction convention and it costs nothing, the analyst chooses the split, and
$n_{\mathrm m}=\lceil n/2\rceil$ serves, but it is \textbf{used}}: every variance bound for the
$\varphi$-step is
denominated in $n_{\mathrm m}$ and every
rate in $n$, so the two must be tied. It is consumed at \cref{lem:opnoise} Step~5,
\cref{prp:stopfeas}(i) and \cref{prp:phirate} Step~1 (N2.3), each of which converts an
$n_{\mathrm m}^{-1/2}$ or $n_{\mathrm m}/L_x$ into the corresponding statement in $n$. Without
\eqref{eq:foldbal} no bound for the $\varphi$-step is false, but nothing is stated in $n$ either.
\item[(2)] Three things follow from the truncation and all three are used. (a) The truncation can only decrease
$\|\hat\varrho-\varrho\|_{L^{2}}$, because $\varrho(\Delta Z)$ lies in \eqref{eq:trunset}
almost surely by A4(i) and the metric projection onto a closed convex set containing the target
is a pointwise contraction; so A4(iii) is preserved. \emph{Convexity is what the argument uses, not
one-dimensionality}: at $\mathsf L=2$ the set \eqref{eq:trunset} is the interval
$[\varepsilon_\varrho,1-\varepsilon_\varrho]$ and the sentence is the one earlier rounds printed. (b) With
$\widehat{\widetilde W}_{i,l}:=\un\{W_i=l\}-\hat\varrho_l(\Delta Z_i)$ as \cref{def:hatT} defines it, we have
$\varepsilon_\varrho\le\big|\widehat{\widetilde W}_{i,l}\big|\le1$ \emph{surely}, for every $i$ and every
$l$, because $\un\{W_i=l\}\in\{0,1\}$ while $\hat\varrho_l(\Delta Z_i)\in[\varepsilon_\varrho,1-\varepsilon_\varrho]$:
the two cases $W_i=l$ and $W_i\ne l$ give $1-\hat\varrho_l\ge\varepsilon_\varrho$ and
$\hat\varrho_l\ge\varepsilon_\varrho$ respectively. \emph{This is the degeneracy guard}: a row and a
column of \eqref{eq:gram} vanish only if $\widehat{\widetilde W}_i$ vanishes as a vector, which the bound
excludes at every $l$. (c) \emph{$\hat\varrho$ is a probability vector}, its implied base cell
$\hat\varrho_0=1-\sum_{l\ge1}\hat\varrho_l$ being at least $\varepsilon_\varrho$, which is what
\eqref{eq:trunset} buys over the box $[\varepsilon_\varrho,1-\varepsilon_\varrho]^{\mathsf L-1}$, on
which the implied base cell may be negative. \emph{Clauses \textup{(a)} and \textup{(b)} would hold on
the box too}, each using only $\hat\varrho_l\in[\varepsilon_\varrho,1-\varepsilon_\varrho]$; what needs
\eqref{eq:trunset} is that $\hat\varrho$ estimate a probability vector at all, hence that $\hat\varrho_0$
and $\hat\Omega$ mean anything and that the reported $\min_l\hat\varrho_l$ be the overlap diagnostic it
is taken for. \emph{A consequence with no consumer in the estimator as adopted}, recorded because the
application reports it and a normalised variant would need it:
$\lambda_{\min}(\hat\Omega_i)\ge\varepsilon_\Omega$ surely, by \eqref{eq:omegafloor} read at
$\hat\varrho(\Delta Z_i)$. \emph{At $\mathsf L=2$ the box and \eqref{eq:trunset} coincide and clause
\textup{(c)} says nothing new}, which is why no earlier round had to distinguish them.
% ARCHAEOLOGY, ROUNDS 2-4.  Round 2 printed a further clause here bounding the
% NORMALISED contrast, with "at most two entries are nonzero" and the Euclidean
% bound min{sqrt 2, sqrt(sans-L - 1)} eps_varrho^{-1}.  BOTH WERE WRONG: at
% W_i = 0 every entry equals -1/hat-varrho_0, so the Euclidean norm there is
% sqrt(sans-L - 1)/hat-varrho_0, and the round-2 constant understates at
% sans-L >= 4 -- AT SECTION 6's OWN DESIGN -- by sqrt(3)/sqrt(2).  Round 3
% corrected it in place.  ROUND 4 REMOVES THE CLAUSE, the estimator no longer
% carrying the normalised contrast; the corrected arithmetic moves to rem:htread,
% where the reading now lives.  Round 3 also grounded eq:trunset on the need for
% Omega-hat^{-1} to exist; THAT GROUND GOES WITH THE NORMALISATION, and clause
% (c) states the one that survives it.
\item[(3)] \emph{If only a smaller lower bound $\varepsilon\in(0,\varepsilon_\varrho]$ is known
to the analyst, truncating at $\varepsilon$ serves equally and every constant below carries
$\varepsilon$ in place of $\varepsilon_\varrho$}, provided $\varepsilon<1/\mathsf L$, without
which \eqref{eq:trunset} is empty at $\varepsilon$ and there is nothing to project onto. At
$\mathsf L=2$ that proviso is $\varepsilon<\tfrac12$ and is implied by
$\varepsilon\le\varepsilon_\varrho$; at $\mathsf L\ge3$ it is also implied, since A4(i) already puts
$\varepsilon_\varrho<1/\mathsf L$. It is stated because a reader who reads clause \textup{(3)} alone,
as a licence to pick any smaller $\varepsilon$, must not read it as a licence to pick a larger one.
\item[(4)] \eqref{eq:hatT} is the object A7 names.
\end{enumerate}
\end{remark}

\begin{remark}[How the kernel block C3 is read]\label{rem:app3read}
Two clauses, both of which stood inside the statement of C3. Neither is a condition, and the block
restricts exactly the same kernels without them.
\begin{enumerate}
\item[(1)] \emph{What \eqref{eq:rescaled} is for.} $\Theta_{I}$ is the object every bound below C3 is
stated for, and the object \eqref{eq:intker} integrates.
\item[(2)] \emph{How clause \textup{(ii)} is read.} The family is indexed by the \textbf{realised
truncation interval} $I\in\mathfrak I$ --- the interval, in scaled coordinates, over which the kernel
is actually integrated at the evaluation point. That index encodes both \emph{how far} the nearer
endpoint is, in bandwidths, and \emph{which} endpoint it is; the scalar boundary index of earlier
rounds, and the reflection convention it needed, are absorbed into it. Support is no longer a
condition: what \textup{(ii)} requires is the \textbf{envelope} $\bar\Theta$, and the defining property
\eqref{eq:truncmom} is a \textbf{truncated} moment condition, as before.
\end{enumerate}
% ARCHAEOLOGY.  Clause (2) read, through dispatch 25, "The boundary index is a
% distance in BANDWIDTHS".  That was the round-12 repair of a genuinely ambiguous
% scalar; the interval index makes the distance one of two things the index
% carries, and the reflection the other.  Clause (1) is unchanged but for the
% symbol.
\end{remark}

% ===========================================================================
%  rem:htread IS NEW IN DISPATCH 25, ROUND 4, AND IT IS THE ONLY NUMBERED
%  ENVIRONMENT THE DISPATCH ADDS.  This file is position 5 of 8, so every
%  theorem-like number in results_s3_supp.tex, results_supp.tex and proofs.tex
%  ADVANCES BY ONE.  Nothing breaks -- every reference is a \cref -- but printed
%  numbers move, and the file's own contract requires that to be said.
%
%  WHY IT EXISTS.  Rounds 2 and 3 put Omega^{-1} into the estimator, on the
%  operator and then on both sides.  THE AUTHOR REVERSED THAT; the estimator
%  carries the raw contrast.  What the two rounds established about what the
%  contrast MEANS is worth keeping, and this is where it is kept.  IT IS AN
%  INTERPRETATION, NOT A METRIC, and the remark says so in its first sentence so
%  that no later reader mistakes it for a specification.
%
%  IT ALSO CARRIES THE CORRECTION OF THE DISPATCH'S WORST ERROR.  Round 3 stated
%  the moments of the normalised contrast backwards at both ends: it computed the
%  BOUNDED quadratic form and cited it where the DIVERGENT squared norm belonged,
%  and it asserted non-integrability where the true failure is undefinedness.
%  Both are corrected here, with the arithmetic printed, because a remark that
%  exists to explain an IPW object must not repeat the standard IPW error.
% ===========================================================================

\begin{remark}[The inverse-probability reading of the contrast, and why the estimator does not use it]\label{rem:htread}
This remark is interpretation. Nothing in the assumption block, in $\psi$ or in the operator is
stated at the object it introduces, and the paper does \textbf{not} estimate at it.
\begin{enumerate}
\item[(1)] \emph{The normalised contrast.} Wherever every cell is non-null, so that $\Omega(\Delta Z)$
is invertible, set
\begin{equation}\label{eq:omeganorm}
\widetilde W^{\Omega}\;:=\;\Omega(\Delta Z)^{-1}\widetilde W ,
\qquad\text{so that}\qquad
\widetilde W^{\Omega}_l\;=\;\frac{\un\{W=l\}}{\varrho_l(\Delta Z)}-\frac{\un\{W=0\}}{\varrho_0(\Delta Z)} ,
\end{equation}
the second equality being \eqref{eq:omegainv}. \emph{This is the vector of inverse-probability
{\upshape(}Horvitz--Thompson{\upshape)} contrasts of each cell against the base}, and at $\mathsf L=2$ it
is $(W-\varrho)/\{\varrho(1-\varrho)\}=\un\{W=1\}/\varrho-\un\{W=0\}/(1-\varrho)$, the identity a binary
design writes down without comment.
\item[(2)] \emph{What it is for.} By \eqref{eq:omegaprefac} the raw contrast has conditional forward
residual $\Omega(z)d(u\mid z)$, a matrix mixture of the cell-versus-base differences, while
$E[\widetilde W^{\Omega}\un\{\Delta U\le u\}\mid\Delta Z=z]=d(u\mid z)$ is the difference itself. \emph{So
\eqref{eq:omeganorm} is the reading under which the restriction is most directly legible}, and the
quadratic form $\widetilde W'\Omega^{-1}\widetilde W=(1-\varrho_W)/\varrho_W$ is base-cell-free where the
componentwise weighting $\operatorname{diag}(\varrho_l(1-\varrho_l))^{-1}$ is not.
\item[(3)] Two moments, and they are different objects. Conditionally on $\Delta Z$, at a point
where every cell is non-null,
\[
E\big[\widetilde W'\Omega^{-1}\widetilde W\,\big|\,\Delta Z\big]=\operatorname{tr}(\Omega^{-1}\Omega)=\mathsf L-1 ,
\qquad
E\big[\big\|\Omega^{-1}\widetilde W\big\|^{2}\,\big|\,\Delta Z\big]=\operatorname{tr}\big(\Omega^{-1}\big)
=\sum_{l\ge1}\frac1{\varrho_l}+\frac{\mathsf L-1}{\varrho_0} .
\]
The first is bounded by $\mathsf L-1$ whatever the cell probabilities are; the second diverges as
any cell probability tends to zero. At $\mathsf L=4$ with a thin non-base cell at $0.05$ and the other
three at $0.31\overline{6}$ they are $3.00$ and $35.79$, against $\varepsilon_\varrho^{-1}=20$: \emph{the
second is not controlled by the level bound}, exceeding it here by a factor of $1.79$. Under A4(i) it is
at most $2(\mathsf L-1)/\varepsilon_\varrho$, so the \emph{worst case} is a factor $2(\mathsf L-1)$, 
$6$ at $\mathsf L=4$, which the configuration above does not attain.
% ROUND-5 CORRECTION.  Round 4 wrote "exceeding it by a factor of order sans-L"
% and attached that to THIS example, whose realised ratio is 35.79/20 = 1.79.
% The order-sans-L statement is the WORST CASE, 2(sans-L - 1), reached only as a
% cell probability approaches the floor with the others spread differently.
% Attaching a worst case to a worked example as though it were the example's
% value is the species of over-claim this block polices elsewhere.
\item[(4)] \emph{And the first absolute moment is free of $\varrho$ altogether}:
$E[|\widetilde W^{\Omega}_l|\mid\Delta Z]=\varrho_l/\varrho_l+\varrho_0/\varrho_0=2$, exactly, at every
$l$ and every cell configuration. This is the standard inverse-probability pattern, the mean is
untouched and the variance diverges, and it is the reason (3)'s two displays must not be substituted
for one another.
\item[(5)] Where $\Omega$ is singular, $\widetilde W^{\Omega}$ is \textbf{undefined}, not
non-integrable. $\Omega$ has no inverse there and \eqref{eq:omeganorm} names nothing; where it is
defined, clause (4) shows it is integrable with a $\varrho$-free bound. The two failures are
different and the corpus must not conflate them.
\item[(6)] Why the estimator does not use any of this. Normalising would replace
$\mathcal A_{\mathrm{cond}}$ by a different operator, and A2's range condition, A$2'$, A1c and H1, 
the last through $\sigma_j(\mathcal A_{\mathrm{cond}})\asymp j^{-a}$, are all stated at
$\mathcal A_{\mathrm{cond}}$, with $\mu_{\text{\textup{H1}}}=s/(2a)$ and the smoothness-budget table's
two daggered lines downstream of $a$, \eqref{eq:frontier} is \emph{not}, its $\mu$ being primitive
at \textup{(P-a.3)}.
% REPAIRED 2026-08-09.  This clause printed TWO errors, both struck elsewhere in
% this file and both surviving here: a bare "mu = s/(2a)", against rem:pasrc(5)'s
% binding reading and against the striking of that identity recorded at :1927,
% :1931 and :1938; and the claim that eq:frontier is downstream of `a', which
% round 12 struck -- eq:frontier's mu is PRIMITIVE at (P-a.3) and does not move
% with `a' in either direction.
%
% The identical pair was found and repaired at :964 by the budget-table re-read,
% which correctly declined to touch THIS site because rem:htread is Phase-1
% settled text and outside the re-read's declared scope.  It named the two-token
% repair instead of making it silently, and the orchestrator applied it here.
% That is the scope discipline working as intended: the twin was declared, not
% quietly fixed and not quietly left.
\emph{A4's note (c) already records
that the two operators must not be assumed to share range or kernel}, and the mechanism is the same one
that defeats the weight-class argument in the paragraph \emph{A4}: $\mathcal W$'s $z$-sections are
real-analytic while $\Omega$ is only H\"older. \emph{The gain would be interpretational and the price is
the whole smoothness-and-ill-posedness budget}, so the paper keeps the raw contrast and keeps the reading
here.
\item[(7)] \emph{What is \textbf{not} claimed.} $\Omega^{-1}$ is not the efficient weight: under A1 the
conditional variance of the indicator moment is $\Omega(z)F(u\mid z)$, so the efficient weight differs
from $\Omega^{-1}$ by a scalar depending on both $u$ and $z$, and no efficiency calculation for this
estimator exists in this corpus. \emph{Nor is any conditioning improvement claimed}: $\Omega^{-1}$
up-weights thin cells, which is where the application's own design sits.
\end{enumerate}
\end{remark}
